\documentclass[a4paper]{ctexart}
\usepackage{amsmath,amsthm,amssymb,graphicx,mathrsfs,physics,cancel,tikz-cd}
\usepackage{enumitem}
\usepackage[
    a4paper,
    left=2.5cm,
    right=2.5cm,
    top=3cm,
    bottom=3cm,
]{geometry}
\usepackage[colorlinks=true, linkcolor=blue, citecolor=blue, urlcolor=blue,hidelinks]{hyperref}
\renewcommand{\contentsname}{Contents}
\title{Calculus}
\author{Star}

\begin{document}
	\maketitle
	\newpage
	\tableofcontents
	\newpage
	\newcommand{\M}{\mathbf{M}}
    \newcommand{\R}{\mathbb{R}}
    \newcommand{\C}{\mathbb{C}}
    \newcommand{\Q}{\mathbb{Q}}
    \newcommand{\N}{\mathbb{N}}
    \newcommand{\Z}{\mathbb{Z}}
	\theoremstyle{definition}
	\renewcommand{\proofname}{Proof}
	\section{Preparations: Ten Axioms}
	\subsection{Fields and the Field Axioms}
	\newtheorem{mydef}{Definition}[subsection]
	\newtheorem{mypro}[mydef]{Proposition}
	\newtheorem{mythm}[mydef]{Theorem}
	\newtheorem{myrmk}[mydef]{Remark}
	\newtheorem{mynot}[mydef]{Notation}
	\newtheorem{myex}[mydef]{Example}
	\newtheorem{myax}[mydef]{Axiom}
	%\newtheorem{proof}{Proof}
	\newtheorem{mymthd}[mydef]{Method}
	\newtheorem{mylma}[mydef]{Lemma}
	\begin{mydef}
		A field $\mathcal{K}$ is a set equipped with addition and multiplication, denoted by $(\mathcal{K},+,\cdot)$, satisfying the following axioms(also known as the $\mathbf{Field\ Axioms}$):
		\begin{enumerate}
			\item Associative Laws: let $a,b,c \in \mathcal{K}$, $a+b+c=a+(b+c) \quad a\cdot b\cdot c=a\cdot (b\cdot c) $
			\item Existence of Identity Elements: let $a\in K$, $a+0_{\mathcal{K}}=a\quad a\cdot 1_{\mathcal{K}}=a$, where $0_{\mathcal{K}}$ and $1_\mathcal{K}$ is called the additive identity and multiplicative identity.
			\item Existence of Negatives: for any $a\in \mathcal{K}$, there exists one element in $\mathcal{K}$ denoted by $(-a)$ such that $a+(-a)=0_\mathcal{K}$.
			\item Existence of Reciprocal: for any $a\in \mathcal{K}$ and $a\neq 0_\mathcal{K}$, there exists one element in $\mathcal{K}$ denoted by $a^{-1}$ such that $a\cdot a^{-1} = 1_\mathcal{K}$. 
			\item Commutative Laws: let $ a,b\in \mathcal{K}$, $a+b=b+a \quad a\cdot b=b\cdot a$
			\item Attributive Law: let $a,b,c \in \mathcal{K}$, $a\cdot (b+c)=a\cdot b + a\cdot c$
		\end{enumerate}
	\end{mydef}	
	\begin{myrmk}
		Closure is contained in the defined addition and multiplication.
	\end{myrmk}
	\begin{mythm}
		Let $a,b,c \in \mathcal{K}$, if $a+c=b+c$, then $a=b$. This is called the cancellation law.
	\end{mythm}
	\begin{mythm}
		$0_\mathcal{K}$ and $1_\mathcal{K}$ in $\mathbf{Field\ Axiom\ 2}$ is unique.
	\end{mythm}
	\begin{myex}
		Note that $\mathbb{C,R,Q}$ are fields, but $\mathbb{N}$ is not a field because $2 \in \mathbb{N}$ but its reciprocal $\frac{1}{2} \notin \mathbb{N}$.This is the main difference between $\mathbb{Q}$ and $\mathbb{N}$.
	\end{myex}
	\begin{mydef}\label{def_of_ss}
		Let $\mathcal{F}$ be a field. A set $\mathcal{K}$ is a subset of $\mathcal{F}$. We shall say $\mathcal{K}$ a subfield of $\mathcal{F}$ when $\mathcal{K}$ satisfies the followings:
		\begin{enumerate}
			\item If $x,y$ are elements of $\mathcal{K}$, then  $x+y\in \mathcal{K}$ and $x\cdot y$ $\in \mathcal{K}$.
			\item If $x\in \mathcal{K}$, then $(-x)\in \mathcal{K}$. If furthermore $x\neq 0$ , $x^{-1}$ $\in \mathcal{K}$.
			\item $0_\mathcal{K}$ and $1_\mathcal{K}$ are in $\mathcal{K}$.
		\end{enumerate}
	\end{mydef}
	\begin{myrmk}
		Definition~\ref{def_of_ss} provides a simpler way to prove a set is a field.
	\end{myrmk}
	\begin{mypro}
		Let $\mathcal{K}$ be a field and $A,B$ are subfields of $\mathcal{K}$, then $A \cap B$ is also a subfield of $\mathcal{K}$.
	\end{mypro}
	\begin{mypro}
		Let $\mathcal{K}$ be a field and $A,B$ are subfields of $\mathcal{K}$, then $A \cup B$ is also a subfield of $\mathcal{K}$ if and only if $A \subseteq B$ or $B \subseteq A$.
	\end{mypro}
%	\begin{mydef}
%		Let $A,B$ be two sets. We define the sum of $A$ and $B4 as below:
%		\begin{equation}
%			A+B:=\{x\, |\, x=a+b,a\in A,\, b \in B\}\notag
%		\end{equation}
%	\end{mydef}
%	\begin{mypro}
%		Let $\mathcal{K}$ be a field and $A,B$ are subfields of $\mathcal{K}$, then $A+B$ is also a subfield.
%	\end{mypro}
	\begin{myex}
		Note that $\mathbb{R}$ is a subfield of $\mathbb{C}$.
	\end{myex}
	\subsection{The Order Axioms}
	\begin{mydef}
		We shall assume that there exists a set $\mathbb{R}^+ \subset \mathbb{R}$, called the set of positive real numbers, which satisfies the followings(also known as the $\mathbf{Order\ Axioms}$):
		\begin{enumerate}
			\item For any $x,y\in \mathbb{R}^+$, we have $x+y \in \mathbb{R}^+$ and $x\cdot y \in \mathbb{R}^+$.
			\item For any real $x \neq 0$, either $x\in \mathbb{R}^+$ or $(-x)\in \mathbb{R}^+$, but not both. 
			\item $0 \notin \mathbb{R}^+$.
		\end{enumerate}
	\end{mydef}
	\begin{mypro}
		$\mathbb{R}^+$ is non-empty since $1 \in \mathbb{R}^+$.
	\end{mypro}
	\begin{mynot}
		For any $a,b\in \mathbb{R}$, we define the symbols $>$, $<$, $\geq$, $\leq$ as below:
			\item $x>y$ means that $x-y$ is positive. \centering
			\item $x<y$ means that $y>x$. \centering
			\item $x\ge y$ means that $x>y$ or $x=y$. \centering
			\item $x \le y$ means that $y \ge x$. \centering
	\end{mynot}
	\begin{mythm}\label{trichotomy law}
		For any $a,b\in \mathbb{R}$, exactly one of the three relations holds: $a>b$, $a<b$ and $a=b$. This is called the trichotomy law.
	\end{mythm}
	\begin{myrmk}
		Theorem~\ref{trichotomy law} is always used to prove real number equality, while another commonly-used way is to prove both $a\ge b$ and $a \le b$ hold. The latter one is similar to the method of proving set equality.
	\end{myrmk}
	\begin{mythm}
		For any $a,b,c\in \mathbb{R}$, if $a<b$ and $b<c$, then $a<c$. This is called the transitive law.
	\end{mythm}
	\subsection{The Completeness Axiom}
	\begin{mydef}\label{def_of_upb}
		Let $S$ be an nonempty set of real numbers, for any $x$ $\in$ S, if we have	$x\le B$, then we say $S$ is bounded above by $B$ and $B$ is a upper bound of $S$. Furthermore, if $B\in S$, then we say $B$ the maximum of $S$, denoted by $\max S$. 
	\end{mydef}
	\begin{myrmk}
		Similarly, we can define lower bound and minimum of a nonempty set.
	\end{myrmk}
	\begin{mydef}
		A number $B$ is called a least upper bound (a supremum) of a nonempty set $S$, denoted by $\sup S$, if $B$ satisfies the followings:
		\begin{enumerate}
			\item $B$ is an upper bound of $S$.
			\item There's no number smaller than $B$ being an upper bound of $S$.
		\end{enumerate}
	\end{mydef}
	\begin{myrmk}
		Similarly, we can define a greatest lower bound (a infimum) of a nonempty set.
	\end{myrmk}
	\begin{myax}
		Every nonempty set $S$ of real numbers which is bounded above has a supremum, i.e., there exists such a number B that $B=\sup S$. This is called the $\mathbf{ Completeness\ Axiom}$.
	\end{myax}
	\begin{myrmk}
		The completeness axiom indicates a crucial difference between $\mathbb{Q}$ and $\mathbb{R}$ because $\mathbb{Q}$ is not complete. For instance, it's impossible to find a least upper bound in $\mathbb{Q}$ for the set $\{ x\in \mathbb{Q}:x<\sqrt{2}\}$.
	\end{myrmk}
	\begin{mythm}
		Every nonempty set of real numbers which is bounded below has a infimum.
	\end{mythm}
	\begin{mydef}
		A set $S$ of real numbers is bounded if there exists a real number $M\ge 0$ such that
		\begin{align*}
			|x|\le M\quad \text{for all } x \in S.
		\end{align*}
	\end{mydef}
	\begin{mythm}
		The set $\mathbb{P}$ of positive integers $1,2,3...$ is not bounded above.
	\end{mythm}
	\begin{mythm}
		For any $x\in \mathbb{R}$, there exists an positive integer $n$ such that $n>x$.
	\end{mythm}
	\begin{mythm}
		If $x>0$ and $y$ is an arbitrary real number, there exists an positive integer $n$ such that $nx>y$.
	\end{mythm}
	\begin{mythm}
		If three real numbers $x,y,a$ satisfy:
		\begin{equation}
			a\le x \le a+\frac{y}{n}\notag
		\end{equation}
		for every integer $n\ge 1$, then $x=a$.
	\end{mythm}
	\begin{mypro}\label{pro_r_dense}
		Suppose $x$ and $y$ are arbitrary real numbers. If $x<y$, then there exists a rational number r such that $x<r<y$. Moreover, there exists infinite many different rational numbers $r_1,r_2,...$ such that $x<r_i<y$ for all i $\in \mathbb{Z}^+$.
	\end{mypro}
	\begin{myrmk}
		Proposition~\ref{pro_r_dense} is often described by saying that rational numbers are $\mathbf{dense}$ in the real numbers.
	\end{myrmk}
	\begin{mypro}\label{pro_ir_dense}
		Suppose $x$ and $y$ are arbitrary real numbers. If $x<y$, then there exists a irrational number z such that $x<z<y$. Moreover, there exists infinite many different irrational numbers $z_1,z_2,...$ such that $x<z_i<y$ for all i $\in \mathbb{Z}^+$.
	\end{mypro}
	\begin{myrmk}
		Proposition~\ref{pro_ir_dense} is often described by saying that irrational numbers are $\mathbf{dense}$ in the real numbers.
	\end{myrmk}
	\begin{mypro}
		If $x$ is a rational number, $x\neq 0$, and $y$ is an irrational number, then $x+y,x-y,xy,\frac{x}{y},\frac{y}{x}$ are all irrational.
	\end{mypro}
	\begin{mypro}
		$\pi+e,\pi-e$ can't be both rational. This also holds for $e+\pi$ and $ \frac{e}{\pi}$. If we know that both $e$ and $\pi$ are transcendental numbers, then we can also prove that $\pi+e$ and $\pi e$, $\pi e$ and $\frac{e}{\pi}$ can't be both rational.
	\end{mypro}
	\begin{mythm}\label{pro1_sup}
		Let $S$ be an nonempty set of real numbers. If $S$ has a supremum, then there exists an element $x\in S$  such that for any  $h\in \mathbb{R}^+$:
		\begin{equation}
			x>\sup S-h\notag
		\end{equation} 
		If $S$ has an infimum, then there exists an element $y$ $\in$ S such that for any $h$ $\in \mathbb{R}^+$:
		\begin{equation}
			y<\inf S+h\notag
		\end{equation}
	\end{mythm}
	\begin{myrmk}\label{two_pro_sup}
		Theorem~\ref{pro1_sup} shows two properties of supremum (infimum) vividly and precisely. On one hand, it's greater (smaller) than or equal to any element of $S$. On the other hand, it's smaller (greater) than or equal to any upper bounds (lower bounds). This theorem also shows the uniqueness of supremum (infimum).
	\end{myrmk}
	\begin{mydef}
		Let $A,B$ be two sets. We define the sum of $A$ and $B$ as below:
		\begin{equation}
			A+B:=\{x\, |\, x=a+b,a\in A,\, b \in B\}\notag
		\end{equation}
	\end{mydef}
	\begin{mythm}
		Let $A,B$ be two non-empty sets of real numbers. If both $A$ and $B$ are bounded above, then:
		\begin{equation}
			\sup (A+B)=\sup A+\sup B\notag
		\end{equation}
		If both A and B are bounded below, then:
		\begin{equation}
			\inf(A+B)=\inf A+\inf B\notag
		\end{equation}
	\end{mythm}
	\begin{mythm}
		Let $A,B$ be two non-empty sets of real numbers. If for any element $x\in A$ and $y\in B$, we have $x\le y$. Then:
		\begin{equation}
			\sup A \le \inf B\notag
		\end{equation}
	\end{mythm}
	\subsection{Inductive Sets and Induction}
	\begin{mydef}
		We call a set $S$ an inductive set if it satisfies the followings:
		\begin{enumerate}
			\item $1\in S$. 
			\item if $x\in$ S, then $x+1 \in S$.
		\end{enumerate}
		Then we could define positive integers more rigorously: let $\mathcal{F}$ be the collection of all inductive sets, then:
		\begin{equation}
			\mathbb{Z}^+ := \bigcap_{I\in \mathcal{F}}I \notag
		\end{equation}
	\end{mydef}
	\begin{myrmk}
		Note that $\mathbb{Z}^+$ is the subset of all inductive sets. After we define $\mathbb{Z}^+$, it's easy to define integers $\mathbb{Z}$ as below:
		\begin{equation}
			\mathbb{Z}:=\mathbb{Z}^+ \cup \{0\} \cup \{x:(-x)\in \mathbb{Z}^+\}\notag
		\end{equation}
		Then we can define rational numbers $\mathbb{Q}$ as below:
		\begin{equation}
			\mathbb{Q}:=\{x:x=\frac{p}{q},\text{ where } p,q \in \mathbb{Z}, q \neq 0 \text{ and } \gcd(p,q)=1\}\notag
		\end{equation}
		However, it's harder to construct real numbers though we have infinite ways. We shall only prove the existence of a square root because it's instructive to see how $\mathbf{Axiom\ 10}$ differentiates real numbers from rational numbers. We will finish our proof in section $1.5$.
	\end{myrmk}
	\begin{mythm}\label{pcp_ind}
		Let $S$ be a set of positive integers which has the following two properties:
		\begin{enumerate}
			\item $1\in S$.
			\item If an integer $k \in S$, then so is $k+1$.
		\end{enumerate}
		Then every positive integer is in the set $S$.
	\end{mythm}
	\begin{mymthd}\label{prove_by_indu}
		Let $A(n)$ be an assertion involving an integer $n$. We conclude that $A(n)$ is true for every $n\ge n_1$ if we can perform the following two steps:
		\begin{enumerate}
			\item Prove that $A(n_1)$ is true.
			\item Let $k$ be an arbitrary but fixed integer $\ge n_1$. Assume that $A(k)$ is true and prove that $A(k+1)$ is true
		\end{enumerate}
	\end{mymthd}
	\begin{myrmk}
		Theorem~\ref{pcp_ind} is called the principle of mathematical induction. When we're proving by induction, we're applying Theorem~\ref{pcp_ind} to the set $S:=\{n\in \mathbb{Z}^+:A(n+n_1-1)\ is\ true.\}$.
	\end{myrmk}
	\begin{myrmk}
		Method~\ref{prove_by_indu} is also called the first method of proof by induction. If we assume that $A(n)$ is true for every integer $n\le k$ and prove that $A(k+1)$ is true, this method is called the second method of proof by induction.
	\end{myrmk}
	\begin{mythm}\label{pcp_wod}
		Every nonempty set of positive integers has a smallest member.
	\end{mythm}
	\begin{proof}
		We shall prove by contradiction.
		
		Let $T$ a nonempty set of positive integers which doesn't has a smallest member, we define a set $S$:
		\begin{equation}
			S:=\{x:x<t\text{ for all } t \in T\}\notag
		\end{equation}
		Then $1$ shouldn't be in $T$ because $1$ is the smallest positive integer. Thus, $1 \in S$. 
		
		For a positive integer $k$, if $k\in S$, that is, $k<t$ for all $t\in T$. We shall consider $k+1$. For these two statements $k+1\in S$ and $k+1 \notin S$, only one of them holds. If we could prove both of them are false, then we reach a contradiction. 
		
		If $k+1 \notin S$, then there exists $t_1\in T$ such that $k+1\ge t_1$. Since $T$ has no smallest member, there exists $t_2 < t_1$ and we duduce $k+1>t_2$, that is, $k \ge t_2$, which contradicts with $k>t$ for all $t\in T$. 
		
		If $k+1\in S$, then we shall perform the principle of induction to $S$, that is, every positive integer is in $S$. Since $T$ is nonempty, there exists one positive integer $t\in T$. However, this $t$ is also in $S$, then we deduce $t<t$ by the definition of $S$, which is impossible.
		
		\fbox{Hence, the nonempty set $T$ must have a smallest member.}
	\end{proof}
	\begin{myex}
		Here's one example of proving by well-ordered principle.\\ Let n be a positive integer.
		\begin{equation}
			1+2+\cdots+n=\frac{n(n+1)}{2}\notag
		\end{equation}
		It's trivial that when $n=1$, the equality holds. Then there're two possible statements:
		\begin{enumerate}
			\item For all positive integer, the equality holds.
			\item There exists at least one positive integer such that the equality doesn't hold.
		\end{enumerate}
		We shall prove statement 2 is wrong so that statement 1 naturally holds. We apply the well-ordered principle to the collection of all positive integer that satisfy statement 2, thus we get one positive integer $k>1$ and for $k-1$ the equality holds. Now we add $k$ to both sides of equality and find the equality also holds for $k$, which leads to a contradiction. Thus statement 1 holds.
	\end{myex}
	\begin{myrmk}
		Theorem~\ref{pcp_wod} is called the well-ordered principle, which is one important property of positive integers since it doesn't hold for arbitrary sets of integers. For example, $\mathbb{Z}$ doesn't have smallest member.
	\end{myrmk}
	\begin{mynot}
		We define the summation notation as below:
		\begin{enumerate}
			\item $\displaystyle\sum_{i=1}^1a_i=a_1$
			\item $\displaystyle\sum_{i=1}^{n+1}a_{i}=\left(\sum\limits_{i=1}^n a_i \right) + a_{n+1}$, $n\ge 1$
		\end{enumerate}
	\end{mynot}
	\begin{myrmk}
		It's easy to conclude that proof by induction and definition by induction involve the same underlying idea, that is, recursion.
	\end{myrmk}
	\begin{myex}
		We shall see another example of inductive definition. Suppose $x$ a positive real number and consider the following inductive definition.
		\begin{itemize}
			\item Let $a_0$ denote the largest integer $\le x$.
			\item Having chosen $a_0$, we let $a_1$ denote the largest integer such that $a_0+\frac{a_1}{10}\le x$.
			\item More generally, having chosen $a_0,a_1,\cdots,a_{n-1}$, we let $a_n$ denote the largest integer such that
			\begin{align*}
				\sum_{i=0}^{n}\frac{a_i}{10^i}\le x
			\end{align*}
		\end{itemize}
		Define set $S$ to be the set of all numbers obtained in the above way, i.e.
		\begin{align*}
			S:=\left\{\sum_{i=0}^{n}\frac{a_i}{10^i}\ :\ n\in \mathbb{Z}^+\right\} \cup \left\{0 \right\}
		\end{align*}
		We have the following propositions:
		\begin{enumerate}
			\item $S$ is bounded above and $x=\sup S$.
			\item Let $a_0,a_1,a_2,\cdots,$ be as the definition of $S$. We define the notation $a_0.a_1a_2a_3 \cdots$ to be the $\mathbf{decimal\ expansion}$ of $x$. Then if we replace all ``$\le$'' with ``$<$'', then one real number might have two different decimal expansions. This is merely a reflection of the fact that two different sets of of real numbers can have the same supremum.
			\item Note that $x$ satisfies the following inequality:
			\begin{align*}
				\sum_{i=0}^{n-1}\frac{a_i}{10^i}+\frac{a_n}{10^n}<x<\sum_{i=0}^{n-1}\frac{a_i}{10^i}+\frac{a_n+1}{10^n}
			\end{align*}
			This gives us a way to approximate $x$ from above and below, which could also serve as a way to construct irrational numbers from rational numbers.
		\end{enumerate}
	\end{myex}
	\subsection{Square Root}
	\begin{mythm}
		Every nonnegative real number a has one unique square root, denoted by $a^{1/2}$ or $\sqrt a$.
	\end{mythm}
	\begin{proof}
		Suppose $S:=\{x \in \mathbb{R}:\ x^2\le a\}$, $b:=\sup S$. We shall finish the proof within two steps:
		\begin{enumerate}
			\item Justify $S$ nonempty and $b$ well-defined.
			\item Prove $b^2=a$ by the trichotomy law.
		\end{enumerate}
		
		Note that $ \displaystyle \left(\frac{a}{a+1}\right)\in S$, thus S is nonempty. Since $(a+1)^2$ is an upper bound, $b$ is well-defined by the completeness axiom.
		
		First, we shall prove $b^2>a$ is false. If we could find a real number $c$ such that $b>c$ and $c^2>a$, then $c$ should be an upper bound smaller than $b$, which contradicts with the definition of supremum. Now we shall try to construct such $c$. Suppose $c=b-t$ where $t>0$, then:
		\begin{align*}
			c^2-a=(b-t)^2-a
			=b^2-2bt+t^2-a>b^2-2bt-a=0  \Rightarrow t=\frac{a-b^2}{2b}>0,\ c=b-\frac{a-b^2}{2b}
		\end{align*}
		
		Thus we successfully construct such $c$ and prove that $b^2>a$ doesn't hold. If we want to write the proof in order, we shall write as below:
		\begin{align*}
			c^2=(b-t)^2>b^2-2bt=b^2-2b\cdot \frac{a-b^2}{2b}=a
		\end{align*}
		
		Second, we shall prove $b^2<a$ is false. If we could find a real number $c$ such that $b<c$ and $c^2<a$, then $c$ should be in $S$ but $c$ is greater than $b$, which contradicts with the definition of supremum. Now we shall try to construct such $c$. Suppose $c=b+t$ where $0<t<b$, then:
		\begin{align*}
			c^2-a&=(b+t)^2-a
			=b^2+2bt+t^2-a<b^2+3bt-a=0\Rightarrow t=\frac{a-b^2}{3b}>0
			\\ t&=\min\left(\frac{a-b^2}{3b},b\right) \text{ satisifying all limitations.}
		\end{align*}
		
		Thus we successfully construct such $c$ and prove that $b^2<a$ doesn't hold. If we want to write the proof in order, we shall write as below:
		\begin{align*}
			c^2=(b+t)^2=b^2+2bt+t^2<b^2+3bt<b^2+3b\cdot \frac{a-b^2}{3b}=a
		\end{align*}
		
		\fbox{Hence, we have $b^2=a$.}
	\end{proof}
	\begin{myrmk}
		In this proof we witness the property of supremum, or let's say, the least upper bound as we mentioned in Remark~\ref{two_pro_sup}. You should feel surprised about this proof now thanks to my detailed instruction.
	\end{myrmk}
	\subsection{Absolute Value and Inequalities}
	\begin{mydef}
		We define the absolute value of a real number $x$ as below:
		\begin{align*}
			|x|:=\begin{cases}x,\ x\ge0
				\\ -x,\ x<0
			\end{cases}
		\end{align*}
	\end{mydef}
	\begin{mythm}
		If $a\ge 0$, $|x|\le a$ if and only if $-a\le x \le a$.
	\end{mythm}
	\begin{mythm}
		If $x,y\in \mathbb{R}$, then $|x+y| \le |x|+|y|$. The equality sign holds if and only if $xy\ge 0$.
	\end{mythm}
	\begin{myrmk}
		This is called the triangle inequality, which could be generalized to vectors or complex numbers. Take $x=a-b,\ y=b-c$, then we have $|a-c|\le |a-b|+|b-c|$, which is more commonly used.
	\end{myrmk}
	\begin{mythm}
		For arbitrary real numbers $a_1,a_2,\cdots,a_n$, we have:
		\begin{align*}
			\left|\sum_{i=1}^n a_i\right|\le \sum_{i=1}^n |a_i|
		\end{align*}
	\end{mythm}
	\begin{mythm}
		For arbitrary real numbers $a_1,a_2,\cdots,a_n$ and $b_1,b_2,\cdots,b_n$, we have:
		\begin{align*}
			\left(\sum_{i=1}^{n}a_i^2\right)\left(\sum_{i=1}^{n}b_i^2\right) \ge \left(\sum_{i=1}^{n}a_ib_i\right)^2
		\end{align*}
		The equality sign holds if and only if there exists a real number $x$ such that $a_ix+b_i=0$ for any $i\in \{1,2,\cdots n\}$.
	\end{mythm}
	\begin{myrmk}
		This is called Cauchy-Schwarz Inequality. The following two inequalities are more commonly used:
		\begin{align*}
			\left(\sum_{i=1}^{n}\frac{a_i^2}{b_i}\right)\left(\sum_{i=1}^{n}b_i\right) \ge \left(\sum_{i=1}^{n}a_i\right)^2\\
			\left(\sum_{i=1}^{n}\frac{a_i}{b_i}\right)\left(\sum_{i=1}^{n}a_ib_i\right) \ge \left(\sum_{i=1}^{n}a_i\right)^2
		\end{align*}
	\end{myrmk}
	\newpage
	\section{Integral Calculus}
	\subsection{Basic Concepts}
	\begin{mydef}
		Let $f$ be a correspondence. If each element of $\mathscr{D}_f$ has at most one image under $f$, we say that $f$ is a function.
	\end{mydef}
	\begin{myrmk}
		If you're confused about the definition of function above, you should refer to the Fundamental Algebra and Analysis.
	\end{myrmk}
%	\begin{mydef}
%		Let f be a function. If the domain of f is a collection of sets and function values are real numbers, we call f a set function.
%	\end{mydef}
	\begin{mydef}
		Suppose $f$ is a function (from $\mathbb{R}$ to $\mathbb{R}$) and $f(x)\ge 0$ for any $x\in \left[ a,b\right] $ where $a<b$. We define the set
		\begin{align*}
			S:=\left\{ (x,y)\in \mathbb{R}^2\ |\ a\le x \le b\text{ and } 0\le y\le f(x)\right\}
		\end{align*}
		to be the ordinate set of $f$.
	\end{mydef}
	\begin{myrmk}
		Definitely $S$ defined above is a plane region. Thus, we can naturally raise these questions: can we compute the area of $S$, or more fundamentally, what is the area of $S$, does it even make sense to talk about the area of $S$? These questions motivate the axiomatic definition of area and measurable sets in $\mathbb{R}^2$. (Here a measurable set is a set that we can assign an area to.)
	\end{myrmk}
	\begin{mydef}
		We assume that there exists $\mathscr{M}$ a collection of subsets of $\mathbb{R}^2$ and a function $A$ with $\mathrm{Dom}(A)=\mathscr{M}$ and $\mathrm{Im}(A)\in \mathbb{R}$ such that the following properties all hold:
		\begin{enumerate}
			\item Nonnegative property: For any $S\in \mathscr{M}$, $A(S)\ge 0$.
			\item Additive property: If $S,T \in \mathscr{M}$, then $S\cap T,S\cup T \in \mathscr{M}$. Additionally, $A(S\cup T)=A(S)+A(T)-A(S\cap T)$.
			\item Difference property: If $S,T \in \mathscr{M}$ and $S \subseteq T$, then $T-S$ (also $T\backslash S$) $\in \mathscr{M}$. Additionally, $A(T-S)=A(T)-A(S)$.
			\item Invariance under congruence: If $S\in \mathscr{M}$ and if $T\in \mathbb{R}^2$ is congruent to $S$, then $T \in \mathscr{M}$ and $A(T)=A(S)$.
			\item Choice of scale: Every rectangle $R$ is in $\mathscr{M}$. If the adjecent edges of $R$ are of length $w$ and $h$, then $A(R)=wh$.
			\item Exhaustion Property: Suppose $Q\subseteq \mathbb{R}^2$. If there exists a unique $c\in \mathbb{R}$ such that $A(S)\le c \le A(T)$ for all step regions $S$ and $T$ with $S\subseteq Q \subseteq T$, then $Q\in \mathscr{M}$ and $A(Q)=c$.
		\end{enumerate}
		We call such $A$ an area function and each $S\in \mathscr{M}$ a measurable set.
	\end{mydef}
	\begin{mydef}[Additional] Some definitions are lost above and here we add it.
		\begin{enumerate}
			\item Two sets $S$ and $T$ in $\mathbb{R}^2$ are called \textbf{congruent} if there exists a \textbf{one-to-one correspondence} (or a \textbf{bijection}) between $S$ and $T$ such that the distance between points are preserved, i.e., if the two points 
			$s_1 = (a,b), \; s_2 = (c,d)$ in $S$ are mapped to the two points 
			$t_1 = (x,y), \; t_2 = (p,q)$ in $T$ under such a bijection, then we must have
			\[
			\sqrt{(a - c)^2 + (b - d)^2} = \sqrt{(x - p)^2 + (y - q)^2}.
			\]
			\item A set $R \subseteq \mathbb{R}^2$ is called a \textbf{rectangle} if it is congruent to a set of the form
			\[
			\{ (x,y) \in \mathbb{R}^2 \mid 0 \le x \le w \text{ and } 0 \le y \le h \},
			\]
			where $w \ge 0$ and $h \ge 0$.
			\item A set $S \subseteq \mathbb{R}^2$ is called a \textbf{step region} if it is the union of a finite collection of 
			adjacent rectangles whose bases all rest on the $x$-axis.
		\end{enumerate}
	\end{mydef}
	\begin{mythm} Here are some theorems about some basic cases.
		\begin{enumerate}
			\item $S_1:=\left\{ (a,b)\right\} $ where $a,b\in \mathbb{R}$ is measurable and $A(S_1)=0$.
			\item $S_2:=\left\{ (a_1,b_1),(a_2,b_2),\cdots,(a_n,b_n)\right\}$ where $n\in \mathbb{Z}^+$ and $a_i,b_i\in \mathbb{R}$ for $i=1,2,\cdots, n$ is measurable with $A(S_2)=0$.
			\item If $S_3$ is the union of a finite collection of line segments in $\mathbb{R}^2$, then $S_3$ is measurable with $A(S_3)=0$.
			\item Suppose $S_4$ is the triangle with vertices $(0,0),(w,0),(\frac{w}{2},h)$ where $w,h>0$, $S_4$ is measurable and $A(S_4)=\frac{1}{2}wh$.
		\end{enumerate}
	\end{mythm}
	\begin{mydef}
		A \textbf{partition} $P$ of a given closed interval $[a,b]$ is a collection of points $x_0,x_1,x_2,\cdots,x_{n}$ satisfying:
		\[
		a=x_0<x_1<x_2<\cdots<x_n=b
		\]
		and we use the symbol
		\[
		P=\{x_0,x_1,x_2,\cdots,x_{n}\}
		\]
		to designate this partition. We call $[x_{k-1},x_k]$ the $k$th closed subinterval and $(x_{k-1},x_k)$ the $k$th open interval.
	\end{mydef}
	\begin{myrmk}
		Though we use set notation to designate a partition, the order \textbf{does} matter.
	\end{myrmk}
	\begin{mydef}
		A function $s$, whose domain is a closed interval $[a,b]$, is called a \textbf{step function} if there is a partition $P=\{x_0,x_1,\cdots,x_n\}$ of $[a,b]$ such that for each $k=1,2,\cdots,n$, there is a real number $s_k$ satisfying
		\[
		s(x)=s_k \quad \text{if} \quad x_{k-1}<x<x_k.
		\]
	\end{mydef}
	\begin{myex}
		We define a function $D(x)$ on $\mathbb{R}$:
		\[
		D(x):=
		\begin{cases}
			1, & \text{if } x \in \mathbb{Q}, \\[4pt]
			0, & \text{if } x \in \mathbb{R} \setminus \mathbb{Q}.
		\end{cases}
		\]
		which is called the Dirichlet Function. It's easy to verify that $D(x)$ is not a step function. Even if we confine Dirichlet Function on $[0,1]$, it is not a step function.
	\end{myex}
	\begin{myrmk}
		At each of the endpoints $x_{k-1}$ and $x_k$ the function must have some well-defined value, but this need not be the same as $s_k$.
	\end{myrmk}
	\begin{mydef}
		For a given partition $P$ of $[a,b]$, a \textbf{refinement} is a new partition $P'$ satisfying $P\subseteq P'$ (order should be ensured). Then $P'$ is said to be finer than P.
	\end{mydef}
	\begin{mydef}
		For two partitions $P_1,P_2$ of $[a,b]$, the union of $P_1$ and $P_2$ is called the \textbf{common refinement} of $P_1$ and $P_2$.
	\end{mydef}
	\begin{mythm}
		The sum and product of two given step functions are also step functions.
	\end{mythm}
	
	\subsection{Integral for Step Functions}
	\begin{mydef}
		Let $s$ be a step function defined on $[a,b]$, and let $P={x_0,x_1,\cdots,x_n}$ be a partition of $[a,b]$ such that:
		\[
		s(x)=s_k \quad \text{if}\quad x_{k-1}<x<x_k,\quad k=1,2,\cdots,n
		\]
		The integral of $s$ from $a$ to $b$, is defined as:
		\[
		\int_{a}^{b}s(x)\,\mathrm{d}x=\sum_{k=1}^{n}s_k\cdot (x_k - x_{k-1})
		\]
	\end{mydef}
	\begin{myrmk}
		The definition is constructed so that the integral of a nonnegative step function is equal to the area of its ordinate set.
	\end{myrmk}
	\begin{mythm}
		Here are some properties of the integral of a step function. If not stated, $s(x)$ and $t(x)$ are step functions well defined on $[a,b]$.
		\begin{enumerate}
			\item Additive Property:
			\[
			\int_{a}^{b}\left[s(x)+t(x)\right]\,\mathrm{d}x=\int_{a}^{b}s(x)\,\mathrm{d}x+\int_{a}^{b}t(x)\,\mathrm{d}x.
			\]
			\item Homogeneous Property: 
			\[
			\int_a^b c \cdot s(x) \, \mathrm{d}x = c \int_a^b s(x) \, \mathrm{d}x \quad \text{for every real } c.
			\]
			\item Linearity Property: Combining Additive Property and Homogeneous Property, we have
			\[
			\int_a^b \left[ c_1 s(x) + c_2 t(x) \right] \, \mathrm{d}x
			= c_1 \int_a^b s(x) \, \mathrm{d}x + c_2 \int_a^b t(x) \, \mathrm{d}x.
			\]
			for every real $c_1$ and $c_2$.
			\item Comparison Theorem: If $s(x)<t(x)$ for every $x$ in $[a,b]$, then
			\[
			\int_a^b s(x) \, \mathrm{d}x < \int_a^b t(x) \, \mathrm{d}x.
			\]
			\item Additivity with respect to the interval of integration:
			\[
			\int_a^c s(x) \, \mathrm{d}x + \int_c^b s(x) \, \mathrm{d}x = \int_a^b s(x) \, \mathrm{d}x 
			\quad \text{if} \quad a < c < b.
			\]
			\item  Invariance under translation:
			\[
			\int_a^b s(x) \, \mathrm{d}x = \int_{a+c}^{b+c} s(x - c) \, \mathrm{d}x 
			\quad \text{for every real } c.
			\]
			\item Expansion or Contraction of the interval of integration:
			\[
			\int_{ka}^{kb} s\!\left(\frac{x}{k}\right) \, \mathrm{d}x = 
			k \int_a^b s(x) \, \mathrm{d}x 
			\quad \text{for every } k > 0.
			\]
		\end{enumerate}
	\end{mythm}
	\begin{myrmk}
		Since the proofs of properties above are quite similar, we shall only prove Property 2 and 7 here as examples.
	\end{myrmk}
	\begin{proof}[Property 2]
		Let \( P = \{x_0, x_1, \ldots, x_n\} \) be a partition of \([a, b]\) such that
		\[
		s(x) = s_k \quad \text{if}\quad x_{k-1} < x < x_k\ (k = 1, 2, \ldots, n). 
		\]
		Then
		\[
		\int_a^b c \cdot s(x) \, \mathrm{d}x
		= \sum_{k=1}^{n} c \cdot s_k \cdot (x_k - x_{k-1})
		= c \sum_{k=1}^{n} s_k \cdot (x_k - x_{k-1})
		= c \int_a^b s(x) \, \mathrm{d}x.
		\]
	\end{proof}
	\begin{proof}[Property 7]
	Let \( P = \{x_0, x_1, \ldots, x_n\} \) be a partition of \([a, b]\) such that
	\[
	s(x) = s_k \quad \text{if}\quad x_{k-1} < x < x_k\ (k = 1, 2, \ldots, n). 
	\]
	Let \[ t(x) = s\left(\frac x k\right) \quad  \text{if}\quad  ka \le x \le kb. \]  
	Then \[ t(x) = s_i \quad \text{if} \quad k x_{i-1} < x< k x_i.\]
	Hence \( P' = \{k x_0, k x_1, \ldots, k x_n\} \) is a partition of \([ka, kb]\), and \(t\) is a step function whose integral is  
	\[
	\int_{ka}^{kb} t(x) \, \mathrm{d}x
	= \sum_{i=1}^{n} s_i \cdot (k x_i - k x_{i-1})
	= k \sum_{i=1}^{n} s_i \cdot (x_i - x_{i-1})
	= k \int_a^b s(x) \, \mathrm{d}x.
	\]		
	\end{proof}
	\begin{mydef}
		If $a>b$, then
		\[
		\int_a^b s(x) \, \mathrm{d}x := -\int_b^a s(x)\, \mathrm{d}x.
		\]
	\end{mydef}
	\subsection{Integral for More General Functions}
	\begin{mydef}
		If $f$ is a real function, we say $f$ is bounded on $[a,b]$ if there exists a number $M>0$ such that  
		\[
		|f(x)|\le M \quad \text{for every}\ x\in [a,b].
		\]
	\end{mydef}
	\begin{mydef}
		If $f$ is a function defined and bounded on $[a,b]$, and $s,t$ are arbitrary step functions satisfying
		\[ s(x)\le f(x)\le t(x) \quad \text{for every}\ x\in[a,b].\]
		And for every pair of $s$ and $t$, there exists one and only one number $I$ such that
		\[\int_a^b s(x)\,\mathrm{d}x\le I \le \int_a^b t(x)\,\mathrm{d}x \quad \text{for every}\ x\in[a,b],\]
		then $I$ is defined as the integral of $f$ and denoted by
		\[ I=\int_a^b f(x)\,\mathrm{d}x.\]
		When such $I$ exists, we say $f$ is integrable on $[a,b]$.
	\end{mydef}
	\begin{myrmk}
		Note that this definition is quite similar to exhaustion property. We will see they work together in the following text.
	\end{myrmk}
	\begin{mydef}
		Let $f$ be a real function bounded on $[a,b]$ and $s,t$ be two step functions satisfying $s\le f\le t$, we define
		\[
		S:=\left\{\int_a^b s(x)\,\mathrm{d}x\, \middle|\, s\le f\right\},
		\quad
		T:=\left\{\int_a^b t(x)\,\mathrm{d}x\, \middle|\, f\le t\right\}.
		\]
		Since $f$ is bounded, $S$ and $T$ are nonempty sets. By $s\le t$, we have 
		\[
		\int_a^b s(x)\,\mathrm{d}x \le \sup S\le \inf T\le \int_a^b t(x)\,\mathrm{d}x.
		\]
		We say $\sup S$ the lower integral of $f$ and $\inf T$ the upper integral of $f$ denoted by
		\[
		\underline{I}(f) = \sup \left\{ \int_a^b s(x) \, \mathrm{d}x \,\middle|\, s \le f \right\}, 
		\qquad 
		\overline{I}(f) = \inf \left\{ \int_a^b t(x) \, \mathrm{d}x \,\middle|\, f \le t \right\}.
		\]
	\end{mydef}
	\begin{mythm}
		Every function $f$ which is bounded on $[a,b]$ has a lower integral and a upper integral satisfying
		\[
		\int_a^b s(x) \, \mathrm{d}x 
		\le \underline{I}(f) 
		\le \overline{I}(f) 
		\le \int_a^b t(x) \, \mathrm{d}x.
		\]
		for all step functions $s$ and $t$ with $s\le f\le t$. The function $f$ is integrable if and only if when
		\[
		\int_a^b f(x) \, \mathrm{d}x = \underline{I}(f) = \overline{I}(f).
		\]
	\end{mythm}
	\begin{myrmk}
		Seemingly, all bounded functions are integrable. But actually this is not true, otherwise we won't introduce such complex notations above. Here's a counter example showing bounded functions can also be non-integrable.
	\end{myrmk}
	\begin{myex}
		Recall the definition of Dirichlet Function. We shall confine it on $[0,1]$. By the dense property of rational numbers and irrational numbers, it's easy to prove
		\[
		\underline{I}(f) = 0 < 1 = \overline{I}(f) \quad \text{so} \quad \underline{I}(f) \neq \overline{I}(f).
		\]
		Hence, the Dirichlet Function is non-integrable.
	\end{myex}
	\begin{mythm}
	Let $f$ be a nonnegative function, integrable on an interval $[a,b]$, and let $Q$ denote the ordinate set of $f$ over $[a,b]$. Then $Q$ is measurable and
	\[
	A(Q)=\int_a^b f(x)\,\mathrm{d}x.
	\]
	\end{mythm}
	\begin{myrmk}
		Note that we only say ordinate sets of some particular functions are measurable, but at this stage we don't know what kind of ordinate sets are not measurable.
	\end{myrmk}
	\begin{mypro}
		Let
		\[
		\widetilde{Q}:=\left\{(x,y)\in \mathbb{R}^2\ |\ a\le x \le b\ and\ 0\le y< f(x)\right\},
		\]
		then $\widetilde{Q}$ is also measurable and $A\!\left(\widetilde{Q}\right)=A(Q)$.
	\end{mypro}
	\begin{mypro}
		Let 
		\[
		S:= \left\{(x,y)\in \mathbb{R}^2\ |\ a\le x \le b\ and\ y=f(x)\right\},
		\]
		then $S$ is also measurable and $A(S)=0$.
	\end{mypro}
	
	\begin{myrmk}
		After we define the integral of more general functions, there comes two questions naturally: which bounded functions are integrable, and given that a function $f$ is integrable, how do we compute the integral of $f$? 
		
		For the first question, we will only give partial answers requiring only elementary ideas because a complete answer needs more prerequisites. For the second question, we'll introduce some techniques. 
		
	\end{myrmk}
	\begin{mydef}
		A function $f$ is said to be increasing on a set $S$ if $f(x)\le f(y)$ for every pair of points $x$ and $y$ in $S$ with $x<y$. If the strict equality $f(x)<f(y)$ holds for all $x<y$ in $S$, then the function $f$ is said to be strictly increasing on $S$. Similarly, we can define decreasing functions and strictly decreasing functions. 
	\end{mydef}
	\begin{mydef}
		A function is called monotonic if it's increasing or decreasing. Similarly, a function is calle strictly monotonic if it's strictly increasing or decreasing. Furthermore, a function $f$ is called piecewise monotonic if there is a partition $P$ on $[a,b]$ such that $f$ is monotonic on every open subintervals of $P$.
	\end{mydef}
	\begin{myrmk}
		Note that a monotonic function needn't to be ``continuous" intuitively.
	\end{myrmk}
	\begin{mythm}
		If a function $f$ is monotonic on $[a,b]$, then $f$ is integrable.
	\end{mythm}
	\begin{proof}
		We shall prove the theorem for increasing functions. Assume $f$ is increasing on $[a,b]$ and let $P=\left\{x_0,x_1,\cdots,x_n\right\}$ be a partition satisfying
		\[
		x_k - x_{k-1} = \frac{b-a}{n}, \quad k=1,2,\cdots,n.
		\]
		and $s,t$ be step functions satisfying:
		\[
		s(x)=f(x_{k-1})\quad \text{and}\quad t(x)=f(x_k),\quad \text{if}\quad x_{k-1}<x<x_k,\quad k=1,2,\cdots,n.
		\]
		At the subdivision points, we define $s_n$ and $t_n$ so as to preserve the relation $s_n(x)\le f(x)\le t_n(x)$ throughout $[a,b]$. 
		It suffices to prove
		\[
		\underline{I}(f) = \overline{I}(f).
		\]
		First we compute
		\begin{align*}
			\int_a^b t(x)\, \mathrm{d}x - \int_a^b s(x)\,\mathrm{d}x
			&= \sum_{k=1}^n f(x_k)(x_k-x_{k-1}) - \sum_{k=1}^n f(x_{k-1})(x_k-x_{k-1})\\[4pt]
			&= \frac{b-a}{n}\left(f(b)-f(a)\right) \\[4pt]
			&= \frac{C}{n}, \quad \text{where} \quad C:=(b-a)(f(b)-f(a))>0
		\end{align*}
		Since we have
		\[
		\int_a^b s(x)\,\mathrm{d}x \le \underline{I}(f)\le 
		\overline{I}(f) \le \int_a^b t(x)\, \mathrm{d}x,
		\]
		thus
		\[
		0 \le \overline{I}(f) - \underline{I}(f) 
		\le \int_a^b t(x)\, \mathrm{d}x - \int_a^b s(x)\,\mathrm{d}x
		= \frac{C}{n}
		\]
		holds for a positive $C$ and every positive integer $n$. Hence we deduce
		\[
		\underline{I}(f) = \overline{I}(f).
		\]
	\end{proof}
	\begin{myrmk}
		Note that the proof for decreasing functions is analogous because it suffices to exchange $s(x)$ and $t(x)$. Here we see the magic power of Archimedian properties of real numbers again.
	\end{myrmk}
	\begin{mythm}
		Assume $f$ is increasing on $[a,b]$. Let
		\[
		x_k = a + \frac{b-a}{n}\cdot k,\quad k=0,1,\cdots,n.
		\]
		If $I$ is any number satisfying the inequalities
		\[
		\frac{b-a}{n}\sum_{k=0}^{n-1}f(x_k) \le I \le
		\frac{b-a}{n}\sum_{k=1}^{n}f(x_k),
		\]
		for every integer $n\ge 1$,	then
		\[
		I = \int_a^b f(x)\,\mathrm{d}x.
		\]
	\end{mythm}
	\begin{myrmk}
		This inspires us to compute the integral of a general function by approaching from above and below.
	\end{myrmk}
	\begin{mythm}
		Here are some properties of integrals of general functions. If not stated, $f$ and $g$ are functions integrable on $[a,b]$.
		\begin{enumerate}
			\item Linearity: Suppose $c_1,c_2$ constants, then
			\[
			\int_a^b \left[ c_1 f(x) + c_2 g(x) \right] \, \mathrm{d}x
			= c_1 \int_a^b f(x) \, \mathrm{d}x
			+ c_2 \int_a^b g(x) \, \mathrm{d}x.
			\]
			\item Additivity:
			\[
			\int_a^b f(x) \, \mathrm{d}x
			+ \int_b^c f(x) \, \mathrm{d}x
			= \int_a^c f(x) \, \mathrm{d}x.
			\]
			\item Invariance under translation: For every real $c$,
			\[
			\int_a^b f(x) \, \mathrm{d}x
			= \int_{a+c}^{b+c} f(x - c) \, \mathrm{d}x.
			\]
			\item Expansion or Contraction: If $k\neq 0$,
			\[
			\int_a^b f(x) \, \mathrm{d}x
			= \frac{1}{k} \int_{ka}^{kb} f\!\left( \frac{x}{k} \right) \, \mathrm{d}x.
			\]
			\item Comparison: If $g(x)\le f(x)$ for every $x$ in $[a,b]$,
			\[
			\int_a^b g(x) \, \mathrm{d}x
			\le \int_a^b f(x) \, \mathrm{d}x.
			\]
		\end{enumerate}
	\end{mythm}
	\begin{proof}
		We shall prove linearity only, since proofs of other properties are analogous. We decompose the linearity into two parts:
		\begin{enumerate}
			\item[(A)] \[ \int_a^b\left[f(x)+g(x)\right]\,\mathrm{d}x=\int_a^b f(x)\,\mathrm{d}x + \int_a^b g(x)\,\mathrm{d}x
			\]
			\item[(B)] \[\int_a^b cf(x)\,\mathrm{d}x=c\int_a^b f(x)\,\mathrm{d}x\]
		\end{enumerate}
		To prove (A), let
		\[
		I(f)=\int_a^b f(x)\,\mathrm{d}x,\quad 
		I(g)=\int_a^b g(x)\,\mathrm{d}x
		\].
		We shall prove
		\[
		\underline{I}(f+g) = \overline{I}(f+g) = I(f)+I(g).
		\]
		Let $s_1,s_2$ be arbitrary step functions below $f,g$ respectively, then
		\[
		I(f)=\sup \left\{ \int_a^b s_1(x)\,\mathrm{d}x \,\middle|\, s_1\le f\right\},\quad
		I(g)=\sup \left\{ \int_a^b s_2(x)\, \mathrm{d}x \,\middle|\, s_2\le g\right\}.
		\]
		So
		\[
		\begin{aligned}
			I(f)+I(g)
			&=\sup \left\{ \int_a^b s_1(x)\,\mathrm{d}x \,\middle|\, s_1\le f\right\} +
			\sup \left\{ \int_a^b s_2(x)\, \mathrm{d}x \,\middle|\, s_2\le g\right\}\\[6pt]
			&=\sup \left\{\int_a^b \left[s_1(x)+s_2(x)\right]\,\mathrm{d}x\,\middle|\,s_1\le f,s_2\le g \right\}
		\end{aligned}
		\]
		Since $s_1\le f, s_2\le g$, we have $s_1(x)+s_2(x) \le f(x)+g(x)$. Thus
		\[
		\underline{I}(f+g)\ \text{is an upper bound of }\ 
		\left\{\int_a^b \left[s_1(x)+s_2(x)\right]\,\mathrm{d}x\,\middle|\,s_1\le f,s_2\le g \right\}.
		\]
		And an upper bound cannot be less than the least upper bound, i.e. the supremum, so
		\[
		I(f)+I(g) \le \underline{I}(f+g).
		\]
		Similarly we can deduce that
		\[
		I(f)+I(g) \ge \overline{I}(f+g).
		\]
		Together with 
		\[
		\underline{I}(f+g) \le \overline{I}(f+g),
		\]
		we reach the equality
		\[
		\underline{I}(f+g) = \overline{I}(f+g) = I(f+g).
		\]
		To prove (B), it's trivial when $c=0$. Now we focus on the case when $c>0$. First we shall prove one property of supremum and infemum: if $c$ is a positive real number, then we have
		\[
		\sup\left\{cx\,|\,x\in A \right\} = c\cdot\sup\left\{ x\,|\,x\in A\right\},
		\quad \inf \left\{cx\,|\,x\in A \right\} = c\cdot\inf \left\{ x\,|\,x\in A\right\}.
		\]
		By definition,
		\[
		\forall x \in A,\ x\le \sup A \Rightarrow cx\le c\cdot \sup A.
		\]
		Assume $c\cdot \sup A$ is not the supremum of the set $\left\{cx\,|\,x\in A \right\}$, since it's a nonempty set and bounded above, there exists a positive $h$ such that $c\cdot \sup A-h$ is the supremum, then
		\[
		\forall x\in A, cx \le c\cdot \sup A-h 
		\Rightarrow x \le \sup A-\frac{h}{c}.
		\]
		Then the right-hand side of the inequality should be the new supremum of $A$, which leads to a contradiction. Thus we prove
		\[
		\sup\left\{cx\,|\,x\in A \right\} = c\cdot\sup\left\{ x\,|\,x\in A\right\}.
		\]
		It's similar to prove infimum's property. Back to our proof, let $s,s_1,t,t_1$ be step functions with $s_1=cs,t_1=ct$. Then
		\[
		\begin{aligned}
			\underline{I}(cf) &= \sup \left\{ \int_a^b s_1(x) \, \mathrm{d}x \,\middle|\, s_1 \le cf \right\}
			=\sup \left\{ c\int_a^b s(x) \, \mathrm{d}x \,\middle|\, s \le f \right\}
			=cI(f)\\[6pt]
			\overline{I}(cf) &= \inf \left\{ \int_a^b t_1(x) \, \mathrm{d}x \,\middle|\, t_1 \le cf \right\}
			=\inf \left\{ c\int_a^b t(x) \, \mathrm{d}x \,\middle|\, t \le f \right\}
			=cI(f).
		\end{aligned}
		\]
		Therefore, $\underline{I}(cf)=\overline{I}(cf)=cI(f)$. For the case when $c<0$, by using
		\[
		\sup\left\{cx\,|\,x\in A \right\} = c\cdot\inf\left\{ x\,|\,x\in A\right\},
		\quad \inf \left\{cx\,|\,x\in A \right\} = c\cdot\sup \left\{ x\,|\,x\in A\right\},
		\]
		it's analogous to finish the proof.
	\end{proof}
	
	\begin{mythm}
		If $p$ is a positive integer and $b>0$, we have
		\[
		\int_0^b x^p \, \mathrm{d}x = \frac{b^{p+1}}{p + 1}.
		\]
		Moreover, we can deduce the equality also holds for $b\le 0$. Therefore, it can be generalized that
		\[
		\int_a^b x^p \, \mathrm{d}x = \frac{b^{p+1} - a^{p+1}}{p + 1},\qquad a,b\in \mathbb{R},\ p\in \mathbb{Z}^+.
		\]
		We can also use the special symbol
		\[
		P(x)\Big|_a^b :=P(b)-P(a).
		\]
		Then the integral could be rewritten as
		\[
		\int_a^b x^p \, \mathrm{d}x
		= \frac{x^{p+1}}{p + 1} \Big|_a^b
		= \frac{b^{p+1} - a^{p+1}}{p + 1}.
		\]
	\end{mythm}
	\begin{mydef}
		Let $f,g$ be two functions and $f\le g$ on $[a,b]$. We call such set $S$ the region between two graphs if $S$ is defined as
		\[
		S:=\left\{(x,y)\in \mathbb{R}^2\ |\ a\le x \le b,f(x)\le y\le g(x) \right\}.
		\]
	\end{mydef}
	\begin{mythm}
		Assume two functions $f,g$ are integrable with $f\le g$ on $[a,b]$. Then the region $S$ between their graphs is measurable and its area $A(S)$ is given by the integral
		\[
		A(S)=\int_a^b [g(x)-f(x)]\,\mathrm{d}x.
		\]
	\end{mythm}
	\begin{myrmk}
		This theorem didn't require $f,g$ to be nonnegative.
	\end{myrmk}
	\begin{mythm}
		 If $p$ is a positive integer and $a,b>0$, then we have
		 \[
		 \int_a^b x^{\frac{1}{p}} \, \mathrm{d}x = 
		 \left(\frac{p}{p+1} \right)\left(b^{\frac{p+1}{p}}- a^{\frac{p+1}{p}}\right).
		 \]
	\end{mythm}
	\newpage
	\section{Applications of Integration}
	Every time before we do the integration, we should ask one thing: is the to-be-integrated function integrable?
	\subsection{More about Area}
	\begin{mydef}
		A circular disk of radius $r$ is the set of all points inside or on the boundary of a circle of radius $r$.
	\end{mydef}
	\begin{mythm}
		A circular disk is measurable and its area $A(r)$ is given by the integral
		\[
		A(r) = 2\int_{-r}^r \sqrt{r^2-x^2}\,\mathrm{d}x.
		\]
	\end{mythm}
	\begin{mydef}
		We define the number $\pi$ to be the area of a unit disk, i.e., the radius of which is exactly $1$. Symbolically: 
		\[
		\pi := 2\int_{-1}^1 \sqrt{1-x^2}\,\mathrm{d}x.
		\]
	\end{mydef}
	\begin{mypro}
		The area of a circular disk with radius $r$ is given by
		\[
		A(r)=\pi r^2.
		\] 
	\end{mypro}
	\begin{mythm}
		Let $f$ be nonnegative and integrable on $[a,b]$ and let $S$ be its ordinate set. We define a new function $g$ denoted by
		\[
		g(x):=kf\left( \frac{x}{k}\right) \qquad \text{with}\ x\in[ka,kb],\, k>0
		\]
		and its ordinate set $kS$. Then the area of $kS$ is given by
		\[
		A(kS)=k^2A(S).
		\]
	\end{mythm}
	\begin{myrmk}
		Actually this new $g$ is ``similar" to $f$ so we call this definition a similarity transformation.
	\end{myrmk}
	\subsection{Integration formulas for the sine and cosine}
	We omit deriving properties of sine and cosine in this lecture note since I'm lazy and I'm pretty sure you know them very well during your study in high school. I shall only introduce one property as you might know but cannot come up with due to less usage.
	\begin{mythm}
		For $0<x<\frac{1}{2}\pi$, we have
		\[
		0<\cos x< \frac{\sin x}{x} < \frac{1}{\cos x}.
		\]
	\end{mythm}
	\begin{mypro}
		Here are several inequalities important in deriving integral formulas
		for sine and cosine.
		\begin{enumerate}
			\item For any $x\in\mathbb R$,
			\begin{align*}
				2\sin\frac{x}{2}\sum_{k=1}^{n}\cos(kx)=\sin\!\left(\Big(n+\frac12\Big)x\right)-\sin\frac{x}{2}.
			\end{align*}
			
			\item For any $\theta\in\mathbb R$,
			\begin{align*}
				2\sin\theta\sum_{k=1}^{n}\cos(2k\theta)=\sin(2n+1)\theta-\sin\theta
				\\
				2\sin\theta\sum_{k=0}^{n-1}\cos(2k\theta)=\sin(2n-1)\theta+\sin\theta.
			\end{align*}
			
			\item Suppose $\theta$ satisfies $0<2n\theta<\tfrac12\pi$, we have
			\begin{align*}
				\sin(2n+1)\theta-\sin\theta<\frac{\sin\theta}{\theta}\sin(2n\theta).
			\end{align*}
			
			\item Suppose $\theta$ satisfies $0<2n\theta<\tfrac12\pi$. First we have
			\begin{align*}
				\frac{\cos(n-1)\theta}{\cos(n\theta)}>\frac{\sin\theta}{\theta}
			\end{align*}
			and then we can deduce
			\begin{align*}
				\frac{\sin\theta}{\theta}\sin(2n\theta)<\sin(2n-1)\theta+\sin\theta.
			\end{align*}
		\end{enumerate}
	\end{mypro}
	\begin{mythm}
		For any $0<a\le \tfrac12\pi$, we have
		\begin{align*}
			\frac{a}{n}\sum_{k=1}^{n}\cos\frac{ka}{n}<\sin a
			<
			\frac{a}{n}\sum_{k=0}^{n-1}\cos\frac{ka}{n}.
		\end{align*}
	\end{mythm}
	\begin{mythm}
		For every real $a$ we have
		\begin{align*}
			\int_0^a \cos x\,\mathrm{d}x = \sin a,\qquad
			\int_0^a \sin x\, \mathrm{d}x = 1-\cos a.
		\end{align*}
		More generally, for every real $a,b$, we have
		\begin{align*}
			\int_a^b \cos x\,\mathrm{d}x = \sin x\,\Big|^b_a\, ,\qquad
			\int_a^b \sin x\, \mathrm{d}x = (-\cos x)\, \Big|^b_a\, .
		\end{align*}
	\end{mythm}

	\subsection{Polar Coordinates}
	\begin{mydef}
		The polar coordinates of a point $P$ in the plane are an ordered pair $(r,\theta)$ where $r$ is the distance from $P$ to the origin $O$ and $\theta$ is the angle formed by the line segment $OP$ and the positive $x$-axis. Here $r\ge 0$ and $\theta\in \mathbb{R}$. They are related to the rectangular coordinates $(x,y)$ of $P$ by the equations
		\[
		x=r\cos \theta,\quad y=r\sin \theta.
		\]
		Here we call $r$ the radial distance of $P$ and $\theta$ a polar angle.
	\end{mydef}
	\begin{myrmk}
		Note that we call $\theta$ ``a'' polar angle instead of ``the'' polar angle because $\theta+2n\pi$ also represents the same angle for any integer $n$. But the radial distance $r$ is uniquely determined. 
		
		If $r=0$, then by convention $\theta$ can be any real number.
	\end{myrmk}
	\begin{mydef}
		A curve is said to be given in polar coordinates if there is a nonnegative function $f$ defined on an interval $I$ such that every point $P$ on the curve has polar coordinates $(r,\theta)$ satisfying
		\[
		r=f(\theta),\quad \theta\in I.
		\]
		And this equation is called the polar equation of the curve. Similarly, the graph of $f$ in polar coordinates is the set of all points $P$ with polar coordinates $(r,\theta)$ satisfying the above equation.
	\end{mydef}

	\begin{myrmk}
		In some textbooks, the variable $r$ in polar coordinates is allowed to be negative: when $r<0$, the point $(r,\theta)$ is identified with the point $(-r,\theta+\pi)$.  Note that the points $(r,\theta)$ and $(-r,\theta)$ are symmetric with respect to the origin. The additional definition allows the existence of inner loops in the graph of a polar equation. But in this lecture note, we only consider nonnegative functions in polar coordinates.
	\end{myrmk}
	\begin{myex}
		Sketch the curves for
		\[
		f(\theta)=1+c\sin\theta\quad \text{where}\quad\theta\in[0,2\pi],\ c\in\{-2,2\}
		\]
		under two different defintion and compare them.
	\end{myex}
	\begin{mydef}
		Let $f$ be a nonnegative function defined on an interval $[a,b]$, where $0\le b-a\le 2\pi$. The set of all points $P$ in the plane with polar coordinates $(r,\theta)$ satisfying
		\[
		0\le r \le f(\theta),\quad \theta\in [a,b]
		\]
		is called the radial set determined by $f$.
	\end{mydef}
	\begin{mydef}
		Similarly as in rectangular coordinates, a nonnegative function $s$ defined on $[a,b]$ is called a step function in polar coordinates if there exists a partition $P=\left\{\theta_0,\theta_1,\cdots,\theta_n\right\}$ of $[a,b]$ such that
		\[
		s(\theta)=s_k\quad\text{where}\quad \theta \in (\theta_{k-1},\theta_k),\quad k=1,2,\cdots,n.
		\]
		And the area of this sector is given by
		\[
		A=\frac12 s_k^2(\theta_k-\theta_{k-1}).
		\]
		Since $b-a<2\pi$, there is no overlapping between different sectors. The area of the radial set determined by $s$ is given by
		\[
		A(S)=\frac12 \sum_{k=1}^n s_k^2(\theta_k-\theta_{k-1})=\frac12 \int_a^b s^2(\theta)\,\mathrm{d}\theta.
		\]
	\end{mydef}
	\begin{mythm}
		Let $f$ be a nonnegative function on $[a,b]$, where $0\le b-a\le 2\pi$. Assume the radial set $S$ determined by $f$ is measurable. If $f^2$ is integrable, then the area $A(S)$ is given by the integral
		\[
		A(S)=\frac12 \int_a^b f^2(\theta)\,\mathrm{d}\theta.
		\]
	\end{mythm}


	\subsection{Volume}
	Similarly, we can use axiomatic definitions to define volume like we did for area.

	\begin{mydef}
		We assume there exists a class $\mathcal{A}$ of solids and a set function $V$, whose domain is $\mathcal{A}$, with the following properties:
		\begin{enumerate}
			\item Nonnegative property: For each set $S$ in $\mathcal{A}$ we have $V(S) \ge 0$.
			
			\item Additive property: If $S$ and $T$ are in $\mathcal{A}$, then 
			$S \cup T$ and $S \cap T$ are in $\mathcal{A}$, and we have
			\[
				V(S \cup T) = V(S) + V(T) - V(S \cap T).
			\]
			
			\item Difference property: If $S$ and $T$ are in $\mathcal{A}$ with $S \subseteq T$, 
			then $T - S$ is in $\mathcal{A}$, and we have 
			\[
				V(T - S) = V(T) - V(S).
			\]
			
			\item Cavalieri's principle.: If $S$ and $T$ are two Cavalieri solids in $\mathcal{A}$ 
			with $A(S \cap F) \le A(T \cap F)$ for every plane $F$ perpendicular to a given line, 
			then $V(S) \le V(T)$.
			
			\item Choice of scale.: Every box $B$ is in $\mathcal{A}$. 
			If the edges of $B$ have lengths $a, b,$ and $c$, then 
			\[
				V(B) = abc.
			\]
			
			\item Every convex set is in $\mathcal{A}$.
		\end{enumerate}
	\end{mydef}
	\begin{mydef}[Additional]
		Some definitions are lost above and here we add it.
		\begin{enumerate}
			\item A solid $S$ in $\mathcal{A}$ is called a \textbf{Cavalieri solid} if for every plane $F$ perpendicular to a given line $L$, the area $A(S\cap F)$ of the cross section of $S$ by $F$ is well-defined, i.e., the cross section is measurable.
			\item A \textbf{box} or \textbf{rectangular parallelepiped} is any set congruent to a set of the form
			\[
			\{(x,y,z)\in \mathbb{R}^3\ |\ 0\le x \le a,\ 0\le y \le b,\ 0\le z \le c \},
			\]
			where nonnegative real numbers $a,b,c$ are called the lengths of the edges of the box.
			\item A solid $S$ in $\mathcal{A}$ is called a \textbf{convex set} if for every two points $P,Q\in S$, the line segment $PQ$ is also contained in $S$.
			\item A plane set is call \textbf{bounded} if it is a subset of some square in the plane.
		\end{enumerate}
	\end{mydef}
	\begin{mypro}
		Here's some propositions for some basic cases.
		\begin{enumerate}
			\item The empty set $\varnothing$ is in $\mathcal{A}$ and $V(\varnothing)=0$.
			\item Monotonicity property: If $S$ and $T$ are in $\mathcal{A}$ with $S\subseteq T$, then $V(S)\le V(T)$.
			\item Every bounded plane set $S$ in $\mathcal{A}$ has volume zero, i.e., $V(S)=0$.
		\end{enumerate}
	\end{mypro}
	\begin{myrmk}
		Note that the Cavalieri's principle can be applied twice: if $A(S\cap F)=A(T\cap F)$ for every plane $F$ perpendicular to a given line, then $V(S)=V(T)$. In Chinese, we call this property Zu Geng's principle.
	\end{myrmk}
	\begin{mydef}
		A \textbf{right cylindrical solid} is a set congruent to a set $S$ of the form
		\[
		S=\left\{(x,y,z)\in \mathbb{R}^3\ |\ (x,y)\in B,\, a\le z \le b \right\},
		\]
		where $B$ is a bounded plane set and $a,b$ are real numbers with $a\le b$. The base of the right cylindrical solid is the set $B$, and its height is $b-a$.
	\end{mydef}
	\begin{mythm}
		Let $S$ be a right cylindrical solid with base $B$ and height $h$. Then $S$ is in $\mathcal{A}$ and its volume is given by
		\[
		V(S)=A(B)\cdot h=\int_a^b A_S(z)\,\mathrm{d}z
		\]
		More generally, let $R$ be a Cavalieri solid in $\mathcal{A}$ with a cross-sectional area function $A_R$ which is integrable on an interval $[a,b]$ and zero outside $[a,b]$. Then the volume of $R$ is given by
		\[
		V(R)=\int_a^b A_R(z)\,\mathrm{d}z.
		\]
	\end{mythm}
	
	\begin{mythm}
	Let $f$ and $g$ be nonnegative integrable functions on $[a,b]$ satisfying $g(x) \le f(x)$ for all $x \in [a,b]$. If the region bounded by $y = f(x)$ and $y = g(x)$ is revolved about the $x$-axis, then the resulting solid of revolution (if it belongs to $\mathcal{A}$) has volume
	\[
		V = \pi \int_a^b \bigl[f^2(x) - g^2(x)\bigr] \, \mathrm{d}x .
	\]
	In particular, if $g(x) = 0$, the volume is
	\[
		V = \pi \int_a^b f^2(x) \, \mathrm{d}x .
	\]
	\end{mythm}
	\begin{myex}
		 Let $f(x) = \sqrt{r^2 - x^2}$ on $[-r,r]$, the solid is a sphere of radius $r$, and
		\[
		V = \pi \int_{-r}^{r} (r^2 - x^2)\, \mathrm{d}x = \frac{4}{3}\pi r^3.
		\]
	\end{myex}

	\subsection{Indefinite Integral}
	\begin{mydef}
		Let $f$ be a real-valued function integrable on $[a,b]$. The \textbf{average value} of $f$, denoted by $\operatorname{avg}(f)$, on $[a,b]$ is defined to be
		\[
		\operatorname{avg}(f)=\frac{1}{b-a}\int_a^b f(x)\,\mathrm{d}x.
		\]
	\end{mydef}
    \begin{myrmk}
		Note that the definition of average value of a function is consistent with the definition of average value of a finite set of numbers. You might not see how this definition works in this section, but it will be useful in later section.
	\end{myrmk}
    \begin{mythm}
        Assume $f$ and $g$ are continuous on $[a,b]$. If $g$ never changes sign on $[a,b]$, then for some $c\in[a,b]$, we have
        \[
        \int_a^b f(x)g(x)\,\mathrm{d}x = f(c)\int_a^b g(x)\,\mathrm{d}x.
        \]
        This is called the \textbf{Weighted Mean Value Theorem for Integrals}.
        Let $g(x)\equiv 1$, we have the \textbf{Mean Value Theorem for Integrals}:
        \[
        \int_a^b f(x)\,\mathrm{d}x = f(c)(b-a).
        \]
    \end{mythm}
    \begin{myrmk}
        The proof is in Chapter 4.
    \end{myrmk}
	\begin{mydef}
		Let $f$ be a real-valued function integrable on every closed interval contained in an open interval $I$. An \textbf{indefinite integral} of $f$ on $I$ is the function $F$ defined by
		\[
		F(x)=\int_a^x f(t)\,\mathrm{d}t,\quad x\in I,
		\]
		where $a$ is an arbitrary point in $I$.
	\end{mydef}
	\begin{myrmk}
		Actually a different choice of $a$ will only lead to a different indefinite integral by a constant, since
		\[
		\int_b^x f(t)\,\mathrm{d}t = \int_a^x f(t)\,\mathrm{d}t - \int_a^b f(t)\,\mathrm{d}t.
		\]
	\end{myrmk}
	\begin{mythm}
		Suppose $f$ integrable and nonnegative on $[a,b]$, then the indefinite integral $A(f)$ is increasing.
	\end{mythm}
	\begin{mydef}
		A function $g$ is said to be \textbf{convex} on an interval $[a,b]$ if for every $x_1,x_2\in [a,b]$ and $0\le \lambda \le 1$, we have
		\[
		g\left(\lambda x_1 + (1-\lambda)x_2\right) \le \lambda g(x_1) + (1-\lambda)g(x_2).
		\]
		Similarly, $g$ is said to be \textbf{concave} on $[a,b]$ if the above inequality is reversed, i.e.,
		\[
		g\left(\lambda x_1 + (1-\lambda)x_2\right) \ge \lambda g(x_1) + (1-\lambda)g(x_2).
		\]
	\end{mydef}
	\begin{myrmk}
		Geometrically, a function $g$ is convex on $[a,b]$ if the line segment joining any two points on the graph of $g$ lies above or on the graph between these two points. Similarly, $g$ is concave on $[a,b]$ if the line segment joining any two points on the graph of $g$ lies below or on the graph between these two points.
	\end{myrmk}
	\begin{mythm}
		Let $f$ be integrable on $[a,b]$ and let $F$ be an indefinite integral of $f$ on $[a,b]$. If $f$ is increasing on $[a,b]$, then $F$ is convex on $[a,b]$. If $f$ is decreasing on $[a,b]$, then $F$ is concave on $[a,b]$.
	\end{mythm}





	\newpage

	\section{Limits and Continuity of functions}
	\subsection{Basic Concepts}
	\begin{mydef}
		Any open interval containing $p\in \mathbb{R}$ as the midpoint is called a \textbf{neighborhood} of $p$. Suppose $r>0$, then
		\[
		N(p)=(p-r,p+r)=\{ x\in \mathbb{R} : |x-p|<r \}
		\]
		is a neighborhood of $p$ with \textbf{radius} $r$. Moreover, a \textbf{punctured neighborhood} is the deleted neighborhood given by removing its midpoint.
	\end{mydef}
	\begin{mynot}
		We can also denote $N(p)$ as $N(p\,;r)$.
	\end{mynot}
	\begin{mydef}
		Let $f$ be a function defined in some neighborhood of $p$ except possibly not defined at $x=p$. Suppose $A\in \mathbb{R}$. The symbol
		\[
		\lim_{x\to p} f(x)=A \quad \left[\text{or} \quad f(x)\rightarrow A\quad \text{as}\quad x\to p\right] 
		\]
		reads ``the limit of $f(x)$ as $x$ approaches $p$ is equal to $A$". The symbol means that for every neighborhood $N_1(A)$, there exists some neighborhood $N_2(p)$ such that
		\[
		f(x)\in N_1(A)\quad \text{whenever}\quad x\in N_2(p)\quad \text{and}\quad x\neq p.
		\]
	\end{mydef}
	\begin{myrmk}
		In this definition, we only consider $A$ to be a real number and $f$ to be a real-valued function. Besides, the limit has nothing to do with what happens at $x=p$. 
	\end{myrmk}
	\begin{mydef}
		Equivalently, the limit can be defined as
		\[
		\forall \varepsilon >0,\ \exists \delta>0,\ |f(x)-A|<\varepsilon\ \text{whenever}\ 0<|x-p|<\delta,
		\]
		which is called the ``$\varepsilon-\delta$ definition of limit".
	\end{mydef}
	\begin{myrmk}
		Actually you've seen and you will see infinitely many ``whenever" which all mean the same thing: ``for all ...".
	\end{myrmk}
	\begin{myrmk}
		The following symbols are all equivalent:
		\begin{align*}
			&\lim_{x\to p} f(x)=A, \\
			&\lim_{x\to p} \left[ f(x)-A\right]=0, \\
			&\lim_{x\to p} |f(x)-A|=0,\\
			&\lim_{h\to 0} f(p+h)=A, \\
			&\lim_{h\to 0} f(p-h)=A.
		\end{align*}
	\end{myrmk}
	\begin{mydef}
		Let $f$ be a function defined on $(p,b)$ for some $b>p$ and $A\in \mathbb{R}$. The symbol
		\[
		\lim_{x\to p^+}f(x)=A
		\]
		means that for every neighborhood $N_1(A)$, there exists some neighborhood $N_2(p)$ such that
		\[
		f(x)\in N_1(A)\quad \text{whenever}\quad x\in N_2(p)\quad \text{and}\quad x>p,
		\]
		which is called the \textbf{right-side limit}.
		
		Similarly we can define what is called the \textbf{left-side limit}.
		Let $f$ be a function defined on $(a,p)$ for some $a<p$ and $A\in \mathbb{R}$. The symbol
		\[
		\lim_{x\to p^-}f(x)=A
		\]
		means that for every neighborhood $N_1(A)$, there exists some neighborhood $N_2(p)$ such that
		\[
		f(x)\in N_1(A)\quad \text{whenever}\quad x\in N_2(p)\quad \text{and}\quad x<p.
		\]
	\end{mydef}
	\begin{myrmk}
		If we have
		\[
		\lim_{x\to p}f(x)=A,
		\]
		then it's easy to verify that both
		\[
		\lim_{x\to p^+}f(x)=A\quad \text{and}\quad \lim_{x\to p^-}f(x)=A
		\]
		hold. Furthermore, the converse still holds.
	\end{myrmk}
	
	
	\begin{mydef}
		A function $f$ is said to be \textbf{continuous} at $p\in \mathbb{R}$ if it satisfies two conditions:
		\begin{enumerate}
			\item $f$ is defined at $p$.
			\item $\lim\limits_{x\to p}f(x)=f(p).$
		\end{enumerate}
	\end{mydef}
	
	\begin{mydef}
		Let $f$ be a function and $a\in\mathbb{R}$. We say that $f$ has a \textbf{removable discontinuity} at $a$ if the two-sided limit
		\[
		\lim_{x\to a} f(x)=L\in\mathbb{R}
		\]
		exists and is finite, but either $f(a)$ is not defined or $f(a)\neq L$. Equivalently, redefining $f(a):=L$ makes $f$ continuous at $a$.
	\end{mydef}
	
	\begin{myex}
		Let
		\[
		f(x):=
		\begin{cases}
			1, & x \neq 0,\\
			0,  & x = 0.
		\end{cases}
		\]
		Then $f$ has a removable discontinuity at $x=0$.
	\end{myex}
	
	\begin{mydef}
		Let $f$ be a function and $a\in\mathbb{R}$. We say that $f$ has a \textbf{jump discontinuity} at $a$ if both one-sided limits exist and are finite,
		\[
		\lim_{x\to a^-} f(x)=L_- \in \mathbb{R},\qquad
		\lim_{x\to a^+} f(x)=L_+ \in \mathbb{R},
		\]
		but $L_- \neq L_+$. The jump size is $|L_+-L_-|$.
	\end{mydef}
	
	\begin{myex}
		Let
		\[
		f(x)=x-\left[ x\right],
		\]
		then $f$ has a jump discontinuity at every $x=k$ where $k$ is an integer.
	\end{myex}
	
	\begin{mydef}
		Let $f$ be a function and $a\in\mathbb{R}$. We say that $f$ has an \textbf{infinite discontinuity} at $a$ if either
		\[
		\lim_{x\to a^-} f(x)=\pm\infty \quad \text{or} \quad
		\lim_{x\to a^+} f(x)=\pm\infty
		\]
		holds. Equivalently, $f$ is unbounded on every punctured neighborhood of $a$.
	\end{mydef}
	
	\begin{myex}
		Let
		\[
		f(x):=
		\begin{cases}
			\dfrac{1}{x^2}, & x \neq 0,\\[6pt]
			0,  & x = 0.
		\end{cases}
		\]
		Then $\lim\limits_{x\to p}f(x)$ \textbf{does not exist} and so for the left-side and right-side limit. Though we haven't define what ``$\infty$" is, we should learn it intuitively yet.
	\end{myex}

	\begin{mydef}
		Let $f$ be a function and $a\in\mathbb{R}$. 
		We say that $f$ has an \textbf{oscillatory discontinuity} at $a$ if at least one of the one-sided limits
		\[
		\lim_{x\to a^-} f(x), \qquad \lim_{x\to a^+} f(x)
		\]
		fails to exist, and the failure is \emph{not} because the limit diverges to 
		$+\infty$ or $-\infty$. Equivalently, $f$ oscillates without approaching any 
		finite limit on every punctured neighborhood of $a$.
	\end{mydef}

	\begin{myex}
		Let 
		\[
		f(x)=
		\begin{cases}
			\sin\left(\dfrac{1}{x}\right), & x\neq 0,\\[4pt]
			0, & x=0.
		\end{cases}
		\]
		Then $f$ has an oscillatory discontinuity at $x=0$. Indeed,
		neither one-sided limit 
		\[
		\lim_{x\to0^-}\sin\left(\dfrac{1}{x}\right), \qquad 
		\lim_{x\to0^+}\sin\left(\dfrac{1}{x}\right)
		\]
		exists, because the function oscillates infinitely often between $-1$ and $1$
		as $x\to 0$ without approaching any limiting value.
	\end{myex}



	\begin{mythm}
		Let $f$ and $g$ be functions such that 
		\[
		\lim_{x\to p} f(x)=A,\quad \lim_{x\to p} g(x)=B,
		\]
		Then we have the following properties:
		\begin{enumerate}
			\item $\displaystyle \lim_{x\to p} [f(x)+g(x)]=A+B$,
			\item $\displaystyle \lim_{x\to p} [f(x)-g(x)]=A-B$,
			\item $\displaystyle \lim_{x\to p} [f(x)\cdot g(x)]=A\cdot B$,
			\item If $B\neq 0$, then $\displaystyle \lim_{x\to p} \frac{f(x)}{g(x)}=\frac{A}{B}$.
		\end{enumerate}
	\end{mythm}
	\begin{mythm}
		Let $f$ and $g$ be continuous at $p$. Then \[f+g,\, f-g,\, fg \] are also continuous at $p$. If we know $g(p)\neq 0$, then $\dfrac{f}{g}$ is also continuous at $p$.
	\end{mythm}
	\begin{mythm}
		Suppose that $f(x)\le g(x)\le h(x)$ for all $x\neq p$ in some neighborhood $N(p)$ and additionally
		\[
		\lim_{x\to p} f(x)=\lim_{x\to p} h(x)=A,
		\]
		then
		\[
		\lim_{x\to p} g(x)=A.
		\]
		This is called the \textbf{squeeze theorem}.
	\end{mythm}
	\begin{proof}
		For every $\varepsilon>0$, there exists some neighborhood $N_1(p)$ such that
		\[
		f(x)\in N(A\,;\varepsilon),\quad h(x)\in N(A\,;\varepsilon)
		\]
		whenever $x\in N_1(p)$ and $x\neq p$. Let $N_2(p)$ be the neighborhood such that
		\[
		f(x)\le g(x)\le h(x)
		\]
		holds whenever $x\in N_2(p)$ and $x\neq p$. Now let $N_3(p)=N_1(p)\cap N_2(p)$, then for every $x\in N_3(p)$ and $x\neq p$, we have
		\[
		A-\varepsilon < f(x) \le g(x) \le h(x) < A+\varepsilon,
		\]
		i.e., $g(x)\in N(A\,;\varepsilon)$. Thus by definition we have
		\[
		\lim_{x\to p} g(x)=A.
		\]
	\end{proof}
	\begin{myrmk}
		When you choose $f$ and $h$ to apply the squeeze theorem, they're usually not unique as long as they work well.
	\end{myrmk}
	\begin{mythm}
		Assume $f$ is integrable on $[a,x]$ for every $x\in [a,b]$ and define
		\[
		F(x)=\int_a^x f(t)\,\mathrm{d}t,\quad x\in[a,b].
		\]
		Then $F$ is continuous at every $x\in(a,b)$.
	\end{mythm}
	\begin{myrmk}
		This theorem provides us a way to prove the continuity of some complicated functions constructed by integrals.

		We can replace ``continuous at every $x\in(a,b)$'' in the conclusion by ``continuous at every $x\in[a,b]$'' with the following definition of one-sided continuity.
	\end{myrmk}
	\begin{mydef}
		A function $f$ is said to be \textbf{right-side continuous} or \textbf{continuous from the right} at $p$ if it satisfies two conditions:
		\begin{enumerate}
			\item $f$ is defined at $p$.
			\item $\lim\limits_{x\to p^+}f(x)=f(p).$
		\end{enumerate}
		
		Similarly, a function $f$ is said to be \textbf{left-side continuous} or \textbf{continuous from the left} at $p$ if it satisfies two conditions:
		\begin{enumerate}
			\item $f$ is defined at $p$.
			\item $\lim\limits_{x\to p^-}f(x)=f(p).$
		\end{enumerate}
	\end{mydef}
	\begin{myrmk}
		Moreover, the above theorems all hold for one-sided limits and one-sided continuity.
	\end{myrmk}
	\begin{myex}
		Here are some basic examples.
		\begin{enumerate}
			\item Note that
			\[
			\sin x = \int_0^x \cos t\,\mathrm{d}t,\quad
			\cos x = 1 - \int_0^x \sin t\,\mathrm{d}t,
			\]
			we can conclude that both $\sin x$ and $\cos x$ are continuous on $\mathbb{R}$.
			\item Define
			\[
			f(x):=
			\begin{cases}
				\dfrac{\sin x}{x}, & x \neq 0,\\[6pt]
				1,  & x = 0.
			\end{cases}
			\]
			Then $f$ is continuous at $x=0$ since
			\[
			0<\cos x<\frac{\sin x}{x}<\frac{1}{\cos x}\quad \text{for}\quad 0<|x-0|<\frac{\pi}{2}.
			\]
			and by squeeze theorem we have
			\[
			\lim_{x\to 0} \frac{\sin x}{x}=1=f(0).
			\]
			Moreover, $f$ is continuous on $\mathbb{R}$.
		\end{enumerate}
	\end{myex}
	\subsection{Major Theorems for Continuous functions}
	\begin{mythm}
		Assume $v$ is continuous at $p$ and $u$ is continuous at $q=v(p)$. Then the composite function $f=u\circ v$ is continuous at $p$.
	\end{mythm}
	\begin{mylma}
		Let $f$ be continuous at $c$ and suppose that $f(c)\neq 0$.
		Then there exists $\delta>0$ such that for the neighborhood $N(c\,;\delta)$ we have
		\[
		\forall x\in N(c\,;\delta),\ f(x)\,f(c)>0 .
		\]
	\end{mylma}
	\begin{proof}
		Since $f$ is continuous at $c$, for $\varepsilon = \dfrac{|f(c)|}{2}>0$, there exists some $\delta>0$ such that
		\[
		|f(x)-f(c)|<\varepsilon = \frac{|f(c)|}{2}
		\]
		holds whenever $x\in N(c\,;\delta)$. Thus we have
		\[
		-\frac{|f(c)|}{2} < f(x)-f(c) < \frac{|f(c)|}{2},
		\]
		i.e.,
		\[
		f(c)-\frac{|f(c)|}{2} < f(x) < f(c)+\frac{|f(c)|}{2}.
		\]
		If $f(c)>0$, then we have
		\[
		\frac{f(c)}{2} < f(x) < \frac{3f(c)}{2},
		\]
		which implies $f(x)>0$. If $f(c)<0$, then we have
		\[
		\frac{3f(c)}{2} < f(x) < \frac{f(c)}{2},
		\]
		which implies $f(x)<0$. In both cases, we have $f(x)f(c)>0$.
	\end{proof}
	\begin{mythm}
		Let $f$ be a function continuous on the closed interval $[a,b]$. If $f(a)$ and $f(b)$ have opposite signs, then there exists some $c\in (a,b)$ such that
		\[
		f(c)=0.
		\]
		This is called the \textbf{Bolzano's Theorem}.
	\end{mythm}
	\begin{proof}
		Without loss of generality, we assume $f(a)<0$ and $f(b)>0$. Let
		\[
		S=\{ x\in [a,b] : f(x)\le 0 \}.
		\]
		Since $a\in S$, $S$ is nonempty. And since $S\subseteq [a,b]$, $S$ is bounded above by $b$. By the completeness axiom of real numbers, let
		\[
		c=\sup S.
		\]
		We shall prove that $f(c)=0$. Suppose on the contrary that $f(c)\neq 0$. Then by the above lemma, there exists some $\delta>0$ such that for every $x\in N(c\,;\delta)$, we have
		\[
		f(x)f(c)>0.
		\]
		If $f(c)>0$, then for every $x\in N(c\,;\delta)$, we have $f(x)>0$. Note that there exists some $x_1\in S$ such that
		\[
		c-\delta < x_1 < c\quad \text{and} \quad f(x_1) > 0,
		\]
		since $c=\sup S$. But we know $x_1\in S$ and thus $f(x_1) \le 0$, which leads to a contradiction.
		
		If $f(c)<0$, then for every $x\in N(c\,;\delta)$, we have $f(x)<0$. Note that there exists some $x_3\in [a,b]\setminus S$ such that
		\[
		c < x_3 < c+\delta, \quad \text{and} \quad f(x_3) < 0,
		\]
		since $c=\sup S$. But we know $x_3\notin S$ and thus $f(x_3) > 0$, which leads to a contradiction.

		Hence we must have $f(c)=0$.
	\end{proof}
	\begin{myrmk}
		One-sided continuity should be considered at the endpoints $a$ and $b$ when applying Bolzano's Theorem. Otherwise, the conclusion might not hold.
	\end{myrmk}
	\begin{mythm}
		Let $f$ be continuous at each point of a closed interval $[a,b]$. For two arbitrary points $x_1,x_2\in[a,b]$ with $f(x_1)\neq f(x_2)$ and any number $L$ between $f(x_1)$ and $f(x_2)$, there exists some $c$ between $x_1$ and $x_2$ such that
		\[
		f(c)=L.
		\]
		This is called the \textbf{Intermediate Value Theorem (IVT)}.
	\end{mythm}
	\begin{myrmk}
		The Intermediate Value Theorem is a immediate consequence of Bolzano's Theorem. And such $c$ above might not be unique.
	\end{myrmk}
	\begin{mypro}
		If $n$ is a positive integer and if $a>0$, then there is exactly one positive real number $x$ such that
		\[
		x^n = a.
		\]
	\end{mypro}
	\begin{myrmk}
		Of course we can derive the existence and uniqueness of $n$th root of any real number by using the completeness axiom of real numbers. But here using the Intermediate Value Theorem is more efficient.
	\end{myrmk}
	\begin{mythm}
		Assume $f$ is strictly increasing and continuous on an interval $[a,b]$. Let $c=f(a)$ and $d=f(b)$ and let $g$ be the inverse function of $f$ defined on $[c,d]$. Then
		\begin{enumerate}
			\item $g$ is strictly increasing on $[c,d]$,
			\item $g$ is continuous on $[c,d]$.
		\end{enumerate}
	\end{mythm}
	\begin{proof}
		\begin{enumerate}
			\item For any $y_1,y_2\in [c,d]$ with $y_1<y_2$, let
			\[
			x_1=g(y_1),\quad x_2=g(y_2).
			\]
			Since $f$ is strictly increasing, we have
			\[
			y_1=f(x_1)<f(x_2)=y_2,
			\]
			which implies that $x_1<x_2$. Thus $g$ is strictly increasing on $[c,d]$.
			\item Choose arbitrary $y_0\in (c,d)$. To prove $g$ is continuous at $y_0$, we need to show that for every $\varepsilon>0$, there exists some $\delta>0$ such that
			\[
			|g(y)-g(y_0)|<\varepsilon\quad \text{whenever}\quad |y-y_0|<\delta.
			\]
			Let 
			\[
			x_0=g(y_0),\quad f(x_0)=y_0.
			\]
			Suppose $\varepsilon>0$ is given. Since $f$ is increasing, we let \[\delta=\min\{f(x_0+\varepsilon)-f(x_0), f(x_0)-f(x_0-\varepsilon)\}>0.\] Then we have
			\[
			f(x_0 - \varepsilon) < y < f(x_0 + \varepsilon) \quad \text{whenever}\quad |y-y_0|<\delta.
			\]
			Since $g$ is strictly increasing, we have
			\[
			x_0 - \varepsilon < g(y) < x_0 + \varepsilon,
			\]
			which implies that
			\[
			|g(y)-g(y_0)|<\varepsilon.
			\]

		\end{enumerate}
	\end{proof}
	\begin{myrmk}
		Suppose $f(x)=x^2$ on $[-a,a]$ where $a>0$ is not invertible, but $f(x)$ is invertible on $[-a,0]$ or on $[0,a]$. In some textbooks, we express ``the inverse'' of $f(x)$ on $[0,a^2]$ as ``double function''.
	\end{myrmk}

	\begin{mydef}
		Suppose $f$ is a real-valued function on a set $S\subseteq \mathbb{R}$. Then $f$ is said to have an \textbf{absolute maximum} on $S$ if
		\[
		\exists c \in S\ \forall x\in S,\ f(x)\le f(c).
		\]
		And $f(c)$ is calle the \textbf{absolute maximum value} of $f$ on $S$. Similarly, $f$ is said to have an \textbf{absolute minimum} on $S$ if
		\[
		\exists d \in S\ \forall x\in S,\ f(x)\ge f(d).
		\]
		And $f(d)$ is calle the \textbf{absolute minimum value} of $f$ on $S$.
	\end{mydef}
	\begin{mylma}
		Suppose $f$ is continuous on $[a,b]$, then $f$ is bounded on $[a,b]$.
	\end{mylma}
	\begin{proof}
		We shall prove by contradiction and using the method of \textbf{successive bisection}. Suppose $f$ is unbounded on $[a,b]$. Let
		\[ c= \frac{a+b}{2},\]
		then we know among the two subintervals $[a,c]$ and $[c,b]$, at least on one of the two we know $f$ is unbounded. Pick the subinterval on which $f$ is unbounded (if $f$ is unbounded on both, then we pick the left one, i.e., $[a,c]$) and denote it by $[a_1,b_1]$. We continue this process and keep bisecting the chosen subintervals to get a sequence of subintervals:
		\[
		[a_1,b_1]\supseteq [a_2,b_2]\supseteq \cdots [a_n,b_n] \supseteq \cdots
		\]
		where $n$ is any positive integer and on which $f$ is unbounded. Note that the length of $[a_n,b_n]$ is given by
		\[
		b_n - a_n = \frac{b-a}{2^n}
		\]
		Let
		\[
		A:=\{a,a_1,a_2,\cdots,a_n,\cdots\}\neq \varnothing
		\]
		and $A$ is bounded above by $b$. Then we let $\alpha=\sup A$. Since $\alpha \in [a,b]$ and $f$ is continuous at $\alpha$, there exists some $\delta>0$ such that given $\varepsilon = 1$, we have
		\[
		|f(x)-f(\alpha)|<1\quad \text{whenever}\ 
		\begin{cases}
		x\in (\alpha-\delta, \alpha + \delta),\ &\text{if}\ \alpha\in [a,b]\\[4pt]
		x \in [\alpha, \alpha + \delta),\ &\text{if}\ \alpha=a\\[4pt]
		x \in (\alpha - \delta, \alpha],\ &\text{if}\ \alpha=b
		\end{cases}
		\]
		By triangle inequality we have
		\[
		|f(x)|-|f(\alpha)| \le |f(x)-f(\alpha)|<1.
		\]
		So we know $f$ is bounded on
		\(
		\begin{cases}
			(\alpha-\delta, \alpha + \delta), & \text{if}\ \alpha\in (a,b)\\[4pt]
			[\alpha, \alpha + \delta), & \text{if}\ \alpha=a\\[4pt]
			(\alpha - \delta, \alpha], & \text{if}\ \alpha=b
		\end{cases}
		\)
		
		\noindent Since
		\[
		a\le a_1 \le a_2 \le \cdots \le a_n \le \cdots,
		\]
		we know there exists some $k$ larger enough such that
		\[
		\begin{cases}
			[a_k,b_k] \subseteq (\alpha-\delta, \alpha + \delta), & \text{if}\ \alpha\in (a,b)\\[4pt]
			[a_k,b_k] \subseteq [\alpha, \alpha + \delta), & \text{if}\ \alpha=a\\[4pt]
			[a_k,b_k] \subseteq (\alpha - \delta, \alpha], & \text{if}\ \alpha=b
		\end{cases}
		\]
		This means $f$ unbounded on $[a_k,b_k]$ but
		\[
		|f(x)| \le 1+|f(\alpha)|,\quad \forall x\in [a_k,b_k],
		\]
		which is a contradiction. Thus $f$ is bounded on $[a,b]$.
	\end{proof}

	\begin{mythm}
		Suppose $f$ is continuous on $[a,b]$, then there exists $c\in [a,b]$ and $d\in [a,b]$ such that
		\[
		\forall x \in [a,b] \quad f(d)\le f(x)\le f(c).
		\]
		This is called the \textbf{Extreme Value Theorem (EVT)}.
	\end{mythm}
	\begin{proof}
	By the above lemma, since $f$ is continuous on $[a,b]$, it is bounded on $[a,b]$.
	Thus the set 
	\[
	S := \{f(x): x\in[a,b]\}
	\]
	is a nonempty bounded subset of $\mathbb{R}$, and therefore has a supremum
	\[
	M := \sup S.
	\]

	We first prove that $f$ attains the value $M$. Suppose, towards a contradiction,
	that $f(x)<M$ for every $x\in[a,b]$. Then for all $x\in[a,b]$ the quantity
	$M-f(x)$ is positive. Define a new function
	\[
	g(x):=\frac{1}{M-f(x)},\qquad x\in[a,b].
	\]
	Because $f$ is continuous and $M-f(x)>0$ everywhere, the function $g$ is also continuous and bounded on $[a,b]$; that is,
	there exists $K>0$ such that
	\[
	g(x)\le K,\qquad \forall x\in[a,b].
	\]
	Substituting the definition of $g$, we obtain
	\[
	\frac{1}{M-f(x)}\le K
	\qquad\Longrightarrow\qquad
	f(x)\le M-\frac{1}{K},\qquad \forall x\in[a,b].
	\]
	Thus $M-\dfrac{1}{K}$ is an upper bound for $S$, but 
	\[
	M-\frac{1}{K}<M,
	\]
	which contradicts the definition of $M=\sup S$. Therefore our assumption was 
	false, and there must exist $c\in[a,b]$ such that $f(c)=M$.

	To obtain the minimum, apply the same argument to the continuous function 
	$-f$ on $[a,b]$. There exists $d\in[a,b]$ such that
	\[
	-f(d)=\sup\{-f(x):x\in[a,b]\}
	\quad\Longrightarrow\quad
	f(d)=\min\{f(x):x\in[a,b]\}.
	\]

	Thus $f$ attains both its maximum and minimum on $[a,b]$, completing the proof.
	\end{proof}
	\begin{myrmk}
		Actually if we define 
		\[
		\sup f:=\sup \{f(x):x\in [a,b]\}\quad \text{and} \quad
		\inf f:=\inf \{f(x):x\in [a,b]\}
		\]
		then in the above theorem we have
		\[
		c=\sup f,\quad d=\inf f.
		\]
	\end{myrmk}
	\begin{mypro}
		Combining \textbf{EVT} and \textbf{IVT}, we have the following : If $f$ is continuous on $[a,b]$, then the range of $f$ on $[a,b]$ is $[\inf f, \sup f]$.
	\end{mypro}
	\begin{mydef}
		Suppose $f$ is continuous on $[a,b]$, we call $(\sup f - \inf f)$ the \textbf{leap} or \textbf{span} of $f$ on $[a,b]$.
	\end{mydef}
	\begin{mypro}
		If $f$ is continuous on $[a,b]$ and $[c,d]\subseteq [a,b]$, then the leap on $[c,d]$ cannot exceed the leap of $f$ on $[a,b]$.
	\end{mypro}
	\begin{mythm}
		Suppose $f$ is continuous on $[a,b]$. Then for any $\varepsilon>0$, there exists a partition $P=\{x_0,x_1,\cdots,x_n\}$ of $[a,b]$ such that the leap of $f$ in every closed interval $[x_{k-1},x_k]$ is less than $\varepsilon$.

		This is called the \textbf{the Theorem of uniform continuity} or \textbf{the small leap/span Theorem}.
	\end{mythm}
	\begin{proof}
		We shall prove by contradiction. Recall how we use the method of successive bisection to prove that a continuous function on a closed interval is bounded. Now we do similarly to choose a sequence of subintervals:
		\[
		[a_1,b_1]\supseteq [a_2,b_2]\supseteq \cdots [a_n,b_n] \supseteq \cdots
		\]
		where $n$ is any positive integer and on which the leap of $f$ is at least $\varepsilon_0$. Let 
		\[
		\alpha = \sup \{a,a_1,a_2,\cdots,a_n,\cdots\} \in [a,b]
		\]
		with
		\[
		a\le a_1 \le a_2 \le \cdots \le a_n \le \cdots.
		\]
		By continuity, there exists some $\delta>0$ such that in an interval
		\[
		\begin{cases}
		(\alpha-\delta, \alpha + \delta), & \text{if}\ \alpha\in (a,b)\\[4pt]
		[\alpha, \alpha + \delta), & \text{if}\ \alpha=a\\[4pt]
		(\alpha - \delta, \alpha], & \text{if}\ \alpha=b
		\end{cases}
		\]
		the leap of $f$ is less than $\varepsilon_0$. Note that the length of $[a_n,b_n]$ is given by
		\[
		b_n - a_n = \frac{b-a}{2^n}.
		\]
		So there exists some $k$ larger enough such that
		\[
		\begin{cases}
			[a_k,b_k] \subseteq (\alpha-\delta, \alpha + \delta), & \text{if}\ \alpha\in (a,b)\\[4pt]
			[a_k,b_k] \subseteq [\alpha, \alpha + \delta), & \text{if}\ \alpha=a\\[4pt]
			[a_k,b_k] \subseteq (\alpha - \delta, \alpha], & \text{if}\ \alpha=b
		\end{cases}
		\]
		This means the leap of $f$ on $[a_k,b_k]$ is less than $\varepsilon_0$, which is a contradiction. Thus we can keep bisecting until the leap of $f$ on every subinterval is less than $\varepsilon_0$.
	\end{proof}
	\begin{mythm}
		Suppose $f$ is continuous on $[a,b]$, then $f$ is integrable on $[a,b]$.
	\end{mythm}
	\begin{proof}
	Obviously $f$ is bounded on $[a,b]$, so $f$ has both an upper integral $\overline{I}(f)$ and a lower integral $\underline{I}(f)$. We shall prove that
	\[
	\overline{I}(f)=\underline{I}(f),
	\]
	which implies that $f$ is integrable on $[a,b]$.

	Choose an integer $N\ge 1$ and let $\varepsilon=1/N$. By the small--span theorem for continuous functions, for this $\varepsilon$ there exists a partition
	\[
	P=\{x_0, x_1,\dots,x_n\}
	\]
	of $[a,b]$ into $n$ subintervals such that the span of $f$ in every subinterval	is less than $\varepsilon$. Denote by $M_k(f)$ and $m_k(f)$ the absolute maximum
	and minimum values of $f$ on the $k$th subinterval $[x_{k-1},x_k]$. Then
	\[
	M_k(f)-m_k(f) < \varepsilon
	\qquad \text{for each } k=1,\dots,n.
	\]
	Now define two step functions $s_n$ and $t_n$ on $[a,b]$ by
	\[
	s_n(x)=
	\begin{cases}
	m_k(f), & \text{if } x_{k-1}<x\le x_k,\\[4pt]
	m_1(f), & x=a,
	\end{cases}
	\ 
	t_n(x)=
	\begin{cases}
	M_k(f), & \text{if } x_{k-1}\le x < x_k,\\[4pt]
	M_n(f), & x=b.
	\end{cases}
	\]
	Then clearly
	\[
	s_n(x)\le f(x)\le t_n(x), \qquad \forall\, x\in[a,b].
	\]
	Moreover,
	\[
	\int_a^b s_n
	= \sum_{k=1}^n m_k(f)(x_k-x_{k-1}), 
	\qquad
	\int_a^b t_n
	= \sum_{k=1}^n M_k(f)(x_k-x_{k-1}).
	\]
	The difference of these two integrals is
	\[
	\int_a^b t_n - \int_a^b s_n
	< \varepsilon \sum_{k=1}^n (x_k-x_{k-1})
	= \varepsilon(b-a)
	= \frac{b-a}{N}.
	\]
	On the other hand, the upper and lower integrals of $f$ satisfy
	\[
	\int_a^b s_n \le \underline{I}(f) \le \overline{I}(f) 
	\le \int_a^b t_n.
	\]
	Hence we obtain
	\[
	0 \le \overline{I}(f) - \underline{I}(f)
	\le \int_a^b t_n - \int_a^b s_n
	< \frac{b-a}{N}.
	\]
	Since this holds for every integer $N\ge 1$, we must have
	\[
	\overline{I}(f)=\underline{I}(f).
	\]
	Thus $f$ is integrable on $[a,b]$.
	\end{proof}
	\begin{mythm}
		If $f$ is continuous on $[a,b]$, then
		\[
		\int_a^b f(x)\, \dd x= f(c)(b-a) \quad \text{for some}\quad c\in [a,b],
		\]
		i.e., there exists $c\in [a,b]$ such that
		\[
		f(c)=\operatorname{avg}(f)=\frac{1}{b-a}\int_a^b f(x)\, \dd x.
		\]
	\end{mythm}
	\begin{proof}
		Let $m=\inf f$ and $M=\sup f$. Since $m\le f(x) \le M$ for all $x\in [a,b]$, we have
		\[
		m=\frac{1}{b-a}\int_a^b m\, \dd x \le \frac{1}{b-a}\int_a^b f(x)\, \dd x \le \frac{1}{b-a}\int_a^b M\, \dd x = M.
		\]
		By \textbf{EVT} and \textbf{IVT}, we know there exists $c\in [a,b]$ such that $f(c)=\operatorname{avg}(f)$.
	\end{proof}
    \begin{myrmk}
        You can also expand the above theorem to what is so-called \textbf{Weighted Mean Value Theorem for Integrals} mentioned in Chapter 3. The proof is left to you as an exercise (since it's quite analogous to the above proof).
    \end{myrmk}
	\newpage
	\section{Differential Calculus}
	\subsection{Basic Concepts}
	\begin{mydef}
		Let $f$ be a function. If
		\[
		\lim _{h\to 0} \frac{f(x+h)-f(x)}{h}=L \in \mathbb{R},
		\]
		then we call it the \textbf{derivative} of $f$ and denote it by	$f'(x)$. Furthermore, the $n$-th derivative of $f$ is defined by
		\[
		f^{(n)}(x)= \left[f^{(n-1)}(x)\right]' \quad \text{for}\quad n=1,2,3,\cdots
		\]
		and $f^{(0)}(x)=f(x)$.
	\end{mydef}

	\begin{myex}
		We're going to prove some basic derivatives using the definition.
		\begin{enumerate}
			\item Suppose $n$ is a positive integer and $x\ge 0$. Consider $f(x)=x^n$, then
			\[
			\begin{aligned}
				f'(x) &= \lim_{h\to 0} \frac{(x+h)^n - x^n}{h}\\ &= \lim_{h\to 0} \sum_{k=0}^{n-1} (x+h)^k x^{n-1-k}\\ &= \sum_{k=0}^{n-1} x^k x^{n-1-k}\\ &= n x^{n-1}.
			\end{aligned}
			\]
			\item Consider $f(x)=\sin x$, then
			\[
			\begin{aligned}
				f'(x) &= \lim_{h\to 0} \frac{\sin(x+h) - \sin x}{h}\\[6pt] 
				&= \lim_{h\to 0} \dfrac{2\sin\left(\dfrac{h}{2}\right)\cos\left(x+\dfrac{h}{2}\right)}{h}\\[6pt] 
				&= \left[ \lim_{h \to 0} \dfrac{\sin\left(\dfrac{h}{2}\right)}{\dfrac{h}{2}}\right] \left[\lim_{h \to 0} \cos\left(x+\dfrac{h}{2}\right)\right]  \\ 
				&= \cos x.
			\end{aligned}
			\]
			\item Consider $f(x)=\cos x$, then
			\[
			\begin{aligned}
				f'(x) &= \lim_{h\to 0} \frac{\cos(x+h) - \cos x}{h}\\[6pt] 
				&= \lim_{h\to 0} \dfrac{-2\sin\left(\dfrac{h}{2}\right)\sin\left(x+\dfrac{h}{2}\right)}{h}\\[6pt] 
				&= -\left[ \lim_{h \to 0} \dfrac{\sin\left(\dfrac{h}{2}\right)}{\dfrac{h}{2}}\right] \left[\lim_{h \to 0} \sin\left(x+\dfrac{h}{2}\right)\right]  \\ 
				&= -\sin x.
			\end{aligned}
			\]
			\item Suppose $n$ is a positive integer and $x>0$. Consider $f(x)=x^{\frac{1}{n}}$ and for convenience let 
			\[
			u=\left( x+h \right)^{\frac{1}{n}} \quad\text{and}\quad v=x^{\frac{1}{n}}.
			\]
			Hence we have
			\[
			\begin{aligned}
				f'(x) &= \lim_{h\to 0} \frac{(x+h)^{\frac{1}{n}} - x^{\frac{1}{n}}}{h}\\[6pt] 
				&= \lim_{h\to 0} \dfrac{u - v}{u^n - v^n} \\[6pt] 
				&= \lim_{u\to v} \dfrac{1}{u^{n-1} + u^{n-2}v + \cdots + uv^{n-2} + v^{n-1}} \\[6pt]
				&= \dfrac{1}{n x^{\frac{n-1}{n}}}\\[6pt]
				&= \dfrac{1}{n} x^{\frac{1}{n}-1}.
			\end{aligned}
			\]
			For the case $x=0$, we can see $f(x)$ is defined at $x=0$ but $f'(x)$ is not defined.
		\end{enumerate}
	\end{myex}
	\begin{mythm}
		If function $f$ has a derivative at $x$, then $f$ is continuous at $x$.
	\end{mythm}
	\begin{proof}
		Let
		\[
		f'(x)=\lim_{h\to 0} \frac{f(x+h)-f(x)}{h}.
		\]
		So
		\[
		\lim_{h\to 0} f(x+h) = \lim_{h\to 0} \left[ \frac{f(x+h)-f(x)}{h} \cdot h + f(x) \right] = f'(x)\cdot 0 + f(x) = f(x).
		\]
	\end{proof}
	\begin{mythm}
		Suppose $f$ and $g$ are two functions defined on a common interval. If both $f$ adn $g$ have derivatives at $x=a$, then so do the following functions:
		\[
		f+g,f-g,fg, \frac{f}{g}\ \text{(provided that $g(a)\neq 0$)}.
		\]
		Additionally, their derivatives satisfy the following:
		\begin{enumerate}
			\item[(1)] $(f+g)'=f'+g'$,
			\item[(2)] $(f-g)'=f'-g'$,
			\item[(3)] $(fg)'=f'g+fg'$,
			\item[(4)] $\left(\dfrac{f}{g}\right)'=\dfrac{f'g-fg'}{g^2}$.
		\end{enumerate}
	\end{mythm}
	\begin{proof}
		We only prove (3) and (4) here.
		\begin{enumerate}
			\item[(3)] We have
			\[
			\begin{aligned}
				(fg)'(a) &= \lim_{h\to 0} \frac{f(a+h)g(a+h) - f(a)g(a)}{h}\\[6pt]
				&= \lim_{h\to 0} \left[ g(a+h)\cdot \frac{f(a+h)-f(a)}{h} + f(a)\cdot \frac{g(a+h)-g(a)}{h} \right]\\[6pt]
				&= g(a)\cdot f'(a) + f(a)\cdot g'(a).
			\end{aligned}
			\]
			\item[(4)] We have
			\[
			\begin{aligned}
				\left(\frac{f}{g}\right)'(a) &= \lim_{h\to 0} \frac{\dfrac{f(a+h)}{g(a+h)} - \dfrac{f(a)}{g(a)}}{h}\\[6pt]
				&= \lim_{h\to 0} \frac{f(a+h)g(a) - f(a)g(a+h)}{h g(a+h) g(a)}\\[6pt]
				&= \lim_{h\to 0} \left[ \frac{g(a)\cdot [f(a+h)-f(a)]}{h g(a+h) g(a)} - \frac{f(a)\cdot [g(a+h)-g(a)]}{h g(a+h) g(a)} \right]\\[6pt]
				&= \frac{g(a)\cdot f'(a)}{g^2(a)} - \frac{f(a)\cdot g'(a)}{g^2(a)}\\[6pt]
				&= \frac{f'g - fg'}{g^2}\bigg|_{x=a}.
			\end{aligned}
			\]
		\end{enumerate}
	\end{proof}
	\begin{myrmk}
		Using these properties and induction, we can expand
		\[
		\left( x^n\right)' = n x^{n-1}
		\]
		for all rational numbers $n$.
	\end{myrmk}
	\begin{mynot}
		Actually the following notations are also commonly used for derivatives:
		\begin{enumerate}
			\item Leibniz's notation:
			\[
			\frac{\dd y}{\dd x},\ \frac{\dd^2 y}{\dd x^2},\ \frac{\dd^3 y}{\dd x^3},\ \cdots,\ \frac{\dd^n y}{\dd x^n},
			\]
			where $\dd x$ and $\dd y$ are called \textbf{differentials}.
			\item Lagrange's notation:
			\[
			f',\ f'',\ f''',\cdots,\ f^{(n)};
			\]
			\item Newton's notation (mainly used in physics):
			\[
			\dot{y},\ \ddot{y},\ \dddot{y},\ \cdots;
			\]
			\item Arbogast's notation:
			\[
			Df,\ D^2f,\ D^3f,\ \cdots,\ D^n f,
			\]
			where $D$ is called the \textbf{differential operator}. It's easy to verify that $D$ is linear.
		\end{enumerate}
	\end{mynot}

	\subsection{Chain Rule for Differentiation}
	\begin{mythm}
		Suppose $u$ and $v$ are functions and the function $f$ is given by
		$f=u\circ v$. Suppose $v'(x)$ and $u'(y)$ where $y=v(x)$ both exist. Then $f'(x)$ exists and is given by
		\[
		f'(x) = u'(v(x)) \cdot v'(x).
		\]
	\end{mythm}
	\begin{proof}
		For convenience, let
		\[
		y=v(x),\quad \Delta y = v(x+h) - v(x).
		\]
		% Then if we assume $\Delta y\neq 0$,  we have
		% \[
		% \begin{aligned}
		% 	f'(x) &= \lim_{h\to 0} \frac{u(v(x+h)) - u(v(x))}{h}\\[6pt]
		% 	&= \lim_{h\to 0} \frac{u(y + \Delta y) - u(y)}{h}\\[6pt]
		% 	&= \lim_{h\to 0} \left[ \frac{u(y + \Delta y) - u(y)}{\Delta y} \cdot \frac{\Delta y}{h} \right]\\[6pt]
		% 	&= \left[ \lim_{\Delta y \to 0} \frac{u(y + \Delta y) - u(y)}{\Delta y} \right] \cdot \left[ \lim_{h\to 0} \frac{\Delta y}{h} \right]\\[6pt]
		% 	&= u'(v(x)) \cdot v'(x).
		% \end{aligned}
		% \]
		Let
		\[
		g(t)= 
		\begin{cases}
		0, & \text{if}\ t=0,\\[6pt]
		\dfrac{u(y+t) - u(y)}{t}-u'(y), & \text{if}\ t\neq 0,
		\end{cases}
		\]
		since we know for sure
		\[
		\lim_{t\to 0} g(t) = \lim_{t\to 0} \left[ \frac{u(y+t) - u(y)}{t} - u'(y) \right]=0.
		\]
		Then we have
		\[
		u(y+t)-u(y)=
		\begin{cases}
		0\left[u'(y) + g(t)\right], & \text{if}\ t=0\\[6pt]
		t\left[u'(y) + g(t)\right], & \text{if}\ t\neq 0
		\end{cases}
		= t\left[u'(y) + g(t)\right]
		\]
		no matter $t=0$ or not. Thus plugging in $t=\Delta y$, we know
		\[
		f'(x)=\lim_{h\to 0} \frac{u(y + \Delta y) - u(y)}{h} = \lim_{h\to 0} \frac{\Delta y}{h} \left[u'(y) + g(\Delta y)\right].
		\]
		Note that
		\[
		\lim_{h\to 0} \frac{\Delta y}{h}=v'(x)
		\]
		as discussed before. Also we know $g$ is continuous at $0$, so
		\[
		\lim_{h\to 0} \left[g(\Delta y)+ u'(y)\right] = g(0) + u'(y) = u'(y).
		\]
		Hence we have
		\[
		f'(x) = u'(v(x)) \cdot v'(x).
		\]
	\end{proof}
	\begin{myrmk}
		In Leibniz's notation, suppose $y=v(x)$ and $z=u(y)$, then we write
		\[
		\frac{\dd y}{\dd x}:=v'(x)
		\quad \text{and}\quad
		\frac{\dd z}{\dd y}:=u'(y).
		\]
		We also write
		\[
		\frac{\dd z}{\dd x} := f'(x)
		\quad \text{where}\quad f=u \circ v.
		\]
		Then the chain rule can be written as
		\[
		\frac{\dd z}{\dd x} = \frac{\dd z}{\dd y} \cdot \frac{\dd y}{\dd x} \not\equiv \frac{\dd z}{\bcancel{\dd y}}\cdot \frac{\bcancel{\dd y}}{\dd x}.
		\]
		Here, $\not \equiv$ means ``not identically equal to", which indicates that the differentials $\dd y$ are \emph{not} actually quantities that can be cancelled out.
	\end{myrmk}

	\subsection{Theorems of Derivatives}
	\begin{mydef}
		A function $f$ defined on a set $S \subseteq \mathbb{R}$ is said to have a \textbf{relative maximum} at $c\in S$ if there exists some open interval $I$ containing $c$ such that
		\[
		\forall x \in I \cap S,\ f(x) \le f(c).
		\]
		Similarly, we can define \textbf{relative minimum} of $f$ at $d\in S$.
	\end{mydef}
	\begin{myrmk}
		We also say $f$ has an \textbf{extremum} at $c$ if it has either a relative maximum or a relative minimum at $c$.
	\end{myrmk}
	\begin{mythm}
		Suppose function $f$ is defined on an open interval $(a,b)$ containing $c$. If $f$ has a relative extremum at $c$ and if $f'(c)$ exists, then
		\[
		f'(c)=0.
		\]
	\end{mythm}
	\begin{proof}
		Let
		\[
		g(x)=\begin{cases}
		\dfrac{f(x)-f(c)}{x-c}, & \text{if}\ x \in (a,b), x\neq c,\\[6pt]
		f'(c), & \text{if}\ x=c.
		\end{cases}
		\]
		It suffices to show that $g(c)=0$. Note
		\[
		\lim_{x \to c} \frac{f(x)-f(c)}{x-c}=f'(c)
		\]
		since $f'(c)$ exists, thus $g$ is continuous at $c$. Now suppose $g(c)>0$, then there exists some $\delta>0$ such that
		\[
		g(x)>0,\quad \forall x\in (c-\delta, c+\delta) \subseteq (a,b).
		\]
		If $x\in(c, c+\delta)$, then from $g(x)>0$ we derive $f(x)>f(c)$. If $x\in(c-\delta, c)$, then from $g(x)>0$ we derive $f(x)<f(c)$. Note this contradicts to $f$ having an extremum at $c$. Similarly, we can prove that $g(c)<0$ also leads to a contradiction. Thus we must have $g(c)=0$, i.e., $f'(c)=0$.
	\end{proof}
	\begin{mythm}
		Let $f$ be continuous on $[a,b]$ and differentiable on $(a,b)$. If $f(a)=f(b)$, then there exists some $c\in (a,b)$ such that
		\[
		f'(c)=0.
		\]
		This is called \textbf{Rolle's Theorem}.
	\end{mythm}
	\begin{proof}
		Suppose $f'(x)\neq 0$ for all $x\in (a,b)$. By \textbf{EVT}, we know
		\[
		\sup f = f(\alpha) \quad \text{and}\quad \inf f = f(\beta)\quad \text{for some}\quad \alpha, \beta \in [a,b].
		\]
		Note we know $f$ have extremum at $\alpha$ and $\beta$. If $\alpha, \beta \in (a,b)$, then by the above theorem we have
		\[
		f'(\alpha)=f'(\beta)=0,
		\]
		which contradicts our assumption. Thus $\alpha,\beta\in \{a,b\}$ and $f(\alpha)=f(\beta)$. Thus we know $f$ is constant on $[a,b]$, which implies
		\[
		f'(x)= 0,\quad \forall x\in (a,b),
		\]
		contradicting our assumption again. Then we complete the proof.
	\end{proof}
	\begin{mythm}
		Let $f$ be continuous on $[a,b]$ and differentiable on $(a,b)$. Then there exists some $c\in (a,b)$ such that
		\[
		f'(c)=\frac{f(b)-f(a)}{b-a}.
		\]
		This is called \textbf{Lagrange's Mean Value Theorem} or \textbf{the Mean Value Theorem (MVT)} for derivatives.
	\end{mythm}
	\begin{proof}
		Define a new function
		\[
		g(x) = f(x) - \left[ \frac{f(b)-f(a)}{b-a} (x - a) + f(a) \right].
		\]
		Clearly $g$ is continuous on $[a,b]$ and differentiable on $(a,b)$. Note that
		\[
		\begin{cases}
			g(a)=f(a)-f(a)=0 ,\\[6pt]
			g(b)=f(b)-\left[ \dfrac{f(b)-f(a)}{b-a} (b - a) + f(a) \right]=0.
		\end{cases}
		\]
		By \textbf{Rolle's Theorem}, there exists some $c\in (a,b)$ such that
		\[
		g'(c) = f'(c) - \frac{f(b)-f(a)}{b-a} = 0,
		\]
		which implies
		\[
		f'(c) = \frac{f(b)-f(a)}{b-a}.
		\]
	\end{proof}
	\begin{mythm}
		Suppose $f$ and $g$ are both continuous on $[a,b]$ and differentiable on $(a,b)$, then there exists some $c\in (a,b)$ such that
		\[
		f'(c)\left[ g(b)-g(a) \right] = g'(c) \left[ f(b)-f(a) \right].
		\]
		This is called \textbf{Cauchy's Mean Value Theorem}.
	\end{mythm}
	\begin{proof}
		Define a new function
		\[
		h(x) = f(x) \left[ g(b)-g(a) \right] - g(x) \left[ f(b)-f(a) \right].
		\]
		Clearly $h$ is continuous on $[a,b]$ and differentiable on $(a,b)$. Note that
		\[
		\begin{cases}
			h(a)=f(a)\left[ g(b)-g(a) \right] - g(a)\left[ f(b)-f(a) \right],\\[6pt]
			h(b)=f(b)\left[ g(b)-g(a) \right] - g(b)\left[ f(b)-f(a) \right].
		\end{cases}
		\Rightarrow h(b)-h(a)=0.
		\]
		By \textbf{Rolle's Theorem}, there exists some $c\in (a,b)$ such that
		\[
		h'(c) = f'(c)\left[ g(b)-g(a) \right] - g'(c)\left[ f(b)-f(a) \right] = 0,
		\]
		which implies
		\[
		f'(c)\left[ g(b)-g(a) \right] = g'(c)\left[ f(b)-f(a) \right].
		\]
	\end{proof}
	\begin{mythm}
		Suppose $f$ is continuous on $[a,b]$ and differentiable on $(a,b)$. Then:
		\begin{enumerate}
			\item [(1)] If $f'(x)=0$ for all $x\in (a,b)$, then $f$ is constant on $[a,b]$;
			\item [(2)] If $f'(x)>0$ for all $x\in (a,b)$, then $f$ is strictly increasing on $[a,b]$;
			\item [(3)] If $f'(x)<0$ for all $x\in (a,b)$, then $f$ is strictly decreasing on $[a,b]$.
		\end{enumerate}
	\end{mythm}
	\begin{mythm}
		Suppose $f$ is continuous on $[a,b]$ and differentiable on $(a,b)$ and $c\in (a,b)$. Then:
		\begin{enumerate}
			\item [(1)] If $f'(x)>0$ for all $x\in (a,c)$ and $f'(x)<0$ for all $x\in (c,b)$, then $f$ has a relative maximum at $c$;
			\item [(2)] If $f'(x)<0$ for all $x\in (a,c)$ and $f'(x)>0$ for all $x\in (c,b)$, then $f$ has a relative minimum at $c$.
		\end{enumerate}
	\end{mythm}
	\begin{myrmk}
		The above two cases still hold even if $f'(c)$ is not well-defined. (We still require $f$ to be continuous at $c$.)
	\end{myrmk}
	\begin{mydef}
		Suppose $f$ is continuous on $[a,b]$ and differentiable on $(a,b)$ and $c\in (a,b)$. If $f'(c)=0$, then we call $c$ a \textbf{critical point} of $f$.
	\end{mydef}
	\begin{mythm}
		Suppose $f$ has a second-derivative $f''$ defined on $(a,b)$ and $c\in (a,b)$ is a critical point of $f$. Then:
		\begin{enumerate}
			\item [(1)] If $f''(c)>0$, then $f$ has a relative minimum at $c$;
			\item [(2)] If $f''(c)<0$, then $f$ has a relative maximum at $c$;
			\item [(3)] If $f''(c)=0$, then the test is inconclusive, i.e., $f$ may have a relative maximum, a relative minimum or neither at $c$.
		\end{enumerate}
	\end{mythm}
	\begin{mythm}
		Suppose $f$ is continuous on $[a,b]$ and differentiable on $(a,b)$. Then:
		\begin{enumerate}
			\item If $f'$ is increasing on $(a,b)$, or equivalently $f''(x) \geq 0$ for all $x \in (a,b)$ provided that $f''$ is well-defined, then $f$ is convex on $[a,b]$.
			\item If $f'$ is decreasing on $(a,b)$, or equivalently $f''(x) \leq 0$ for all $x \in (a,b)$ provided that $f''$ is well-defined, then $f$ is concave on $[a,b]$.
		\end{enumerate}
	\end{mythm}
	\begin{myrmk}
		Note that a flat line is both convex and concave (in our text book's definition).
	\end{myrmk}

	\newpage
	\section{Integration and Differentiation}
	\subsection{Two Fundamental Theorems for Calculus}
	\begin{mythm}
		Suppose $f$ is a function that is integrable on $[a,x]$ for every $x\in [a,b]$. Let $c\in [a,b]$ and define function $F$ on $[a,b]$ by
		\[
		F(x) = \int_c^x f(t)\, \dd t.
		\]
		Then for any $x\in (a,b)$, if $f$ is continuous at $x$, then $F'(x)$ exists and
		\[
		F'(x) = f(x)\quad\text{or}\quad \frac{\dd}{\dd x} \left( \int_c^x f(t)\, \dd t \right) = f(x).
		\]
		This is called the \textbf{First Fundamental Theorem of Calculus (FTC1)}.
	\end{mythm}
	\begin{proof}
		Let $x\in (a,b)$ be fixed. Consider the quotient in the definition of $F'(x)$:
		\[
		\begin{aligned}
			\frac{F(x+h)-F(x)}{h} 
			&= \frac{1}{h} \left( \int_c^{x+h} f(t)\, \dd t - \int_c^x f(t)\, \dd t \right) \\
			&= \frac{1}{h} \int_x^{x+h} f(t)\, \dd t\\
			&= \frac{1}{h} \int_x^{x+h} \left[ f(x)+\left(f(t)-f(x)\right)\right] \dd t\\
			&= f(x) + \frac{1}{h} \int_x^{x+h} \left( f(t)-f(x) \right) \dd t
		\end{aligned}
		\]
		Note it suffices to show 
		\[
		G(h):= \frac{1}{h} \int_x^{x+h} \left( f(t)-f(x) \right) \dd t \to 0 \quad \text{as}\quad h\to 0.
		\]
		Given any $\varepsilon>0$, since $f$ is continuous at $x$, there exists some $\delta>0$ such that
		\[
		|f(t)-f(x)| < \frac{\varepsilon}{2}, \quad \forall t\in (x-\delta, x+\delta)\subseteq (a,b).
		\]
		\textbf{Claim:} Suppose $f$ is a function integrable on $[a,b]$, then we have
		\[
		\int_a^b |f(t)| \dd t \geq \left| \int_a^b f(t) \dd t \right|.
		\]
		\begin{proof}
			Note that
			\[
			|f(x)| \ge \max\left\{ f(x), -f(x) \right\}.
			\]
			So we have
			\[
			\int_a^b |f(t)| \dd t \ge \max\left\{\int_a^b f(x)\dd t, -\int_a^b f(x)\dd t\right\} = \left|\int_a^b f(t) \dd t\right|.
			\]
			Thus the claim holds.
		\end{proof}
		\medskip

		\textbf{Case 1:} Suppose $0<h<\delta$, then $[x-h, x+h] \subseteq (x-\delta, x+\delta)$. Thus we have
		\[
		|G(h)| = \left| \frac{1}{h} \int_x^{x+h} \left( f(t)-f(x) \right) \dd t \right| \leq \frac{1}{h} \int_x^{x+h} |f(t)-f(x)| \dd t < \frac{\varepsilon}{2} < \varepsilon.
		\]
		So $G(h)\to 0$ as $h\to 0^+$.

		\textbf{Case 2:} Suppose $-\delta < h < 0$, then analogously we have
		\[
		\lim _{h\to 0^-} G(h) = 0.
		\]
		Together we have $G(h)\to 0$ as $h\to 0$. Thus we complete the proof.
	\end{proof}
	\begin{mydef}
		A function $P$ is called a \textbf{primitive} (or an \textbf{antiderivative}) on a function $f$ on an open interval $I$ if $P'(x)=f(x)$ for any $x\in I$. 
	\end{mydef}
	\begin{mythm}
		Assume $f$ is continuous on an open interval $I$ and let $P$ be a primitive of $f$ on $I$. Then for each $x,c \in I$, we have
		\[
		\int_c^x f(t)\, \dd t = P(x) - P(c).
		\]
		This is called the \textbf{Second Fundamental Theorem of Calculus (FTC2)}.
	\end{mythm}
	\begin{myrmk}
		We have 
		\[
		\int_a^b f(x)\, \dd x = P(b) - P(a)
		\]
		provided that $P$ is a primitive of $f$ on some open interval containing $[a,b]$.
	\end{myrmk}
	\begin{mynot}
		Leibniz introduced the notation \textbf{indefinite integral} to denote a general primitive of $f$ or the common format shared by all primitive of $f$:
		\[
		\int f(x)\, \dd x := \int_c^x f(t)\, \dd t + C
		\]
		where $C$ is an arbitrary constant. In practice, we often just write
		\[
		\int f(x)\, \dd x = P(x) + C
		\]
		where $P(x)$ it self \emph{does not} have any constant term.
	\end{mynot}
	\begin{myrmk}
		In some textbooks, since we often write
		\[
		C=\int_d^c f(t)\, \dd t
		\]
		and then 
		\[
		\int f(x)\, \dd x =\int_c^x f(t)\, \dd t + C = \int_d^x f(t)\, \dd t,
		\]
		the defintion of indefinite integral is generalized so that Leibniz's notation is also called an indefinite integral.
	\end{myrmk}

	\subsection{Integration Techniques}
	\begin{mymthd}
		Suppose $u=g(x)$ is a differentiable function whose range is an interval $I$ and $f$ is continuous on $I$. Then
		\[
		\int f(g(x)) g'(x)\, \dd x = \int f(u)\, \dd u.
		\]
		This is called \textbf{integration by substitution}.
	\end{mymthd}
	\begin{mythm}
		Suppose $g$ has a continuous derivative $g'$ on an open interval $I$. Let $J$ be the set of values taken by $g$ on $I$ and assume that $f$ is continuous on $J$. Then for any $x,c\in I$, we have
		\[
		\int_c^x f(g(t)) g'(t)\, \dd t = \int_{g(c)}^{g(x)} f(u)\, \dd u.
		\]
	\end{mythm}

    \begin{proof}
        Since $g$ is continuous, by \textbf{IVT} we know $J=g(I)$ is a well-defined interval. Thus we can define
        \[
        F(u)=\int _{g(c)}^u f(w)\, \dd w,\quad u \in J.
        \]
        Note that $f$ is continuous on $J$, we know $F$ is differentiable on $J$. By \textbf{FTC1}, we have $F'(u)=f(u)$. Define
        \[
        H(t) = F(g(t)),\quad t\in I
        \]
        with
        \[
        H'(t)=F'(g(t))\,g'(t)=f(g(t))\,g'(t),\quad t\in I.
        \]
        By \textbf{FTC2}, we know
        \[
        \begin{aligned}
            \int_c^x f(g(t))g'(t)\,\dd t &= \int_c^x H'(t)\,\dd t = H(x)-H(c)=F(g(x))-F(g(c)) \\ &= \int_{g(c)}^{g(x)} f(w)\,\mathrm{d}w - \int_{g(c)}^{g(c)} f(w)\,\mathrm{d}w = \int_{g(c)}^{g(x)} f(w)\,\mathrm{d}w.
        \end{aligned}
        \]
        Note that the integral value is independent of integral variable. Hence, we complete the proof.
    \end{proof}
	\begin{mymthd}
		Suppose $u=u(x)$ and $v=v(x)$ are two differentiable functions on an open interval $I$. Then
		\[
		\int u(x) v'(x)\, \dd x = u(x) v(x) - \int v(x) u'(x)\, \dd x.
		\]
		For definite integrals, we have
		\[
		\int_a^b u(x) v'(x)\, \dd x = \left[ u(x) v(x) \right]_a^b - \int_a^b v(x) u'(x)\, \dd x.
		\]
		This is called \textbf{integration by parts}.
	\end{mymthd}
	\begin{myrmk}
		With Leibniz's notation, since equivalently (in a sense of notation but it's rigorous by defining what is so called differential later) we have
		\[
		\dd u = u'(x)\, \dd x \quad \text{and}\quad \dd v = v'(x)\, \dd x,
		\]
		we can write integration by parts as
		\[
		\int u\, \dd v = uv - \int v\, \dd u.
		\]
		Leibniz's notation is often more convenient and aesthetically pleasing.
	\end{myrmk}
	\begin{myex}
		Recall
		\[
		\int x^n\, \dd x = \frac{1}{n+1} x^{n+1} + C,\quad \text{does not hold when}\ n=-1.
		\]
		Now we use integration by parts to try solve the case $n=-1$:
		\[
		\begin{aligned}
			\int \frac{1}{x}\, \dd x &= \int 1 \cdot \frac{1}{x}\, \dd x\\[6pt]
			&= x \cdot \frac{1}{x} - \int x \cdot \left( -\frac{1}{x^2} \right) \dd x + C\\[6pt]
			&= 1 + \int \frac{1}{x}\, \dd x + C.
		\end{aligned}
		\]
		Now we see that the constant term $C$ is actually important, otherwise we would derive something nonsense. And we know it's not a wise choice to use integration by parts here. 
	\end{myex}
	\begin{mythm}
		Suppose $g$ is continuous on $[a,b]$ and $f$ has a continuous derivative $f'$ on $[a,b]$.  If $f'$ does not change sign on $[a,b]$, then there exists some $c\in [a,b]$ such that
		\[
		\int_a^b f(x) g(x)\, \dd x = f(a) \int_a^c g(x)\, \dd x + f(b) \int_c^b g(x)\, \dd x.
		\]
	\end{mythm}
	\begin{proof}
		Let
		\[
		G(x) = \int_a^x g(t)\, \dd t.
		\]
		Since $g$ is continuous on $[a,b]$, by the \textbf{First Fundamental Theorem of Calculus}, we know $G'(x)=g(x)$ for any $x\in (a,b)$. By \textbf{integration by parts}, we have
		\[
		\begin{aligned}
			\int_a^b f(x) g(x)\, \dd x &= \int_a^b f(x) G'(x)\, \dd x\\[6pt]
			&= \left[ f(x) G(x) \right]_a^b - \int_a^b G(x) f'(x)\, \dd x\\[6pt]
			&= f(b) G(b) - f(a) \cdot 0 - \int_a^b G(x) f'(x)\, \dd x\\[6pt]
			&= f(b) \int_a^b g(t)\, \dd t - \int_a^b G(x) f'(x)\, \dd x\\[6pt]
			% &= f(b) \int_a^b g(t)\, \dd t - \int_a^b G(x) f'(x)\, \dd x. 
		\end{aligned}
		\]
		Using the \textbf{Mean Value Theorem}, since $f'$ does not change sign on $[a,b]$, there exists some $c\in [a,b]$ such that
		\[
		\int_a^b G(x) f'(x)\, \dd x = G(c) \int_a^b f'(x)\, \dd x =  \left( \int_a^c g(t)\, \dd t \right) [f(b)-f(a)].
		\]
		Thus we have
		\[
		\begin{aligned}
			\int_a^b f(x) g(x)\, \dd x &= f(b) \int_a^b g(t)\, \dd t - \left( \int_a^c g(t)\, \dd t \right) [f(b)-f(a)]\\[6pt]
			&= f(b) \int_c^b g(t)\, \dd t + f(a) \int_a^c g(t)\, \dd t.
		\end{aligned}
		\]
	\end{proof}

	\newpage
	\section{Transcendental Functions}
	\subsection{The Logarithm}
	\begin{mydef}
		If $x$ is a positive real number, then the \textbf{natural logarithm} of $x$ is defined by
		\[
		\ln x = \int_1^x \frac{1}{t}\, \dd t.
		\]
	\end{mydef}
	\begin{mythm}
		The logarithm function has the following properties:
		\begin{enumerate}
			\item [(1)] $\ln 1 = 0$;
			\item [(2)] $(\ln x)' = \dfrac{1}{x}$ for all $x>0$;
			\item [(3)] $\ln (xy) = \ln x + \ln y$ for all $x,y>0$;
			\item [(4)] $\ln x^r = r \ln x$ for all $x>0$ and all positive integers $r$;
			\item [(5)] $\ln \dfrac{1}{x} = - \ln x$ for all $x>0$;
			\item [(6)] $f(x)=\ln x$ is strictly increasing and unbounded on $(0,\infty)$;
			\item [(7)] $f(x)=\ln x$ is concave on $(0,\infty)$.
		\end{enumerate}
	\end{mythm}
	\begin{mythm}
		For every real number $b$, there exists a unique positive real number $a$ such that $\ln a = b$.
	\end{mythm}
	\begin{proof}
		Suppose $b>0$ and pick positive integer $n$ so that $\ln 2^n = n \ln 2 > b$. Since $f(x) = \ln x$ is continuous and strictly increasing on $[1, 2^n]$, by the \textbf{IVT}, there exists some unique $a\in (1, 2^n)$ such that $\ln a = b$.

		The case $b=0$ is trivial since $\ln 1 = 0$ and the case $b<0$ can be proved similarly.
	\end{proof}
	\begin{mydef}
		Particularly, we denote the unique positive real number $e$ such that
		\[
		\ln e = 1,
		\]
		which is called Euler's number or the natural number.
	\end{mydef}
	\begin{myrmk}
		Suppose $f(x)=c \ln x$ for some fixed constant $c$. We know it preserves all properties of $\ln x$. Now we know there exists a unique positive real number $\alpha$ such that
		\[
		c \ln \alpha = 1,\quad (\text{or }f(\alpha)=1,\alpha \neq 1) \implies c = \frac{1}{\ln \alpha}.
		\]
		With this $\alpha$ notation, we can rewrite
		\[
		f(x) = c \ln x = \frac{\ln x}{\ln \alpha},
		\]
		which motivates us to define logarithm with base $\alpha$.
	\end{myrmk}
	\begin{mydef}
		Suppose $b>0$ and $b\neq 1$. For any $x>0$, the \textbf{logarithm with base $b$} is defined by
		\[
		\log_b x = \frac{\ln x}{\ln b}.
		\]
		We usually call $b$ the \textbf{base} of the logarithm.
	\end{mydef}
	\begin{myrmk}
		Note that $\log_e x = \ln x$.
	\end{myrmk}
	\begin{myrmk}
		By injectivity of logarithm function, we know for any $b>0, u\in \mathbb{R}$ and $b\neq 1$, there exists a unique $x>0$ such that $\log_b x = u$. We denote this unique $x$ by
		\[
		b^u.
		\]
		If we adopt the convention that $1^u := 1$, then the power $b^u$ is well-defined for any $b>0$ and $u\in \mathbb{R}$.
	\end{myrmk}
	\begin{myex}
		Here we see how logarithm helps us to integrate some functions.
		\begin{enumerate}
			\item \[
			I = \int \tan x\, \dd x = \int \frac{\sin x}{\cos x}\, \dd x = -\int \frac{(\cos x)'}{\cos x} \, \dd x = -\ln |\cos x| + C.
			\]
			\item \[
			I = \int \ln x \, \dd x = \int 1 \cdot \ln x\, \dd x = x \ln x - \int x \cdot \frac{1}{x}\, \dd x = x \ln x - x + C.
			\]
			\item \[
			I_1 = \int \sin(\ln x)\, \dd x = \int \sin(\ln x) \cdot 1\, \dd x = x \sin(\ln x) - \int x \cdot \frac{1}{x} \cos(\ln x)\, \dd x.
			\]
			Now we need to calculate $\int \cos(\ln x)\, \dd x$:
			\[
			I_2 = \int \cos(\ln x)\, \dd x = \int \cos(\ln x) \cdot 1\, \dd x = x \cos(\ln x) + \int x \cdot \frac{1}{x} \sin(\ln x)\, \dd x.
			\]
			Thus we have
			\[
			\begin{aligned}
				\int \sin(\ln x)\, \dd x = \frac12 \left[ x \sin(\ln x) - x \cos(\ln x) \right] + C.\\
				\int \cos(\ln x)\, \dd x = \frac12 \left[ x \cos(\ln x) + x \sin(\ln x) \right] + C.
			\end{aligned}
			\]
		\end{enumerate}
	\end{myex}
	\begin{myrmk}
		Since
		\[
		\ln x = \int_1^x \frac{1}{t}\, \dd t\quad \text{for all}\ x>0,
		\]
		using substitution $u = 1-t$ we have
		\[
		\ln (1-x) = \int_1^{1-x} \frac{1}{t}\, \dd t = -\int_0^x \frac{1}{1-u}\, \dd u \quad \text{for all}\ x<1.
		\]
		We use a polynomial series to approximate $\dfrac{1}{1-u}$:
		\[
		-\ln (1-x) = \int_0^x \left( 1 + u + \frac{u^2}{1-u} \right) \dd u = x + \frac{x^2}{2} + \int_0^x \frac{u^2}{1-u}\, \dd u.
		\]
		Note when $x$ is near $0$, the polynomial $P(x)=x + \dfrac{x^2}{2}$ is a good approximation of $-\ln(1-x)$. This motivates us to introduce the following theorem.
	\end{myrmk}
	\begin{mythm}
		Let
		\[
		P_n(x) = x + \frac{x^2}{2} + \frac{x^3}{3} + \cdots + \frac{x^n}{n} = \sum_{k=1}^n \frac{x^k}{k}
		\]
		denote the polynomial of degree $n$. Then for any $x<1$ and any positive integer $n$, we have
		\[
		-\ln(1-x) = P_n(x) + \int_0^x \frac{t^n}{1-t}\, \dd t.
		\]
	\end{mythm}
	\begin{myrmk}
		We denote the last term by
		\[
		E_n(x) := \int_0^x \frac{t^n}{1-t}\, \dd t
		\]
		and call it the \textbf{error term}. Now we study the error term and see how good the approximation is.
	\end{myrmk}
	\begin{mythm}
		With the same setup as above, we have the following:
		\begin{enumerate}
			\item [(1)] If $0<x<1$, then
			\[
			\frac{x^{n+1}}{n+1} \leq E_n(x) \leq \frac{1}{(1-x)}\cdot\frac{x^{n+1}}{(n+1)}.
			\]
			\item [(2)] If $x<0$, then
			\[
			0 < (-1)^{n+1} E_n(x) \leq \frac{|x|^{n+1}}{n+1}.
			\]
		\end{enumerate}
	\end{mythm}
	\begin{myrmk}
		\begin{enumerate}
			\item [(1)] If $|x|$ is ``very small'', then $E_n(x)$ is also ``very small'' (no matter what $n$ is).
			\item [(2)] If $|x|$ is fixed and $n$ is ``very large'', then $E_n(x)$ is also ``very small'' and thus $P_n(x)$ is a good approximation of $-\ln(1-x)$.
		\end{enumerate}
	\end{myrmk}
	
	\subsection{Exponential function}
	\begin{mydef}
		For any real number $x$, the \textbf{exponential function} or the \textbf{antilogorithm} is defined by
		\[
		y = e^x = \text{the unique positive real number such that}\ \ln y = x.
		\]
	\end{mydef}
	\begin{mynot}
		We use the notation $\exp(x)$ to denote $e^x$ sometimes.
	\end{mynot}
	\begin{myrmk}
		Note that
		\[
		\operatorname{Dom} (\exp) = \mathbb{R}, \quad \operatorname{Im} (\exp) = (0, +\infty).
		\]
		Since $\ln x$ is continuous and strictly increasing on $(0, +\infty)$, we know $\exp(x)$ is also continuous and strictly increasing on $\mathbb{R}$.
	\end{myrmk}
	\begin{mythm}
		The exponential function has the following properties:
		\begin{enumerate}
			\item [(1)] $e^0 = 1$;
			\item [(2)] $(e^x)' = e^x$ for all $x\in \mathbb{R}$;
			\item [(3)] $e^{x+y} = e^x \cdot e^y$ for all $x,y\in \mathbb{R}$;
			% \item [(4)] $e^{-x} = \dfrac{1}{e^x}$ for all $x\in \mathbb{R}$;
			% \item [(5)] $f(x)=e^x$ is unbounded on $\mathbb{R}$;
			% \item [(6)] $f(x)=e^x$ is convex on $\mathbb{R}$.
		\end{enumerate}
	\end{mythm}
	\begin{proof}
		We only prove (2) here. By the definition of derivative, we have
		\[
		(e^x)' = \lim_{h\to 0} \frac{e^{x+h} - e^x}{h} = e^x \lim_{h\to 0} \frac{e^h - 1}{h}.
		\]
		It suffices to show the limit is $1$. Note that
		\[
		k := e^h - 1 \implies h = \ln (k+1).
		\]
		Thus we have
		\[
		\lim_{h\to 0} \frac{e^h - 1}{h} = \lim_{k\to 0} \frac{k}{\ln (k+1)} = \lim_{k\to 0} \frac{k}{\ln (1+k) - \ln 1} = \frac{1}{\left( \ln (1+k) \right)'}\Big|_{k=0} = 1.
		\]
	\end{proof}
	\begin{mydef}
		For any $a>0$ and $x\in \mathbb{R}$, we also define $a^x := e^{x \ln a}$.
	\end{mydef}
	\begin{mymthd}
		Using exponential function, we can rewrite integration by substitution as
		\[
		\int e^{f(x)} f'(x)\, \dd x = e^{f(x)} + C.
		\]
	\end{mymthd}
	\begin{mydef}
		The \textbf{hyperbolic sine} and \textbf{hyperbolic cosine} functions are defined by
		\[
		\sinh x = \frac{e^x - e^{-x}}{2}, \quad \cosh x = \frac{e^x + e^{-x}}{2}.
		\]
		Similarly, we define the \textbf{hyperbolic tangent} function by
		\[
		\tanh x = \frac{\sinh x}{\cosh x} = \frac{e^x - e^{-x}}{e^x + e^{-x}}.
		\]
	\end{mydef}
    \begin{mythm}
        Suppose $f$ is strictly increasing and continuous on $[a,b]$ and let $g$ be the inverse of $f$. If the derivative $f'(x)$ exists and is non-zero on $(a,b)$, then the derivative $g'(y)$ also exists and is non-zero at the corresponding $y=f(x)$. Moreover, we have
        \[
        g'(y) = \frac{1}{f'(g(y))} \quad \text{or}\quad \frac{\dd}{\dd y} f^{-1}(y) = \frac{1}{f'(f^{-1}(y))}.
        \]
    \end{mythm}
    \subsection{Inverse of Trigometric Functions}
    Note that the usual trigonometric functions are not invertible on their entire domains (due to periodicity).  So to define their ``inverses'', we consider them restricted to some interval on which they are monotonic.
    \begin{mydef}
        Here are the definitions of the three most common inverse trigonometric functions:
        \begin{enumerate}
            \item To define the \textbf{inverse sine function}, we consider $f(x)=\sin x$ for $x \in \left[-\dfrac{\pi}{2}, \dfrac{\pi}{2}\right]$, we denote the inverse function by $g(y)=\arcsin y$ for $y \in [-1,1]$.
            \item To define the \textbf{inverse cosine function}, we consider $f(x)=\cos x$ for $x \in [0, \pi]$, we denote the inverse function by $g(y)=\arccos y$ for $y \in [-1,1]$.
            \item To define the \textbf{inverse tangent function}, we consider $f(x)=\tan x$ for $x \in \left(-\dfrac{\pi}{2}, \dfrac{\pi}{2}\right)$, we denote the inverse function by $g(y)=\arctan y$ for $y \in \mathbb{R}$.
        \end{enumerate}
    \end{mydef}
    \begin{mythm}
        The derivatives of the inverse trigonometric functions are given by:
        \begin{enumerate}
            \item [(1)]
            \[
            g'(y) = \dfrac{1}{f'(g(y))} = \frac{1}{\cos(\arcsin y)} = \frac{1}{\sqrt{1-\sin^2(\arcsin y)}} = \frac{1}{\sqrt{1-y^2}}.
            \]
            \item [(2)]
            \[
            g'(y) = \dfrac{1}{f'(g(y))} = \frac{-1}{\sin(\arccos y)} = \frac{-1}{\sqrt{1-\cos^2(\arccos y)}} = \frac{-1}{\sqrt{1-y^2}}.
            \]
            \item [(3)]
            \[
            g'(y) = \dfrac{1}{f'(g(y))} = \frac{1}{\sec^2(\arctan y)} = \frac{1}{1+\tan^2(\arctan y)} = \frac{1}{1+y^2}.
            \]
        \end{enumerate}
    \end{mythm}
    \begin{mythm}
        The integrals involving inverse trigonometric functions are given by:
        \begin{enumerate}
            \item [(1)]
            \[
            \int \arcsin x\, \dd x = x \arcsin x + \int \frac{x^2}{\sqrt{1-x^2}}\, \dd x = x \arcsin x + \sqrt{1-x^2} + C.
            \]
            \item [(2)]
            \[
            \int \arccos x\, \dd x = x \arccos x - \int \frac{x^2}{\sqrt{1-x^2}}\, \dd x = x \arccos x - \sqrt{1-x^2} + C.
            \]
            \item [(3)]
            \[
            \int \arctan x\, \dd x = x \arctan x - \int \frac{x^2}{1+x^2}\, \dd x = x \arctan x - x + \ln(1+x^2) + C.
            \]
        \end{enumerate}
    \end{mythm}
    \begin{myrmk}
        As for $\arcsec x$, $\arccsc x$ and $\arccot x$, we define them to be:
        \begin{enumerate}
            \item $$ \arcsec x := \arccos \dfrac{1}{x},\quad |x| \geq 1; $$
            \item $$ \arccsc x := \arcsin \dfrac{1}{x},\quad |x| \geq 1; $$
            \item $$ \arccot x := \dfrac{\pi}{2}-\arctan x,\quad x\in\mathbb{R}$$
        \end{enumerate}
    \end{myrmk}

    \subsection{Integration by Partial Fractions}
    % \begin{myex}
    %     Try to calculate
    %     \[
    %     I = \int \frac{2x+5}{x^2+2x-3}\, \dd x.
    %     \]
    %     Note that
    %     \[
    %     \begin{aligned}
    %         I &= \int \frac{2x+5}{(x+3)(x-1)}\, \dd x\\
    %         &= \int \left( \frac{A}{x+3} + \frac{B}{x-1} \right) \dd x\quad (A,B\text{ are constants})\\
    %         &= \int \left( \frac{A(x-1) + B(x+3)}{(x+3)(x-1)} \right) \dd x\\
    %         &= \int \frac{(A+B)x + (3B - A)}{(x+3)(x-1)}\, \dd x.
    %     \end{aligned}
    %     \]
    %     By comparing coefficients, we have
    %     \[
    %     \begin{cases}
    %         A + B = 2,\\
    %         3B - A = 5.
    %     \end{cases}
    %     \implies
    %     \begin{cases}
    %         A = \dfrac74,\\[6pt]
    %         B = \dfrac14.
    %     \end{cases}
    %     \]
    %     Thus we have
    %     \[
    %     \begin{aligned}
    %         I &= \int \left( \frac{7/4}{x+3} + \frac{1/4}{x-1} \right) \dd x\\
    %         &= \frac74 \ln |x+3| + \frac14 \ln |x-1| + C.
    %     \end{aligned}
    %     \]
    % \end{myex}
    % \begin{myex}
    %     Try to calculate
    %     \[
    %     I = \int \frac{x^3+3x}{x^2-2x-3} \, \dd x.
    %     \]
    %     Note that since the degree of the numerator is greater than that of the denominator, we first use polynomial division:
    %     \[
    %     \begin{aligned}
    %         I &= \int \left( x + 2 + \frac{10x + 6}{x^2 - 2x - 3} \right) \dd x\\
    %         &= \int (x + 2)\, \dd x + \int \frac{10x + 6}{(x+1)(x-3)}\, \dd x.
    %     \end{aligned}
    %     \]
    %     Then we do similarly as before.
    % \end{myex}
    \begin{mydef}
        We call a rational function $r(x)=\dfrac{f(x)}{g(x)}$ \textbf{proper} if the degree of $f(x)$ is strictly less than that of $g(x)$. Otherwise, we call it \textbf{improper}.
    \end{mydef}
    \begin{myrmk}
        For an improper rational function, we can always use polynomial division to find polynomial $Q(x)$ and $R(x)$:
        \[
        \frac{f(x)}{g(x)} = Q(x) + \frac{R(x)}{g(x)}
        \]
        where $\dfrac{R(x)}{g(x)}$ is a proper rational function. Here $Q(x)$ is called the \textbf{quotient} and $R(x)$ is called the \textbf{remainder}.
    \end{myrmk}
    \begin{mydef}
        A quadratic function/factor of the form $x^2+bx+c$ is called \textbf{irreducible} over $\mathbb{R}$ if it cannot be further factored as $(x+x_1)(x+x_2)$ where $x_1,x_2\in \mathbb{R}$, i.e. its discriminant $b^2-4c<0$.
    \end{mydef}
    \begin{mythm}
        Every proper rational function can be expressed as a finite sum of a collection of terms each of which takes the form
        \[
        \frac{A}{(x+a)^k}, \quad \frac{Bx + C}{(x^2 + bx + c)^m}
        \]
        where $A,B,C,a,b,c$ are real numbers, $k,m$ are positive integers, and $x^2 + bx + c$ is irreducible over $\mathbb{R}$.
    \end{mythm}
    \begin{mymthd}
        Suppose $r(x)=\dfrac{f(x)}{g(x)}$ is proper.

        \textbf{Case 1:}
        $g(x)$ can be factorized into a collection of distinct linear factors. For example, suppose
        \[
        g(x)=(x-x_1)(x-x_2)\cdots (x-x_n) \quad \text{where } x_i \neq x_j \text{ for } i\neq j \in \left\{ 1, 2, \ldots, n \right\}
        \]
        Then
        \[
        \frac{f(x)}{g(x)}=\frac{A_1}{x-x_1} + \frac{A_2}{x-x_2} + \cdots + \frac{A_n}{x-x_n}\quad \text{for some } A_1, A_2, \ldots, A_n \in \mathbb{R}.
        \]
        and
        \[
        \int \frac{f(x)}{g(x)} \dd x = A_1 \ln |x-x_1| + A_2 \ln |x-x_2| + \cdots + A_n \ln |x-x_n| + C.
        \]
        \indent
        \textbf{Case 2:}
        $g(x)$ can still be factorized into a collection of linear terms, but some of these linear terms are repeated. For example, suppose $(x-a)$ is repeated $p$ times where $p\ge 2$ is an integer, then the corresponding term in the partial fraction decomposition becomes the following sum
        \[
        \frac{B_1}{(x-a)} + \frac{B_2}{(x-a)^2} + \cdots + \frac{B_p}{(x-a)^p}\quad \text{for some } B_1, B_2, \ldots, B_p \in \mathbb{R}.
        \]
        \indent
        \textbf{Case 3:}
        $g(x)$ contains non-repeated irreducible quadratic factors (in addition to linear factors) $x^2 + bx + c$ where $b^2 - 4c < 0$. Then in the partial fraction decomposition, we will have the term
        \[
        \frac{Ax+B}{x^2 + bx + c}\quad \text{for some } A,B \in \mathbb{R}.
        \]
        Now we integrate this term:
        \[
        \int \frac{Ax+B}{x^2 + bx + c}\, \dd x = \frac{A}{2} \ln |x^2 + bx + c| + \frac{2B - Ab}{\sqrt{4c - b^2}} \arctan \left( \frac{2x + b}{\sqrt{4c - b^2}} \right) + C.
        \]
        \indent
        \textbf{Case 4:}
        $g(x)$ contains repeated irreducible quadratic factors of the form $(x^2 + bx + c)^p$ where $p\ge 2$ is an integer. Then in the partial fraction decomposition, we will have the term
        \[
        \frac{A_1 x + B_1}{x^2 + bx + c} + \frac{A_2 x + B_2}{(x^2 + bx + c)^2} + \cdots + \frac{A_p x + B_p}{(x^2 + bx + c)^p}\quad \text{for some } A_i, B_i \in \mathbb{R}.
        \]
        And we do similarly as Case 3 to integrate each term.
    \end{mymthd}

    \newpage
    \section{Polynomial Approximations to functions}
    With the help of a special identity $1-x^n = (1-x)(1+x+x^2+\cdots + x^{n-1})$, we can find a polynomial approximation to $-\ln(1-x)$ in the previous section. However, this may not help to find a polynomial approximation to other functions, which motivates us to find a more general method to find polynomial approximations to functions.
    \subsection{The Taylor Polynomial for a function}
    \begin{mythm}
        Let $f$ be a function with derivatives of order $n$ at the point $x=a$. Then there exists a unique polynomial $P$ of degree at most $n$ such that $P(a)=f(a), P'(a)=f'(a), \ldots, P^{(n)}(a) = f^{(n)}(a)$. Additionally, this polynomial is given by
        \[
        P(x) = \sum_{k=0}^n \frac{f^{(k)}(a)}{k!} (x-a)^k = f(a) + f'(a)(x-a) + \cdots + \frac{f^{(n)}(a)}{n!} (x-a)^n.
        \]
        and it is referred to as the \textbf{Taylor polynomial of degree $n$ for $f$ at $x=a$}.
    \end{mythm}
    \begin{mynot}
        We often denote $P$ in the above theorem by $T_n(f)$ or $T_n(f;a)$ to emphasize the degree and the center of the Taylor polynomial.
    \end{mynot}
    \begin{myex}
        Here are some examples of Taylor polynomials:
        \begin{enumerate}[label=(\arabic*)]
            \item The Taylor polynomial of degree $n$ for $f(x)=e^x$ at $x=0$ is given by
            \[
            T_n(f) = 1 + x + \frac{x^2}{2!} + \cdots + \frac{x^n}{n!} = \sum_{k=0}^n \frac{x^k}{k!}.
            \]
            \item The Taylor polynomial of degree $2n+1$ for $f(x)=\sin x$ at $x=0$ is given by
            \[
            T_{2n+1}(f) = x - \frac{x^3}{3!} + \cdots + (-1)^n \frac{x^{2n+1}}{(2n+1)!} = \sum_{k=0}^n (-1)^k \frac{x^{2k+1}}{(2k+1)!}.
            \]
            \item The Taylor polynomial of degree $2n$ for $f(x)=\cos x$ at $x=0$ is given by
            \[
            T_{2n}(f) = 1 - \frac{x^2}{2!} + \cdots + (-1)^n \frac{x^{2n}}{(2n)!} = \sum_{k=0}^n (-1)^k \frac{x^{2k}}{(2k)!}.
            \]
        \end{enumerate}
    \end{myex}
    \begin{mythm}
        Here are some properties of the Taylor polynomial:
        \begin{enumerate}[label=(\arabic*)]
            \item Linearity: If $c_1,c_2 \in \mathbb R$, then $T_n(c_1 f + c_2 g) = c_1 T_n(f) + c_2 T_n(g)$.
            \item Differentiation: $(T_n(f))' = T_{n-1}(f')$.
            \item Integration: If $\displaystyle F(x)=\int_a^x f(t)\, \dd t$, then \[T_{n+1}(F) = \int_a^x T_n(f(t))\, \dd t.\]
            \item Substitution: If $g(x)=f(cx)$ where $c\in \mathbb R$ is a constant, then \[T_n(g)(x;a) = T_n(f)(cx; ca).\]
            Particularly, if $a=0$, then $T_n(g)(x) = T_n(f)(cx)$.
        \end{enumerate}
    \end{mythm}
    \begin{myex}
        With the help of the properties above, we continue to find the Taylor polynomial of degree $n$ for more functions at $x=0$:
        \begin{enumerate}[label=(\arabic*),start=4]
            \item The Taylor polynomial of degree $2n$ for $f(x)=\cosh x$ at $x=0$ is given by
            \[
            T_{2n}(f) = 1 + \frac{x^2}{2!} + \cdots + \frac{x^{2n}}{(2n)!} = \sum_{k=0}^n \frac{x^{2k}}{(2k)!}.
            \]
            \item The Taylor polynomial of degree $2n-1$ for $f(x)=\sinh x$ at $x=0$ is given by
            \[
            T_{2n-1}(f) = x + \frac{x^3}{3!} + \cdots + \frac{x^{2n-1}}{(2n-1)!} = \sum_{k=0}^{n-1} \frac{x^{2k+1}}{(2k+1)!}.
            \]
            \item Recall
            \[
            \frac1{1-x} = 1 + x + x^2 + \cdots + x^n + \frac{x^{n+1}}{1-x}:=P_n(x)+\frac{x^{n+1}}{1-x} \quad \text{for all } x\neq 1.
            \]
            Let $h(x) = x^n g(x)$ where $g(x) = \dfrac{x}{1-x}$, then we have
            \[
            f(x)=h(x) + P_n(x) \quad \text{for all } x\neq 1.
            \]
            Note that by checking the derivatives of $f$ and $P_n$ at $x=0$, we have
            \[
            f^{(k)}(0) = P_n^{(k)}(0) \quad \text{for all } k\in \{0,1,\ldots,n\}.
            \]
            Thus we have $T_n(x) = P_n(x)$. By $\displaystyle \int_0^x \frac{1}{1-t} \dd t = -\ln(1-x)$, we have
            \[
            T_{n+1}\left[ -\ln(1-x) \right] = \int_0^x T_n\left( \frac{1}{1-t} \right) \dd t = \int_0^x P_n(t) \dd t = \sum_{k=1}^{n+1} \frac{x^k}{k}.
            \]
        \end{enumerate}
    \end{myex}
    \begin{mythm}
        Let $P_n$ be a polynomial of degree $n\ge 1$. Let $f$ and $g$ be two functions with derivatives of order $n$ at $x=0$ and assume that $f(x)=P_n(x)+x^n g(x)$ where $g(x)\to 0$ as $x\to 0$. Then $P_n$ is the Taylor polynomial $T_n(f)$ at 0.
    \end{mythm}
    \begin{myex}
        \begin{enumerate}[label=(\arabic*),start=7]
        \item We start again by
        \[
        \frac{1}{1-x} = 1 + x + x^2 + \cdots + x^n + \frac{x^{n+1}}{1-x}
        \]
        but replace $x$ by $(-x^2)$:
        \[
        \frac{1}{1+x^2} = 1 - x^2 + x^4 - \cdots + (-1)^n x^{2n} + \frac{(-1)^{n+1} x^{2n+1}}{1+x^2}\quad \forall x\in \R.
        \]
        Analogously, we let $P_{2n}(x) = 1 - x^2 + x^4 - \cdots + (-1)^n x^{2n}$ and $g(x) = \dfrac{(-1)^{n+1} x}{1+x^2}$, then we have
        \[
        \frac{1}{1+x^2} = P_{2n}(x) + x^{2n} g(x) \quad \text{for all } x\in \R.
        \]
        By the above theorem, we have $T_{2n}\left( \dfrac{1}{1+x^2} \right) = P_{2n}(x)$. By $\displaystyle \int_0^x \frac{1}{1+t^2} \dd t = \arctan x$, we have
        \[
        T_{2n+1}(\arctan x) = \int_0^x T_{2n}\left( \frac{1}{1+t^2} \right) \dd t = \int_0^x P_{2n}(t) \dd t = \sum_{k=0}^n \frac{(-1)^k x^{2k+1}}{2k+1}.
        \]
        \end{enumerate}
    \end{myex}
    \begin{mydef}
        We call the function $E_n(x) := f(x) - T_n(f(x))$ the \textbf{error term} or the \textbf{remainder} of the Taylor polynomial $T_n(f)$.
    \end{mydef}
    \begin{myex}
        \begin{enumerate}[label=(\arabic*),start=8]
            \item Suppose $f$ has continuous second derivative $f''$ in some neighborhood of $a$. Then in this neighborhood of $a$, we know:
            \[
            \begin{aligned}
                E_1(x) &= f(x) - T_1(f) \\
                &= f(x) - \left[ f(a) + f'(a)(x-a) \right] \\
                &= \int_a^x (f'(t) - f'(a))\, \dd t\quad (\text{IBP with } u=f'(t)-f'(a),v=t-x ) \\
                &= \left[ f'(t) - f'(a) \right](t-x)\Big|_{t=a}^x - \int_a^x f''(t)(t-x)\, \dd t\\
                &= \int_a^x f''(t)(x-t)\, \dd t.
            \end{aligned}
            \]
            If, additionally, we know $m\le f''(t) \le M$ for any $t$ in the considered neighborhood of $a$, then we have
            \[
            \frac{m(x-a)^2}{2} \le E_1(x) = \int_a^x f''(t)(x-t)\, \dd t \le \frac{M(x-a)^2}{2}.
            \]
            By the squeeze theorem, we have $\displaystyle \lim_{x\to a} E_1(x) = 0$.
        \end{enumerate}
    \end{myex}
    \begin{mythm}\label{Taylor formula with error term}
        Assume $f$ has a continuous derivative of order $n+1$ in some interval containing $a$. Then for every $x$ in this interval, we have the Taylor formula:
        \[
        f(x) = T_n(f) + E_n(x) = \sum_{k=0}^n \frac{f^{(k)}(a)}{k!} (x-a)^k + \frac{1}{n!} \int_a^x f^{(n+1)}(t) (x-t)^n\, \dd t.
        \]
        If, additionally, there exists $m$ and $M$ such that $m \le f^{(n+1)}(t) \le M$ for any $t$ in the considered interval, then we have the following bounds/estimates for the error term:
        \[
        \frac{m(x-a)^{n+1}}{(n+1)!} \le E_n(x) = \frac{1}{n!} \int_a^x f^{(n+1)}(t) (x-t)^n\, \dd t \le \frac{M(x-a)^{n+1}}{(n+1)!} \ (\forall x>a).
        \]
        and
        \[
        \frac{m(a-x)^{n+1}}{(n+1)!} \le (-1)^{n+1} E_n(x)  \le \frac{M(a-x)^{n+1}}{(n+1)!} \quad \forall x<a.
        \]
        \end{mythm}
        \begin{myex}
            We go back to the example of $\displaystyle f(x) = e^x=\sum_{k=0}^n \frac{x^k}{k!} + E_n(x)$ and $a=0$. Note for any $x\in [b,c]$, we know $e^b \le e^x \le e^c$. Particularly, if $b=0$ and $c=x=1$, then we have
            \[
            \frac{1}{(n+1)!} \le E_n(1) \le e\cdot \frac{1}{(n+1)!} < \frac{3}{(n+1)!}.
            \]
            Thus we have $\displaystyle e=\sum_{k=0}^n \frac{1}{k!} + E_n(1)$.
        \end{myex}
        \begin{myrmk}
            Recall the error term
            \[
            E_n(x) = \frac{1}{n!} \int_a^x f^{(n+1)}(t) (x-t)^n\, \dd t.
            \]
            Note that for any $t\in [a,x]$, the term $(x-t)^n$ never changes sign. Back in Theorem~\ref{Taylor formula with error term}, we require $f^{(n+1)}$ to be continuous on the considered interval. By the weighted Mean Value Theorem for integrals, we can find some $c\in [a,x]$ such that
            \[
            \begin{aligned}
                E_n(x) &= \frac{1}{n!} \int_a^x f^{(n+1)}(t) (x-t)^n\, \dd t\\
                &= \frac{1}{(n+1)!} f^{(n+1)}(c) \int_a^x (x-t)^n\, \dd t\\
                &= \frac{1}{(n+1)!} f^{(n+1)}(c) (x-a)^{n+1}\quad \text{for some } c\in [a,x].
            \end{aligned}
            \]
            This is called the \textbf{Lagrange form of the error term}.

            Suppose $f^{(n+1)}$ is continous on some interval $[a-r,a+r]$ where $r>0$, then $f^{(n+1)}$ is bounded on $[a-r,a+r]$ and there exists $M>0$ such that $|f^{(n+1)}(t)| \le M$ for any $t\in [a-r,a+r]$. Then
            \[
            |E_n(x)| = M \cdot \frac{|x-a|^{n+1}}{(n+1)!} \quad \text{for any } x\in [a-r,a+r].
            \]
            For $x\neq a$, we know $\displaystyle 0 \le \big|\frac{E_n(x)}{(x-a)^{n}}\big| \le \frac{M}{(n+1)!}|x-a|$. By the squeeze theorem, we know $\displaystyle \frac{E_n(x)}{(x-a)^{n}} \to 0$ as $x \to a$. We say that the error term $E_n$ is of small order that $(x-a)^n$ as $x\to a$, and we often write $E_n(x) = o\left( (x-a)^n \right)$ as $x\to a$.
        \end{myrmk}
        \begin{mydef}
            Assume $g(x)\neq 0$ for $x\neq a$ in some interval containing $a$. The notation $f(x) = o(g(x))$ as $x\to a$ means that $\displaystyle \lim_{x\to a} \frac{f(x)}{g(x)} = 0$.
        \end{mydef}
        \begin{myex}
            We can rewrite Taylor's formula in o-notation as
            \[
            f(x) = T_n(f) + o\left( (x-a)^n \right) \quad \text{as } x\to a.
            \]
        \end{myex}
        \begin{mythm}
            We deduce some useful properties of the o-notation: as $x\to a$,
            \begin{enumerate}[label=(\arabic*)]
                \item $o(g(x)) \pm o(g(x)) = o(g(x)).$
                \item $o(cg(x)) = o(g(x))$ for any real constant $c\neq 0$.
                \item $f(x) \cdot o(g(x)) = o(f(x)g(x))$ for $f(x)\neq 0$.
                \item $o(o(g(x))) = o(g(x))$.
                \item $\displaystyle \frac{1}{1+g(x)} = 1 - g(x) + o(g(x))$ if $g(x)\to 0$ as $x\to a$.
            \end{enumerate}
        \end{mythm}
        \begin{proof}
            We only prove (1) and (5) here since the others are analogous.
            \begin{enumerate}[label=(\arabic*)]
                \item Let $f_1(x) = o(g(x))$ and $f_2(x) = o(g(x))$, then $\displaystyle \lim_{x\to a} \frac{f_1(x)}{g(x)} = 0$ and $\displaystyle \lim_{x\to a} \frac{f_2(x)}{g(x)} = 0$. Thus we have
                \[
                \lim_{x\to a} \frac{f_1(x) \pm f_2(x)}{g(x)} = \lim_{x\to a} \frac{f_1(x)}{g(x)} \pm \lim_{x\to a} \frac{f_2(x)}{g(x)} = 0.
                \]
                Thus $o(g(x)) + o(g(x)) = f_1(x) + f_2(x) = o(g(x))$.
                \item[(5)] Recall
                \[
                \frac{1}{1+g(x)} = 1 - g(x) + \frac{g(x)^2}{1+g(x)}\quad \text{for any } g(x)\neq -1.
                \]
                Thus we have
                \[
                \lim_{x\to a} \frac{\frac{1}{1+g(x)} - (1 - g(x))}{g(x)} = \lim_{x\to a} \frac{g(x)^2}{(1+g(x))g(x)} = \lim_{x\to a} \frac{g(x)}{1+g(x)} = 0.
                \]
                Thus we have $\displaystyle \frac{1}{1+g(x)} = 1 - g(x) + o(g(x))$.
            \end{enumerate}
        \end{proof}
        \begin{myex}
            We know 
            \[
            \cos x = 1 - \frac{x^2}{2!} + \cdots + \frac{(-1)^n x^{2n}}{(2n)!} + o(x^{2n+1}) \quad \text{as } x\to 0.
            \]
                Let $g(x) = \cos x -1$, then $g(x)\to 0$ as $x\to 0$. By the properties of o-notation, we have
                \[
                \frac{1}{\cos x} = \frac{1}{1 + g(x)} = 1 - g(x) + o(g(x)) = 1 +\frac{x^2}{2} + o(x^2) + o(x^3) \quad \text{as } x\to 0.
                \]
            Note $o(x^2)+o(x^3)=o(x^2)$, so we have
            \[
            \tan x = \frac{\sin x}{\cos x} = \left( x - \frac{x^3}{3!} + o(x^4) \right) \left( 1 + \frac{x^2}{2} + o(x^3) \right) = x + \frac{x^3}{3} + o(x^3) \quad \text{as } x\to 0.
            \]
        \end{myex}
        \begin{myex}
            With Taylor's formula, we can calculate the following limit:
            \begin{enumerate}
                \item \(\displaystyle
            \lim_{x\to 0} \frac{a^x - b^x}{x} \quad \text{where } a,b>0.
            \)

            Note that
            \[
            e^x = 1 + x + o(x) \quad \text{as } x\to 0.
            \]
            Thus we have
            \[
            a^x - b^x = e^{x\ln a} - e^{x\ln b} = x(\ln a - \ln b) + o(x) \quad \text{as } x\to 0.
            \]
            By definition of o-notation, we have
            \[
            \lim_{x\to 0} \frac{a^x - b^x}{x} = \ln a - \ln b.
            \]
                \item \(\displaystyle
            \lim_{x\to 0} \frac{\ln(1+ax)}{x} \quad \text{where } a\in \mathbb{R}.
            \)

            If $a\neq 0$, note that
            \[
            \ln(1+ax) = ax + o(x) \quad \text{as } x\to 0.
            \]
            Thus we have
            \[
            \frac{\ln(1+ax)}{x} = a + o(1) \quad \text{as } x\to 0 \implies \lim_{x\to 0} \frac{\ln(1+ax)}{x} = a.
            \]
            But, what if $a=0$? 
            \end{enumerate}
        \end{myex}

        \subsection{Indeterminate forms and L'Hôpital's Rule}
        \begin{mythm}
            Assume $f$ and $g$ have derivatives $f'(x)$ and $g'(x)$ for every $x$ in an open interval $(a,b)$ and suppose that $\displaystyle \lim_{x\to a^+} f(x) = \lim_{x\to a^+} g(x) = 0$. If $g'(x) \neq 0$ for every $x\in (a,b)$ and the limit $\displaystyle \lim_{x\to a^+} \frac{f'(x)}{g'(x)}$ exists and have value $L$, then we have $\displaystyle \lim_{x\to a^+} \frac{f(x)}{g(x)} = \frac{f'(x)}{g'(x)} = L$.
        \end{mythm}
        \begin{proof}
            Note that L'ôpital's Rule does require $f,g,f',g'$ to be well-defined for $x\in (a,b)$. So we shall ``extend'' the definition of $f$ and $g$ ``a bit'' by letting
            \[
            F(x) = \begin{cases}
            f(x) & \text{if } x\in (a,b)\\
            0 & \text{if } x=a
            \end{cases}
            \quad \text{and} \quad
            G(x) = \begin{cases}
            g(x) & \text{if } x\in (a,b)\\
            0 & \text{if } x=a
            \end{cases}
            \]
            We can apply Cauchy's Mean Value Theorem to the interval $[a,x]$ for $F$ and $G$ to get
            \[
            \frac{F(x) - F(a)}{G(x) - G(a)} = \frac{F'(c)}{G'(c)} \quad \text{for some } c\in (a,x).
            \]
            Simply note that $F(a) = G(a) = 0$, we have
            \[
            \frac{f(x)}{g(x)} = \frac{F(x)}{G(x)} = \frac{F'(c)}{G'(c)} = \frac{f'(c)}{g'(c)} \quad \text{for some } c\in (a,x).
            \]
            Note $g(x) \neq 0$ for $x\in (a,b)$ otherwise $G(x) = 0 = G(a)$ and by Role's Theorem, there exists some $d\in (a,x)$ such that $G'(d) = 0 = g'(d)$, which is a contradiction. Thus as $x\to a^+$, we have $c\to a^+$ and $\displaystyle \lim_{x\to a^+} \frac{f(x)}{g(x)} = \lim_{c\to a^+} \frac{f'(c)}{g'(c)} = L$.
        \end{proof}
        \begin{myrmk}
            The "left-sided limit" and the "two-sided limit" versions of L'Hôpital's Rule also hold provided that the derivatives are defined in the needed open interval.
            
            More versions of L'Hôpital's Rule can be derived for indeterminate forms of the type $\infty/\infty$, $0\cdot \infty$, $\infty - \infty$, $0^0$, $1^\infty$ and $\infty^0$ or for ``limit notation'' like $\displaystyle \lim_{x\to \infty } \frac{f(x)}{g(x)}$, $\displaystyle \lim_{x\to a} \frac{f(x)}{g(x)} = \infty$ or $\displaystyle \lim_{x\to \infty} \frac{f(x)}{g(x)} = -\infty$.
        \end{myrmk}
        \begin{myex}
            We use L'Hôpital's Rule to calculate some previous limits:
            \begin{enumerate}[label=(\arabic*)]
                \item \(\displaystyle
            \lim_{x\to 0} \frac{a^x - b^x}{x} \quad \text{where } a,b>0.
            \)

            Note that
            \[
            \lim_{x\to 0} a^x = \lim_{x\to 0} b^x = 1.
            \]
            Thus we have
            \[
            \lim_{x\to 0} \frac{a^x - b^x}{x} = \lim_{x\to 0} \frac{a^x \ln a - b^x \ln b}{1} = \ln a - \ln b.
            \]
                \item \(\displaystyle
            \lim_{x\to 0} \frac{\ln(1+ax)}{x} \quad \text{where } a\in \mathbb{R}.
            \)

            We talk about the case $a=0$ here. Note that
            \[
            \lim_{x\to 0} \ln(1+ax) = \ln 1 = 0.
            \]
            Thus we have
            \[
            \lim_{x\to 0} \frac{\ln(1+ax)}{x} = \lim_{x\to 0} \frac{a}{1+ax} = a.
            \]
            \item $\displaystyle \lim_{x\to 0} \frac{x-\tan x}{x-\sin x} = \lim_{x\to 0}\frac{1-\sec^2 x}{1-\cos x} =-\lim_{x\to 0} \frac{\cos x+1}{\cos^2 x} -2$.
            \end{enumerate}
        \end{myex}
        \begin{mydef}
            We introduce extended definitions of limits and ``limits with infinity'' here:
            \begin{itemize}
                \item The symbol
                \[
                \lim_{x\to \infty} f(x) = L \quad \text{or} \quad f(x) \to L \quad \text{as } x\to \infty
                \]
                means that for every $\varepsilon > 0$, there exists some $M>0$ such that \[|f(x) - L| < \varepsilon \quad \text{for any } x > M. \]
                \item The symbol $\displaystyle \lim_{x\to a} f(x) = \infty$ means that for every $M>0$, there exists some $\delta > 0$ such that $f(x) > M$ whenever $0 < |x-a| < \delta$. Similarly, the symbol $\displaystyle \lim_{x\to a} f(x) = -\infty, \displaystyle \lim_{x\to a^+} f(x) = \pm \infty$ and $\displaystyle \lim_{x\to a^-} f(x) = \pm \infty$ can be defined analogously.
                \item The symbol $\displaystyle \lim_{x\to \infty} f(x) = \infty$ means that for every $M>0$, there exists some $N>0$ such that $f(x) > M$ whenever $x > N$. Similarly, the symbol $\displaystyle \lim_{x\to \infty} f(x) = -\infty$ can be defined analogously.
            \end{itemize}
        \end{mydef}
        \begin{myrmk}
            Note, if $g(t) = f(1/t)$ and $\displaystyle \lim_{t\to 0^+} g(t) = A$, then $\displaystyle \lim_{x\to \infty} f(x) = A$. Similarly, if $\displaystyle \lim_{t\to 0^-} g(t) = B$, then $\displaystyle \lim_{x\to -\infty} f(x) = B$.
        \end{myrmk}
        \begin{myrmk}
            With the extended definitions of limits, we can also apply L'Hôpital's Rule to calculate limits like $\displaystyle \lim_{x\to \infty} \frac{f(x)}{g(x)}$ provided that the derivatives are defined in the needed open interval and the limit $\displaystyle \lim_{x\to \infty} \frac{f'(x)}{g'(x)}$ exists.
        \end{myrmk}
        \begin{mythm}
            Suppose $a,b>0$ then $\displaystyle \lim_{x\to +\infty} \frac{(\ln x)^b}{x^a} = 0$ and $\displaystyle \lim_{x\to +\infty} \frac{x^b}{e^{ax}} = 0$.
        \end{mythm}
        \begin{proof}
            Suppose $c>0$ and $t\ge 1$. Then $\dfrac{1}{t} \le \dfrac{1}{t^{1-c}}$. For any $x>1$, we have
            \[
            0<\ln x = \int_1^x \frac{1}{t} \dd t \le \int_1^x \frac{1}{t^{1-c}} \dd t = \frac{x^c - 1}{c} < \frac{x^c}{c}.
            \]
            Thus we have
            \[
            0 < \frac{(\ln x)^b}{x^a} < \frac{x^{cb-a}}{c^b} \quad \text{for any } x>1.
            \]
            Since $c$ is arbitrary, we choose $c=\dfrac{a}{2b}$ to get $x^{cb-a} = x^{-\frac{a}{2}} \to 0$ as $x\to +\infty$. By the squeeze theorem, we have $\displaystyle \lim_{x\to +\infty} \frac{(\ln x)^b}{x^a} = 0$.

            For the second limit, we let $t=e^x$ and we have
            \[
            \lim_{x\to +\infty} \frac{x^b}{e^{ax}} = \lim_{t\to +\infty} \frac{(\ln t)^b}{t^a} = 0.
            \]
        \end{proof}
    \begin{myex}
        \[
        \lim_{x\to 0} x^x = \lim_{x\to 0} e^{x\ln x} = e^{\lim_{x\to 0} x\ln x} = e^0 = 1.
        \]
    \end{myex}
    \begin{myex}
        \[
        \lim_{x\to \infty} x^{1/x} = \lim_{x\to 0} \frac{1}{x^x} = \frac{1}{\displaystyle \lim_{x\to 0} x^x} = \frac{1}{1} = 1.
        \]
    \end{myex}
    \newpage
    \section{Sequences and Series}
    \subsection{Sequences and Partial Sums}
    \begin{mydef}
        A function $f:\N \to \R$ is called a \textbf{sequence}. The function value $f(n)$ is called the $n$-th term of the sequence. We often denote such a sequence by $\{a_n\}$ or $f(n)$ where $a_n := f(n)$ for any $n\in \N$.
    \end{mydef}
    \begin{mydef}
        A sequence $\{f(n)\}$ is said to have a limit $L\in \R$, or converges to $L$, if for every $\varepsilon > 0$, there exists some $N\in \N$ such that $|f(n) - L| < \varepsilon$ for any $n > N$ and in this case we say the sequence is \textbf{convergent}. If a sequence is not convergent, then we say it is \textbf{divergent}.
    \end{mydef}
    \begin{mydef}
        A sequence $\{f(n)\}$ is \textbf{increasing} if $f(n) \le f(n+1)$ for all $n\in \N$. A sequence $\{f(n)\}$ is \textbf{decreasing} if $f(n) \ge f(n+1)$ for all $n\in \N$. A sequence is \textbf{monotonic} if it is either increasing or decreasing. A sequence $\{f(n)\}$ is \textbf{bounded} if there exists some $M>0$ such that $|f(n)| \le M$ for any $n\in \N$. A sequence is \textbf{unbounded} if it is not bounded.
    \end{mydef}
    \begin{mythm}
        A monotonic sequence converges if and only if it is bounded.
    \end{mythm}
    \begin{proof}
        If $\{f(n)\}$ is increasing, then it converges. Additionally, if $S = \{f(n):n\in \N\}$, then $f(n) \to \sup S = L$. For any $\varepsilon > 0$, $L-\varepsilon$ is not an upper bound of $S$, so there exists some $N\in \N$ such that $f(N) > L-\varepsilon$. For any $n\ge N$, we have $f(n) \ge f(N) > L-\varepsilon$ and $f(n) \le L$, so $|f(n) - L| < \varepsilon$. Thus $\{f(n)\}$ converges to $L$. The case for decreasing sequence is analogous.
    \end{proof}
    \begin{mydef}
        For any sequence $\{a_n\}$, we can generate another sequence $\{S_n\}$ by defining
        \[
        S_n = a_1 + a_2 + \cdots + a_n = \sum_{k=1}^n a_k \quad \text{for any } n\in \N.
        \]
        We call $S_n$ the partial sums of the first $n$ terms of $\{a_n\}$ and $\{S_n\}$ is called an \textbf{(infinite) series}.
    \end{mydef}
    \begin{mynot}
        The following notations are often used for series:
        \[
        \sum a_n = a_1 + a_2 + \cdots = \sum_{n=1}^\infty a_n.
        \]
        If $S_n \to S$ as $n\to \infty$, then we say the series $\{S_n\}$ converges and has the sum $S$. In this case, we often write $\displaystyle \sum_{n=1}^\infty a_n = S$ instead of $S_n \to S$. A series that does not converge is divergent and it has no sum.
    \end{mynot}
    \begin{myex}
        We list some typical examples of series here:
        \begin{enumerate}
            \item If $\displaystyle a_n = \frac{1}{n}$, then
            \[
            S_n = \sum_{k=1}^n \frac{1}{k} > \int_1^{n+1} \frac{1}{t} \dd t = \ln(n+1) \to +\infty \quad \text{as } n\to +\infty.
            \]
            This divergent series is called the \textbf{harmonic series}.
            \item If $\displaystyle a_n = \frac{1}{2^{n-1}}$, then
            \[
            S_n = \sum_{k=1}^n \frac{1}{2^{k-1}} = 2\left( 1 - \frac{1}{2^n} \right) \to 2 \quad \text{as } n\to +\infty.
            \]
            This kind of series that the ratio of two adjacent terms is a constant is called a \textbf{geometric series}.
        \end{enumerate}
    \end{myex}
    \begin{mythm}
        Let $\sum a_n$ and $\sum b_n$ be two convergent series and $\alpha,\beta \in \R$. Then the series $\sum (\alpha a_n + \beta b_n)$ converges and its sum is given by
        \[
        \sum_{n=1}^\infty (\alpha a_n + \beta b_n) = \alpha \sum_{n=1}^\infty a_n + \beta \sum_{n=1}^\infty b_n.
        \]
    \end{mythm}
    \begin{mythm}
        If $\{b_n\}$ is a sequence and another sequence $\{a_n\}$ satisfies $a_n = b_n - b_{n+1}$ for any $n\in \N$, then we call the series $\sum a_n$ a \textbf{telescoping series} and $\sum a_n$ converges if and only if $\{b_n\}$ converges. Furthermore, if $\{b_n\}$ converges to $L$ as $n\to \infty$, then we have
        \[
        \sum_{n=1}^\infty a_n = \sum_{n=1}^\infty (b_n - b_{n+1}) =  b_1 - L.
        \]
    \end{mythm}
    \begin{mythm}
        Let $a_n = x^n$ for $n\in \N$, then the partial sum $\displaystyle S_n = \sum_{k=1}^n x^k$ satisfies
        \[
        S_n = \begin{cases}
            n+1 & \text{if } x=1\\
            \dfrac{1-x^{n}}{1-x} & \text{if } x\neq 1
        \end{cases}
        \]
        We call $\{S_n\}$ the \textbf{geometric series} and $\{S_n\}$ converges if and only if $|x|<1$. When $\{S_n\}$ converges, we have $\displaystyle \sum_{n=1}^\infty x^n = \frac{1}{1-x}$.
    \end{mythm}
    \begin{myex}
        We can do some replacement to get some typical examples of geometric series:
        \begin{enumerate}[label=(\arabic*)]
            \item Replace $x$ by $x^2$ in the geometric series, we have
            \[
            1+ x^2 + x^4 + \cdots = \frac{1}{1-x^2} \quad \text{for } |x|<1.
            \]
            \item For $|x|<1$, we multiply both sides by $x$:
            \[
            x + x^3 + x^5 + \cdots = \frac{x}{1-x^2}.
            \]
            \item Replace $x$ by $-x$:
            \[
            1 - x + x^2 - x^3 + \cdots = \frac{1}{1+x} \quad \text{for } |x|<1.
            \]
        \end{enumerate}
    \end{myex}
    \begin{mydef}\label{def:power-series}
        We call a series of the form
        \[
        \sum_{n=0}^\infty a_n x^n
        \]
        a \textbf{power series} in $x$ and the number $a_n$ is called the coefficient of $x^n$ in the power series.
    \end{mydef}
    \begin{myex}
        Suppose $f$ is a function with derivatives of arbitrary order in a neighborhood of $x=0$. Let $a_n = \dfrac{f^{(n)}(0)}{n!}$ for $n\in \N$, then the power series with coefficients $a_n$ is called the \textbf{Taylor series} of $f$ at $x=0$. Since
        \[
        f(x) = \sum_{n=0}^\infty \frac{f^{(n)}(0)}{n!} x^n + E_n(x),
        \]
        Thus if $E_n(x) \to 0$ as $n\to \infty$ for any $x$ in some neighborhood of $0$, then the Taylor series converges to $f(x)$.
    \end{myex}

    \subsection{Convergence Tests for Series}
    In this subsection we introduce three tests for convergence of series.
    \begin{mythm}
        If $\sum a_n$ converges, then $\displaystyle \lim_{n\to \infty} a_n = 0$.
    \end{mythm}
    \begin{proof}
        Suppose $\displaystyle S_n = \sum_{k=1}^n a_k$. Then $a_n = S_n - S_{n-1}$. Since $\{S_n\}$ converges, we have $\displaystyle \lim_{n\to \infty} a_n = \displaystyle \lim_{n\to \infty} S_n - \displaystyle \lim_{n\to \infty} S_{n-1} = 0$.
    \end{proof}
    \begin{myrmk}
        This is a necessary condition for convergence of series but not a sufficient condition. For example, the harmonic series $\displaystyle \sum_{n=1}^\infty \frac{1}{n}$ diverges but $\displaystyle \lim_{n\to \infty} \frac{1}{n} = 0$.
    \end{myrmk}
    \begin{mythm}
        Suppose $a_n \ge 0$ for any $n\in \N$. Then the series $\sum a_n$ converges if and only if the sequence of partial sums $\{S_n\}$ is bounded.
    \end{mythm}
    \begin{proof}
        We know $\{S_n\} $ is increasing and by the monotone convergence theorem, $\{S_n\}$ converges if and only if $\{S_n\}$ is bounded.
    \end{proof}
    \begin{myrmk}
        This is a necessary and sufficient condition for convergence test.
    \end{myrmk}
    \begin{mythm}
        Suppose $a_n\ge 0$ and $b_n\ge 0$ for any $n\in \N$. If there exists some $c>0$ such that $a_n \le c b_n$ for all $n$, then the convergence of $\sum b_n$ implies the convergence of $\sum a_n$. (or equivalently, the divergence of $\sum a_n$ implies the divergence of $\sum b_n$).
        
        This is called the \textbf{comparison test}.
    \end{mythm}
    \begin{proof}
        Suppose $\displaystyle S_n = \sum_{k=1}^n a_k$ and $\displaystyle T_n = \sum_{k=1}^n b_k$. Then we have $S_n \le c T_n$ for any $n\in \N$. If $\{T_n\}$ converges, then $\{S_n\}$ is bounded and thus $\{S_n\}$ converges.
    \end{proof}
    \begin{myex}
        Recall that we already have $\displaystyle \sum \frac{1}{n^2+n}$ convergent. Note that we have
        \[
        \lim_{n\to \infty} \frac{1/n^2}{1/(n^2+n)} = \lim_{n\to \infty} \left( 1+\frac1n \right) = 1 >0
        \]
        Thus $\displaystyle \sum \frac1{n^2}$ converges. We define \textbf{Riemann zeta function} by
        \[
        \zeta(s) = \sum_{n=1}^\infty \frac{1}{n^s}.
        \]
        We have $\zeta(2) = \displaystyle \sum_{n=1}^\infty \frac{1}{n^2}$ converges. We define the domain $\operatorname{Dom} \left( \zeta \right)$ to be the set of all $s\in \R$ such that $\zeta(s)$ converges.
    \end{myex}
    \begin{mythm}
        Suppose $f$ is a positive decreasing function defined for all $x\ge 1$. For each positive integer $n$, let $S_n = \displaystyle \sum_{k=1}^n f(k)$ and $\displaystyle t_n = \int_1^n f(x)\dd x$. Then the two sequences $\{S_n\}$ and $\{t_n\}$ both converge or both diverge.

        This is called the \textbf{integral test}.
    \end{mythm}
    \begin{proof}
        Geometrically, we have
        \[
        S_n - f(1) = \sum_{k=2}^n f(k) \le \int_1^n f(x)\dd x = t_n \le \sum_{k=1}^{n-1} f(k) = S_{n-1}.
        \]
        Equivalently, we have
        \[
        t_n + f(n) \le S_n \le t_n + f(1).
        \]
    \end{proof} 
    \begin{myex}
        Let $f(x) = x^{-s}$, then
        \[
        t_n = \int_1^n x^{-s} \dd x = \begin{cases}
        \dfrac{n^{1-s} - 1}{1-s} & \text{if } s\neq 1\\[6pt]
        \ln n & \text{if } s=1
        \end{cases}
        \]

        If $s>1$, then $n^{1-s} \to 0$ as $n\to +\infty$, so $t_n \to \dfrac{1}{s-1}$.

        If $s\le 1$, then $t_n$ diverges. 

        Recall the definition of Riemann zeta function $\zeta(s) = \displaystyle \sum_{n=1}^\infty \frac{1}{n^s}$, we have $\zeta(s)$ converges if and only if $s>1$ and $\zeta(s)$ diverges if $s\le 1$.
    \end{myex}
    \begin{myrmk}
        The integral test only tells us about the convergence or divergence of the series but does not give us the value of the sum when the series converges (they are not necessarily equal)
    \end{myrmk}
    \begin{mythm}
        Suppose $a_n \ge 0$ for any $n\in \N$ and additionally, $a_n^{1/n} \to r$ as $n\to \infty$ for some $r\in \R$. Then
        \begin{enumerate}[label=(\arabic*)]
            \item If $r<1$, then $\displaystyle \sum a_n$ converges.
            \item If $r>1$, then $\displaystyle \sum a_n$ diverges.
            \item If $r=1$, then the test is inconclusive/fails.
        \end{enumerate}
        This is called the \textbf{root test}.
    \end{mythm}
    \begin{proof}
        If $r<1$, then there exists some $x$ such that $r<x<1$. Since $a_n^{1/n} \to r$, there exists some $N\in \N$ such that $a_n^{1/n} < x$ for any $n > N$. Thus we have $a_n < x^n$ for any $n > N$. By the comparison test, we have $\displaystyle \sum a_n$ converges.

        If $r>1$, then there exists $\tilde{N}$ such that $a_n^{1/n} >1$ for all $n\ge \sqrt{\tilde{N}}$. Then $a_n>1$ for all $n\ge \tilde{N}$ and thus $\displaystyle \lim_{n\to \infty} a_n$ cannot be $0$ and $\displaystyle \sum a_n$ diverges.

        If $r=1$, it suffices to give two examples.

        \textbf{Example 1:} Suppose $\displaystyle a_n = \left( \frac{1}{n^2} \right) ^{1/n}$, then
        \[
        \lim_{n\to \infty} \left( \frac{1}{n^2} \right) ^{1/n} = \lim_{x\to 0} \left( x^x \right)^2 = 1.
        \]
        However, we know $\displaystyle \sum a_n = \zeta(2)$ is convergent.

        \textbf{Example 2:} Simply letting $b_n=1$ gives a valid example.
        
    \end{proof}
    \begin{mythm}
        Suppose $a_n>0$ for all $n\ge 1$ and additionally $\dfrac{a_{n+1}}{a_n} \to r$ as $n\to \infty$. Then
        \begin{enumerate}[label=(\arabic*)]
            \item If $r<1$, then $\displaystyle \sum a_n$ converges.
            \item If $r>1$, then $\displaystyle \sum a_n$ diverges.
            \item If $r=1$, then the test is inconclusive/fails.
        \end{enumerate}
        This is called the \textbf{ratio test}.
    \end{mythm}
    \begin{mydef}
        A series of the form $\displaystyle \sum_{n=1}^\infty (-1)^{n-1} a_n$ or $\displaystyle \sum_{n=1}^\infty (-1)^{n} a_n$ where $a_n > 0$ is called an \textbf{alternating} series.
    \end{mydef}
    \begin{mythm}
        If $\{a_n\}$ is a decreasing positive sequence with $\displaystyle \lim_{n\to \infty} a_n = 0$, then the alternating series $\displaystyle \sum_{n=1}^\infty (-1)^{n-1} a_n$ converges. In the case where it does converge, let $S$ to denote the sum and $S_n$ denote the partial sum, then
        \[
        0 \le (-1)^n (S-S_n) = a_{n+1} - a_{n+2} + a_{n+3} + \cdots \le a_{n+1} \quad \text{for all } n\ge 1.
        \] 
        This is called the \textbf{Leibniz test} or the \textbf{alternating series test}.
    \end{mythm}
    \begin{proof}
        Note $$S_{2n+2} - S_{2n} = a_{2n+1} - a_{2n+2} \ge 0$$
        and $$ S_{2n+1} - S_{2n-1} = a_{2n+1} - a_{2n} \le 0$$.
        For $\left\{ S_{2n} \right\}$,
        \[
        S_2 \le S_{2n} = a_1 - a_2 + a_3 - \cdots - a_{2n-2} + a_{2n-1} - a_{2n} \le a_1.
        \]
        For $\left\{ S_{2n-1} \right\}$,
        \[
        S_2 \le S_2 + (a_3 - a_4) + \cdots+ (a_{2n-3} - a_{2n-2}) + a_{2n-1} = S_{2n-1} \le S_1.
        \]
        By the monotone convergence theorem, we know $S_{2n} \to S_e$ and $S_{2n-1} \to S_o$ where $S_e, S_o\in \R$ as $n\to \infty$. Since
        \[
        S_e -  S_o = \lim_{n\to \infty} (S_{2n} - S_{2n-1}) = \lim (-1)^{2n-1} a_{2n} = 0,
        \]
        we find $S_e = S_o := S$. And $S_{2n} \to S, S_{2n-1} \to S$ as $n\to \infty$ tells us $S_n\to S$ as $n\to \infty$. Then $\displaystyle \sum_{n=1}^\infty (-1)^{n-1} a_n$ converges.

        Since $\left\{ S_{2n} \right\}$ increases and $\left\{ S_{2n-1} \right\}$ decreases, we have
        \[
        S_{2n} \le S_{2n+2} \le S \le S_{2n+1} \le S_{2n-1}.
        \]
        Thus,
        \[
        0 \le S-S_{2n} \le S_{2n+1} -S_{2n} = a_{2n+1}
        \]
        \[
        0 \le S_{2n-1} - S \le S_{2n-1} - S_{2n} = a_{2n}
        \]
        Conbining the two sets of inequalities we get
        $$ 0 \le (-1)^k (S-S_k) \le a_{k+1}.$$
    \end{proof}
    \begin{myex}
        The harmonic series
        $$ 1+\frac12 +\frac13 + \cdots \to \infty.$$
        However the alternating harmonic series
        $$ 1-\frac12+\frac13-\cdots $$ converges. This motivates to think about the relationship between $\sum |a_n|$ and $\big|\sum a_n\big|$.
    \end{myex}
    \begin{mydef}
        $\sum a_n$ is called \textbf{absolutely convergent} if $\big|\sum a_n\big|$ converges. $\sum a_n$ is called \textbf{conditionally convergent} if $\sum a_n$ converges but  $\big|\sum a_n\big|$ diverges.
    \end{mydef}
    \begin{mythm}
        If $\sum |a_n|$ converges, then $\sum a_n$ also converges with
        \[
        \Big|\sum_{n=1}^\infty a_n\Big| \le \sum_{n=1}^\infty |a_n|
        \]
    \end{mythm}
    \begin{proof}
        Let $b_n = a_n + |a_n| \ge 0$. We will show $\sum b_n$ converges. Note that
        \[
        b_n = \begin{cases}
            0 \quad & \text{if } a_n < 0\\
            2a_n \quad & \text{if } a_n \ge 0
        \end{cases}
        \]
        and thus $0\le b_n \le 2|a_n|$. Thus using the comparison test, we have $\sum b_n$ converges. Since $\sum a_n = \sum b_n - \sum |a_n|$, we have $\sum a_n$ congerges.

        Since $\displaystyle \Big|\sum_{k=1}^n a_k \Big| \le \sum_{k=1}^n |a_k|$ for all $n$, taking the limit as $n \to \infty$ (provided that they both exist), we get
        \[
        \Big|\sum_{n=1}^\infty a_n\Big| \le \sum_{n=1}^\infty |a_n|.
        \]
    \end{proof}
    \begin{myex}
        The alternating harmonic series converges conditionally. The series $\displaystyle \sum_{n=1}^\infty  \frac{(-1)^{n-1}}{n^2}$ converges absolutely.
    \end{myex}
    \begin{mypro}
        If $\sum a_n$ and $\sum b_n$ both are absolutely convergent, then $\sum (\alpha a_n + \beta b_n)$ is also absolutely convergent for any $\alpha,\beta \in \R$.
        $$ \sum_{k=1}^n|\alpha a_n + \beta b_n| \le \alpha \left( \sum_{k=1}^n |a_n| \right) + \beta \left( \sum_{k=1}^n |b_n| \right) \le \alpha \sum_{k=1}^{\infty} |a_n| + \beta \sum_{k=1}^\infty |b_n| =: S.$$
        Since all partial sum are bounded above by $S$, we have $\sum |\alpha a_n + \beta b_n|$ converges and thus $\sum (\alpha a_n + \beta b_n)$ is absolutely convergent.
    \end{mypro}
    \begin{mylma}
        Suppose two sequences $\{a_n\}$ and $\{b_n\}$. Let $\displaystyle A_n = \sum_{k=1}^n a_k$. Then
        \[
        \sum_{k=1}^{n} a_k b_k = A_n b_{n+1} + \sum_{k=1}^n A_k (b_k - b_{k+1}).
        \]
        This is called \textbf{Abel's partial sum formula}. 
    \end{mylma}
    \begin{proof}
        Let $A_0=0$, then $a_k = A_k - A_{k-1}$ and 
        \[
        \begin{aligned}
            \sum_{k=1}^n a_k b_k = \sum_{k=1}^n A_k b_k - \sum_{k=1}^n A_{k-1} b_k &= \sum_{k=1}^n A_k b_k - \left[ \sum_{k=1}^n A_kb_{k+1} - A_nb_{n+1} \right]\\ &= A_n b_{n+1} + \sum_{k=1}^n A_k (b_k - b_{k+1}).
        \end{aligned}
        \]
    \end{proof}
    \begin{myrmk}
        If both $\{A_nb_{n+1}\}$ and $\displaystyle \sum_{k=1}^n A_k (b_k - b_{k+1})$ converge, then $\displaystyle \sum_{k=1}^n a_k b_k$ converges.
    \end{myrmk}
    
    \begin{mythm}
        Suppose $\sum a_n$ is a series whose partial sum $\{A_n\}$ is bounded. Suppose $\{b_n\}$ is a decreasing sequence with $b_n \to 0$ as $n\to \infty$. Then the series $ \sum a_n b_n$ converges.

        This  is called the \textbf{Dirichlet's test}.
    \end{mythm}
    \begin{proof}
        Let $M>0$ be such that $|A_n| \le M$ for any $n\in \N$. We tackle the convergence pf $\{A_n b_{n+1}\}$ and $\displaystyle \sum_{k=1}^n A_k (b_k - b_{k+1}) $ seperately. 
        
        \begin{itemize}
            \item Note that $\{A_nb_{n+1}\}$ converges becuase $0 \le |A_n b_{n+1} | \le M |b_{n+1}| \to 0$ as $n\to \infty$. 
            \item For the convergence $\sum A_k(b_k - b_{k+1})$, we have
            \[
            0 \le |A_k(b_k - b_{k+1})| \le M|b_k - b_{k+1}| = M(b_k - b_{k+1}).
            \]
            $\sum (b_k - b_{k+1})$ converges since it is a telescoping series with $\displaystyle \lim_{n\to \infty} b_n$ exists. Thus $\sum A_k(b_k - b_{k+1})$ converges by the comparison test.
        \end{itemize}
        By Abel's partial sum formula, we have $\sum a_n b_n$ converges.
    \end{proof}
    \begin{mythm}
        If $\sum a_n$ converges and $\{b_n\}$ is monotonic and convergent, then $\sum a_n b_n$ converges.

        This is called the \textbf{Abel's test}.
    \end{mythm}
    \begin{myrmk}
        If $\{b_n\}$ is decreasing and $b_n \to 0$ as $n\to \infty$. Taking $a_n = x^n$, we know the partial sum $\displaystyle A_n = \sum_{k=1}^n x^k$ is bounded for any $x\in [-1,1)$. Particularly, when $x=-1$, Dirichlet's test  says that $\displaystyle \sum_{n=1}^\infty a_nb_n = \sum_{n=1}^\infty (-1)^n b_n$ converges, which is simply the alternating series test or the Leibniz test.
    \end{myrmk}

    \subsection{Improper Integrals}
    \begin{mydef}
        The symbol $\displaystyle \int_a^\infty f(x) \dd x$ means $\displaystyle \lim_{b\to \infty} \int_a^b f(x) \dd x$ and the symbol $\displaystyle \int_{-\infty}^b f(x) \dd x$ means $\displaystyle \lim_{a\to -\infty} \int_a^b f(x) \dd x$. They are referred to as \textbf{improper integrals of the first kind}. We say the improper integral $\displaystyle \int_a^\infty f(x) \dd x$ converges if the limit exists and diverges otherwise. If the limit does exist and equal to $A$, then we call $A$ the value of the improper integral and write $\displaystyle \int_a^\infty f(x) \dd x=A$.

        If $\displaystyle \int_{-\infty}^c f(x) \dd x$ and $\displaystyle \int_c^\infty f(x) \dd x$ both converge for some $c\in \R$, then we say the improper integral $\displaystyle \int_{-\infty}^\infty f(x) \dd x$ also converges. And we say $\displaystyle \int_{-\infty}^\infty f(x) \dd x$ diverges if at least one of $\displaystyle \int_{-\infty}^c f(x) \dd x$ and $\displaystyle \int_c^\infty f(x) \dd x$ diverges.
    \end{mydef}
    \begin{myrmk}
        Note that
        \[
        \int_{-\infty}^\infty f(x) \dd x = \int_{-\infty}^c f(x) \dd x + \int_c^d f(x) \dd x + \int_d^\infty f(x) \dd x \quad \left( \forall c<d\in \R \right).
        \]
        Combining the first two terms and last two terms gives us
        \[
        \int_{-\infty}^c f(x) \dd x + \int_c^\infty f(x) \dd x = \int_{-\infty}^d f(x) \dd x + \int_d^\infty f(x) \dd x.
        \]
        Thus we just need to find one point $c\in \R$ such that $\displaystyle \int_{-\infty}^c f(x) \dd x$ and $\displaystyle \int_c^\infty f(x) \dd x$ both converge to determine the convergence of $\displaystyle \int_{-\infty}^\infty f(x) \dd x$.
    \end{myrmk}
    \begin{myex}
        We introduces some typical examples of improper integrals:
        \begin{enumerate}[label=(\arabic*)]
            \item For a real $s$,
            \[
            \int_1^\infty x^{-s} \dd x = \begin{cases}
            \dfrac{x^{1-s}}{1-s} \Big|_1^\infty = \dfrac{1}{s-1} & \text{if } s>1\\[6pt]
            \ln x \Big|_1^\infty = +\infty & \text{if } s=1\\[6pt]
            \dfrac{x^{1-s}}{1-s} \Big|_1^\infty = +\infty & \text{if } s<1
            \end{cases}.
            \]
            \item For a real positive $a$,
            \[
            \int_{-\infty}^\infty e^{-a|x|} \dd x = \frac{2}{a}
            \]
            \item $\displaystyle \int_{-\infty}^\infty x \dd x$ diverges since $\displaystyle \int_c^\infty x \dd x$ and $\displaystyle \int_{-\infty}^c x \dd x$ diverges for any $c\in \R$.
            
            In this case, we cannot cancel the two divergent integrals like
            \[
            \lim_{b\to \infty} (\int_{-b}^0 x \dd x + \int_0^b x \dd x) = \lim_{b\to \infty} \int_{-b}^b x \dd x = 0.
            \]
            
            \end{enumerate}
    \end{myex}
    \begin{mythm}
        Suppose $f(x)\ge 0$ for all $x\ge a$ and the proper integral $\displaystyle \int_a^b f(x) \dd x$ exists for any $b\ge a$. Then the improper integral $\displaystyle \int_a^\infty f(x) \dd x$ converges if and only if there is an $M>0$ such that $\displaystyle \int_a^b f(x) \dd x \le M$ for any $b\ge a$.
    \end{mythm}
    \begin{mythm}
        Suppose $f(x)\ge 0$ and $g(x)\ge 0$ for all $x\ge a$ and the proper integrals $\displaystyle \int_a^b f(x) \dd x$ and $\displaystyle \int_a^b g(x) \dd x$ exist for any $b\ge a$. If there exists some $c>0$ such that $f(x) \le c g(x)$ for all $x\ge a$, then the convergence of $\displaystyle \int_a^\infty g(x) \dd x$ implies the convergence of $\displaystyle \int_a^\infty f(x) \dd x$. (or equivalently, the divergence of $\displaystyle \int_a^\infty f(x) \dd x$ implies the divergence of $\displaystyle \int_a^\infty g(x) \dd x$). Particularly, if $\displaystyle \int_a^\infty g(x) \dd x$ converges, then $\displaystyle \int_a^\infty f(x) \dd x \le c \int_a^\infty g(x) \dd x$.
    \end{mythm}
    \begin{mythm}\label{limit_test_int}
        Suppose that $f(x)>0$ and $g(x)>0$ for all $x\ge a$ and the proper integrals $\displaystyle \int_a^b f(x) \dd x$ and $\displaystyle \int_a^b g(x) \dd x$ exist for any $b\ge a$. If $\displaystyle \lim_{x\to \infty} \frac{f(x)}{g(x)} = c>0$, then $\displaystyle \int_a^\infty f(x) \dd x$ converges if and only if $\displaystyle \int_a^\infty g(x) \dd x$ converges.
    \end{mythm}
    \begin{myrmk}
        Consider the same set-up in Theorem ~\ref{limit_test_int} but with $f(x)\ge 0$ and $\displaystyle \lim_{x\to +\infty} \frac{f(x)}{g(x)} = c = 0$. Intuitively the conclusion should be weaker: the convergence of $\displaystyle \int_a^\infty g(x) \dd x$ implies the convergence of $\displaystyle \int_a^\infty f(x) \dd x$ but the converse may not hold. A trivial example is $f(x)=0$ and $g(x)=x$.
    \end{myrmk}
    \begin{mydef}
        Suppose $f$ is defined on $(a,b]$ where $b>a$ and the proper integral $\displaystyle \int_x^b f(t) \dd t$ exists for any $x\in (a,b]$. Then the symbol $\displaystyle \int_{a^+}^b f(t) \dd t$ means $\displaystyle \lim_{x\to a^+} \int_x^b f(t) \dd t$ and is referred to as an \textbf{improper integral of the second kind}. We say the improper integral $\displaystyle \int_{a^+}^b f(t) \dd t$ converges if the limit exists and diverges otherwise. If the limit does exist and equal to $A$, then we call $A$ the value of the improper integral and write $\displaystyle \int_{a^+}^b f(t) \dd t = A$.
    \end{mydef}
    \begin{mydef}
        The symbol $\displaystyle \int_{a}^{b^-} f(t) \dd t$ is defined similarly. If both $\displaystyle \int_{a^+}^c f(t) \dd t$ and $\displaystyle \int_c^{b^-} f(t) \dd t$ converge, then we say $\displaystyle \int_{a^+}^{b^-} f(t) \dd t$ also converges and write
        \[
        \int_{a^+}^{b^-} f(t) \dd t = \int_{a^+}^c f(t) \dd t + \int_c^{b^-} f(t) \dd t.
        \]
    \end{mydef}
    
    \subsection{Sequences and series of functions}
    \begin{mydef}
        Suppose for each positive integer $n$, $f_n$ is a function. We then call $\{f_n\}$ a \textbf{sequence of functions}. If there exists a function $f$ and a set $S\subseteq \R$ such that
        \[
        f(x) = \lim_{n\to \infty} f_n(x) \quad \text{for all } x\in S,
        \]
        then we say $f$ \textbf{the limit function of $\{f_n\}$ on $S$} and we say that $\{f_n\}$ \textbf{converges pointwise to $f$ on $S$}. We often simply write
        \[
        f_n \to f \text{ pointwise on } S.
        \]
    \end{mydef}
    \begin{myrmk}
        Here we introduce the ``$\varepsilon -N$'' definition for the pointwise convergence of sequence of functions:
        \[
        \forall x\in S, \forall \varepsilon >0, \exists N\in \N \text{ such that } |f_n(x) - f(x)| < \varepsilon \text{ for all } n > N.
        \]
        Note that the choice of $N$ depends on $x$ and $\varepsilon$ because when $x$ takes different values, we actually have different sequences $\{f_n(x)\}$.(it's a sequence of real numbers now!!)
    \end{myrmk}
    \begin{myex}
        Let $f_n(x) = x^n$ for $x\in [0,1]$ and all $n\in \N$. Then $\{f_n\}$ converges pointwise to the function
        \[
        f(x) = \begin{cases}
        0 & \text{if } x\in [0,1)\\
        1 & \text{if } x=1
        \end{cases}.
        \]
        Note $f_n$ is continuous on $[0,1]$ but the limit function $f$ is not continuous on $[0,1]$. Thus $f$ doesn't preserve the continuity of $f_n$, which implies that the pointwise definition is not strong enough. This is because there exists a particular $x=1$ such that $f_n(1) = 1$ for all $n\in \N$ and thus the convergence is not uniform. However, our definition treats all $x\in [0,1]$ equally. This motivates us to introduce the concept of uniform convergence.
    \end{myex}
    \begin{mydef}
        A sequence of functions $\{f_n\}$ is said to \textbf{converge uniformly} to a function $f$ on a set $S$ if
        \[
        \forall \varepsilon >0, \exists N\in \N \text{ such that } \forall x\in S, |f_n(x) - f(x)| < \varepsilon \text{ for all } n > N.
        \]
        We often simply write
        \[
        f_n \to f \text{ uniformly on } S.
        \]
    \end{mydef}
    \begin{myrmk}
        Note that the choice of $N$ only depends on $\varepsilon$ but not on $x$. Thus we can choose a single $N$ for all $x\in S$ to make sure $|f_n(x) - f(x)| < \varepsilon$ for all $n > N$. In other words, the convergence is uniform in the sense that it is uniform across all $x\in S$.
    \end{myrmk}
    \begin{mythm}
        If $f_n \to f$ uniformly on $S\subseteq \R$ and for each $n\ge 1$, $f_n$ is continuous at $x=a \in S$, then then limit function $f$ is also continuous at $x=a$.
    \end{mythm}
    \begin{proof}
        We need to show that $\displaystyle \lim_{x\to a} f(x) = f(a)$. 

        Given $\varepsilon >0$, since $f_n \to f$ uniformly on $S$, there exists an $N\in \N$ such that $|f_n(x) - f(x)| < \varepsilon/3$ for all $x\in S$ and all $n \ge N$.

        Since $f_N$ is continuous at $x=a$, there exists some $\delta >0$ such that $|f_N(x) - f_N(a)| < \varepsilon/3$ for all $x\in S$ with $|x-a| < \delta$. Note we've chosen fixed $N$ first and then $\delta$ is also fixed. 

        For such $x$ and $n\ge N$, we have
        \[
        |f(x) - f(a)| \le |f(x) - f_N(x)| + |f_N(x) - f_N(a)| + |f_N(a) - f(a)| < \varepsilon/3 \times 3 = \varepsilon.
        \]
    \end{proof}
    \begin{mypro}
        Particularly, if $f_n=\displaystyle \sum_{k=1}^n u_k$ is the sequence of partial sums of another function sequence $\{u_n\}$, then the above theorem can be written as:

        If a series of function $\displaystyle \sum u_n$ converges uniformly to a sum function $f$ on $S\subseteq \R$ and for each $n\ge 1$, $u_n$ is continuous at $x=a \in S$, then the sum function $f$ is also continuous at $x=a$.

        We often just write the conclusion as
        \[
        \lim_{x\to a} \sum_{n=1}^\infty u_n(x) = \lim_{x\to a} f(x) = f(a)= \sum_{n=1}^\infty u_n(a) = \sum_{n=1}^\infty \lim_{x\to a} u_n(x).
        \]
    \end{mypro}
    \begin{mythm}
        Suppose $f_n \to f$ uniformly on $[a,b]$ and for each $n\ge 1$, $f_n$ is continuous on $[a,b]$. Define another function $\{g\}$ as $g_n(x) = \displaystyle \int _a^x f_n(t) \dd t$ for $x\in [a,b]$ and function $g$ as $g(x) = \displaystyle \int_a^x f(t) \dd t$ for $x\in [a,b]$. Then $g_n \to g$ uniformly on $[a,b]$ we have
        \[
        \lim_{n\to \infty} \int_a^x f_n(t) \dd t = \int_a^x \lim_{n\to \infty} f_n(t) \dd t.
        \]
    \end{mythm}
    \begin{proof}
        Given $\varepsilon >0$, there is an $N$ such that $|f_n(x) - f(x)| < \varepsilon/(b-a)$ for all $x\in [a,b]$ whenever $n\ge N$. Then for any $x\in [a,b]$ we have
        \[\begin{aligned}
            |g_n(x) - g(x)| &= \left| \int_a^x f_n(t) \dd t - \int_a^x f(t) \dd t \right|\\
            &\le \int_a^x |f_n(t) - f(t)| \dd t\\
            &< \int_a^x \frac{\varepsilon}{b-a} \dd t = \frac{\varepsilon}{b-a} (x-a) \le \varepsilon.
        \end{aligned}\]
    \end{proof}
    \begin{mypro}
        Particularly, if $f_n=\displaystyle \sum_{k=1}^n u_k$ where $\{u_n\}$ is another function sequence, then the above theorem can be written as:

        Suppose a series of functions $\sum u_n$ converges uniformly to a sum function $f$ on $[a,b]$ and for each $k\ge 1$, $u_k$ is continuous on $[a,b]$. Then we have
        \[
        g_n(x) = \int_a^x \sum_{k=1}^n u_k(t) \dd t,\qquad g(x) = \int_a^x \sum_{k=1}^\infty u_k(t) \dd t
        \]
        Then $g_n \to g$ uniformly on $[a,b]$ and
        \[
        \lim_{n\to \infty} \int_a^x f_n(t) \dd t = \int_a^x \lim_{n\to \infty} f_n(t) \dd t.
        \]
        Note we have the left hand side as
        \[
        \lim_{n\to \infty} \int_a^x f_n(t) \dd t = \lim_{n\to \infty} \int_a^x \sum_{k=1}^n u_k(t) \dd t = \lim_{n\to \infty} \sum_{k=1}^n \int_a^x u_k(t) \dd t = \sum_{k=1}^\infty \int_a^x u_k(t) \dd t,
        \]
        while the right hand side is
        \[
        \int_a^x \lim_{n\to \infty} f_n(t) \dd t = \int_a^x \sum_{k=1}^\infty u_k(t) \dd t.
        \]
        This also suggests that we can integrate term-by-term for a uniformly convergent series of functions, i.e.
        \[
        \sum_{k=1}^\infty \int_a^x u_k(t) \dd t = \int _a^x\sum_{k=1}^\infty u_k(t) \dd t.
        \]

    \end{mypro}
    \begin{myrmk}
        We do not have such nice property (i.e. differentiating term-by-term) for general uniformly convergent series. For example, let $f_n(x) = \displaystyle \frac{\sin nx}{n^2}$. Note that $\sum f_n$ is convergent since $\displaystyle 0\le |f_n(x)| \le \frac{1}{n^2}$ and $\displaystyle \sum \frac{1}{n^2}$ converges. Then by the Weierstrass M-test (which will be introduced later), we have $\sum f_n$ converges uniformly.

        Term-by-term differentiation gives us $\displaystyle \sum f_n'(x) = \sum \frac{\cos nx}{n}$, which diverges at $x=0$ since $\displaystyle \sum \frac{1}{n}$ diverges.
    \end{myrmk}
    \begin{mythm}
        Given a series of functions $\sum u_n$ that converges pointwise to a function $f$ on $S\subseteq \R$. If there is a convergent series of positive numbers $\sum M_n$ such that $|u_n(x)| \le M_n$ for all $x\in S$ and all $n\in \N$, then $\sum u_n$ converges uniformly to $f$ on $S$.
        
        This is called the \textbf{Weierstrass M-test}.
    \end{mythm}
    \begin{proof}
        Given $\varepsilon >0$, we need to show that there is an $N\in \N$ such that $\Big|\displaystyle \sum_{k=1}^n u_k(x) - f(x)\Big| < \varepsilon$ for all $x\in S$ whenever $n \ge N$.

        Suppose $\displaystyle \sum_{k=1}^\infty M_k=A$, then for the given $\varepsilon >0$, there is an $N\in \N$ such that $\displaystyle \Big| \sum_{k=1}^n M_k - A \Big| < \varepsilon$ for all $n \ge N$.

        $$ \Big|f(x) - \sum_{k=1}^n u_k(x)\Big| = \Big|\sum_{k=n+1}^\infty u_k(x)\Big| \le \sum_{k=n+1}^\infty |u_k(x)| \le \Big|\sum_{k=n+1}^\infty M_k \Big| = \Big| A - \sum_{k=1}^n M_k \Big| < \varepsilon. $$
        holds for any $x\in S$ and all $n \ge N$. Thus $\sum u_n$ converges uniformly to $f$ on $S$.
    \end{proof}

    \subsection{Power Series}
    Recall in Definition~\ref{def:power-series}, we have defined the power series $\displaystyle \sum_{n=0}^\infty a_n x^n$. Now we will discuss the convergence of power series.
    \begin{mydef}
        In general, a series of the form $\displaystyle \sum_{n=0}^\infty a_n (x-a)^n$ is also called a \textbf{power series} in $(x-a)$ or a \textbf{power series centered at} $x=a$. We call the $a$ the center and $a_n$ the coefficient of the power series.
    \end{mydef}
    \begin{myex}
        \begin{enumerate}
            \item For the power series $\displaystyle \sum_{n=0}^\infty \frac{x^n}{n!}$, ratio test gives us
            $$ \big| \frac{x^{n+1}}{(n+1)!} \cdot \frac{n!}{x^n} \big| = \frac{|x|}{n+1} \to 0 < 1\quad \text{for } x\neq 0$$
            as $n\to \infty$. The case where $x=0$ is trivial. Thus $\displaystyle \sum_{n=0}^\infty \frac{x^n}{n!}$ converges for all $x\in \R$.
            \item For the power series $\displaystyle \sum_{n=0}^\infty n^23^n x^n$, root test gives us
            $$ \sqrt[n]{|n^23^n x^n|} = 3|x| \quad \text{for } x\in \R \text{ as } n\to \infty$$
            Thus $\displaystyle \sum_{n=0}^\infty n^23^n x^n$ converges for $|x| < \frac{1}{3}$ and diverges for $|x| > \dfrac{1}{3}$. The case where $|x| = \dfrac{1}{3}$ is inconclusive and we need to check it manually. When $x=\dfrac{1}{3}$, we have $\displaystyle \sum_{n=0}^\infty n^2$ which diverges. When $x=-\dfrac{1}{3}$, we have $\displaystyle \sum_{n=0}^\infty (-1)^n n^2$ also diverges.
        \end{enumerate}
    \end{myex}
    \begin{mythm}\label{convergence_power_series}
        If a power series $\sum a_n x^n$ converges for a particular $x_1 \neq 0$, then
        \begin{enumerate}[label=(\alph*)]
            \item the power series converges absolutely for every $x$ with $|x| < |x_1|$;
            \item the power series converges uniformly on every interval $[-R,R]$ with $0<R<|x_1|$.
        \end{enumerate}
    \end{mythm}
    \begin{proof}
        We first prove (b). Note the necessary condition for the convergence of $\sum a_n x^n$ is that $a_n x^n \to 0$ as $n\to \infty$. Particularly, there is an $N\in \N$ such that $|a_n x_1^n| < 1$ whenever $n\ge N$. Then for $x\in [-R,R]$ with $0<R<|x_1|$, we have
        $$ 0\le |a_n x^n| = |a_n||x|^n = |a_n x_1^n| \cdot \left|\frac{x}{x_1}\right|^n < \left|\frac{x}{x_1}\right|^n\le \Big| \frac{R}{x_1}\Big|^n < 1 \quad \text{for } n\ge N. $$
        Then by the comparison test, we have $\sum a_n x^n$ converges absolutely on $[-R,R]$. Also by the Weierstrass M-test(in this case, $M_n = \displaystyle\Big| \frac{R}{x_1}\Big|^n$), we have $\sum a_n x^n$ converges uniformly on $[-R,R]$.

        Now we prove (a). For any $|x|<|x_1|$, we can find some $R$ such that $|x| < R < |x_1|$. Then by (b), we have $\sum a_n x^n$ converges absolutely for every $x$ with $|x| < R$. In particular, $\sum a_n x^n$ converges absolutely for every $x$ with $|x| < |x_1|$.
    \end{proof}
    \begin{mythm}\label{radius_of_convergence}
        If a power series $\sum a_n x^n$ converges for a particular $x_1 \neq 0$ and diverges for a particular $x_2 \neq 0$, then there is a unique $r$ such that
        \begin{enumerate}[label=(\alph*),start=3]
            \item the power series converges absolutely for every $x$ with $|x| < r$;
            \item the power series diverges for every $x$ with $|x| > r$.
        \end{enumerate}
        We call $r$ the \textbf{radius of convergence} of the power series since this $r$ is unique.
    \end{mythm}
    \begin{proof}
        We prove (d) first. Let $A=\{ |x| : \displaystyle \sum a_n x^n \text{ converges at } x\}$. Note $|x_1|\in A$ and $|x_2|\notin A$. By (a) we know $|x_2|$ is an upper bound of $A$. Let $r = \sup A$, then $r$ is the unique number. Note $0<|x_1| \le r \le |x_2|$. Then for any $|x|>r$, $|x|\notin A$ and thus $\sum a_n x^n$ diverges at such $x$.

        Now we prove (c). For $0<|x|<r$, $|x|$ is not an upper bound of $A$ and there exists $|z| \in A$ such that $0<|x| < |z| < r$. By (a), we have $\sum a_n x^n$ converges absolutely for every $x$ with $|x| < |z|$. In particular, $\sum a_n x^n$ converges absolutely for every $x$ with $|x| < r$.
    \end{proof}
    \begin{myrmk}
        \leavevmode
        \begin{itemize}
            \item If the power series $\sum a_n x^n$ converges for all $x\in \R$, then we say the radius of convergence is $r = +\infty$.
            \item If the power series $\sum a_n x^n$ diverges for all $x\neq 0$, then we say the radius of convergence is $r=0$.
        \end{itemize}
    \end{myrmk}
    \begin{myrmk}
        Theorem~\ref{convergence_power_series} and Theorem~\ref{radius_of_convergence} also hold for power series of the form $\displaystyle \sum_{n=0}^\infty a_n (x-a)^n$ by simply replacing $x$ with $x-a$ in the proofs so the power series centers at $x=a$ instead of $x=0$. Let $r$ be the \textbf{radius of convergence}, we often call $(a-r,a+r)$ the \textbf{interval of convergence} of the power series $\displaystyle \sum_{n=0}^\infty a_n (x-a)^n$. Note the convergence at the endpoints $x=a\pm r$ is inconclusive and we need to check it manually.
    \end{myrmk}
    \begin{mydef}
        On the interval of convergence of the power series $\displaystyle \sum_{n=0}^\infty a_n (x-a)^n$, we write $f(x) = \displaystyle \sum_{n=0}^\infty a_n (x-a)^n$ and call $f$ the \textbf{sum function}. We also say that the power series \textbf{represent the function $f$} in its interval of convergence and we call it the \textbf{power series expansion} of $f$ at $x=a$. 
    \end{mydef}
    \begin{myrmk}
        By Theorem~\ref{convergence_power_series} and Theorem~\ref{radius_of_convergence}, given a power series $\displaystyle \sum_{n=0}^\infty a_n (x-a)^n$, let $r$ be the radius of convergence, we know the power series
        \begin{itemize}
            \item converges absolutely on $(a-r,a+r)$;
            \item converges uniformly on $[a-R,a+R]$ for any $0<R<r$;
        \end{itemize}
        Thus the sum function $f$ is continuous on $(a-r,a+r)$. We also know we can integrate the power series term-by-term on $[a-r,a+r]$ to obtain the integral of the sum function $f$. For any $x\in(a-r,a+r)$ we have
        \[
        \int_a^x f(t) \dd t = \int_a^x \sum_{n=0}^\infty a_n (t-a)^n \dd t = \sum_{n=0}^\infty a_n \int_a^x (t-a)^n \dd t = \sum_{n=0}^\infty a_n \frac{(x-a)^{n+1}}{n+1}.
        \]
    \end{myrmk}
    \begin{mythm}\label{diff_power_series}
        Suppose $f(x) = \displaystyle \sum_{n=0}^\infty a_n (x-a)^n$ on $(a-r,a+r)$ where $r$ is the radius of convergence. Then 
        \begin{enumerate}[label=(\alph*)]
            \item $\displaystyle \sum_{n=1}^\infty n a_n (x-a)^{n-1}$ also has the same interval convergence $(a-r,a+r)$
            \item The derivative $f'(x) $ exists for each $x\in (a-r,a+r)$ and $f'(x) = \displaystyle \sum_{n=1}^\infty n a_n (x-a)^{n-1}$ for each $x\in (a-r,a+r)$.
        \end{enumerate}
    \end{mythm}
    \begin{proof}
        Suppose $a=0$. We start with (a). For any $x\in (0,r)$, we know for any $h>0$ with $0<x<x+h<r$,
        \[
        f(x+h) = \sum_{n=0}^\infty a_n (x+h)^n, \quad \frac{f(x+h)-f(x)}{h} = \sum_{n=0}^\infty a_n \frac{(x+h)^n - x^n}{h}.
        \]
        Let $g_n(x) = x^n$, by the \textbf{MVT} to $g_n$ on $[x,x+h]$ (note $g_n$ is continuous on $[x,x+h]$ and differentiable on $(x,x+h)$), there is some $c_n \in (x,x+h)$ such that
        \[
        g_n(x+h) - g_n(x) = h g'(c_n) \quad \text{where } x<c_n<x+h.
        \]
        Thus we have
        \[
        \frac{f(x+h)-f(x)}{h} = \sum_{n=0}^\infty a_n \frac{(x+h)^n - x^n}{h} = \sum_{n=0}^\infty a_n g'(c_n) = \sum_{n=1}^\infty n a_n c_n^{n-1}\quad \text{converges absolutely}.
        \]
        Note that $\displaystyle \sum_{n=1}^\infty n a_n x^{n-1}$ satisfies
        \[
        0 \le | n a_n x^{n-1}| < |n a_n c_n^{n-1}| \quad \text{for all } n\ge 1.
        \]
        Thus $\displaystyle \sum_{n=1}^\infty n a_n x^{n-1}$ converges absolutely for any $x\in (-r,r)$. Similarly, we know it actually converges absolutely for any $x\in (-r,0)$ and thus for any $x\in (-r,r)$.

        Now we prove $\displaystyle \sum_{n=1}^\infty n a_n (x-a)^{n-1}$ diverges for any $x$ with $|x| > r$. Suppose $\displaystyle \sum_{n=1}^\infty n a_n x^{n-1}$ converges absolutely on $(-r_2,r_2)$ where $r_2 > r$ and $0\le |a_n x^n| < |n a_n x^{n-1}|$ when $n>r_2$. For any $x\in (r,r_2)$, we have $\displaystyle \sum_{n=0}^\infty a_n x^n$ diverges and thus $\displaystyle \sum_{n=1}^\infty n a_n x^{n-1}$ also diverges. Thus the differentiated series $\displaystyle \sum_{n=1}^\infty n a_n x^{n-1}$ share the same interval of convergence as the orginal series $\displaystyle \sum_{n=0}^\infty a_n x^n$.

        Now we prove (b). Let $D(x) = \displaystyle \sum_{n=1}^\infty n a_n x^{n-1}$, then for any $x\in (-r,r)$,
        \[
        \int_0^x D(t) \dd t = \sum_{n=1}^{\infty} a_n \int_0^x n t^{n-1} \dd t = \sum_{n=1}^\infty a_n x^n = f(x)-a_0.
        \]
        Note $D(x)$ is continuous on $(-r,r)$. By the \textbf{FTC1}, we have $D(x) = f'(x)$ for any $x\in (-r,r)$.
    \end{proof}
    \begin{myex}
        Recall
        \begin{equation}
            1+x+x^2 + \cdots = \frac{1}{1-x} \quad r=1. \tag{1}
        \end{equation}
        Differentiating (1) term-by-term gives us
        \begin{equation}
            1 + 2x + 3x^2 + \cdots = \frac{1}{(1-x)^2} \quad r=1. \tag{2}
        \end{equation}
        Replace $x$ by $-x$ in (1) gives us
        \begin{equation}
            1-x+x^2 - \cdots = \frac{1}{1+x} \quad r=1. \tag{3}
        \end{equation}
        Integrating (3) term-by-term gives us
        \begin{equation}
            x - \frac{x^2}{2} + \frac{x^3}{3} - \cdots = \ln(1+x) \quad r=1. \tag{4}
        \end{equation}
        Replace $x$ by $x^2$ in (3) gives us
        \begin{equation}
            1-x^2+x^4 - \cdots = \frac{1}{1+x^2} \quad r=1. \tag{5}
        \end{equation}
        Integrating (5) term-by-term gives us
        \begin{equation}
            x - \frac{x^3}{3} + \frac{x^5}{5} - \cdots = \arctan x \quad r=1. \tag{6}
        \end{equation}
        For (4), when $x=1$, we have
        \[
            1 - \frac{1}{2} + \frac{1}{3} - \cdots = \ln 2.
        \]
        For (6), when $x=1$, we have
        \[
            1 - \frac{1}{3} + \frac{1}{5} - \cdots = \frac{\pi}{4}.
        \]
    \end{myex}
    \begin{myrmk}
        Suppose $f(x) = \displaystyle \sum_{n=0}^\infty a_n (x-a)^n$ on its interval of convergence $(a-r,a+r)$. For any $x\in (a-r,a+r)$, $f'(x), f''(x), \cdots, f^{(k)}(x)$ are all well-defined (and can be obtained by keeping differentiating the power series) for any positive integer $k$. We say that $f$ is infinitely differentiable on $(a-r,a+r)$. Note
        $$ f^{(k)}(x) = \sum_{n=k}^\infty n(n-1)\cdots (n-k+1) a_n (x-a)^{n-k} \quad \text{for any } k\ge 1. $$
        In particular,
        $$ f^{(k)}(a) = k! a_k \Rightarrow a_k = \frac{f^{(k)}(a)}{k!} \quad \text{for any } k\ge 0. $$
        Then $f(x) = \displaystyle \sum_{n=0}^\infty \frac{f^{(n)}(a)}{n!} (x-a)^n$ for any $x\in (a-r,a+r)$. In other words, the power series expansion of $f$ at $x=a$ is actually the Taylor series expansion of $f$ at $x=a$ and thus unique.
    \end{myrmk}
    \begin{mythm}
        If two power series $\displaystyle \sum_{n=0}^\infty a_n (x-a)^n$ and $\displaystyle \sum_{n=0}^\infty b_n (x-a)^n$ have the same sum function $f$ in some neighbourhood of $x=a$, then two series are equal (term by term); in fact, we have $a_n = b_n= \displaystyle \frac{f^{(n)}(a)}{n!}$ for all $n\ge 0$.
    \end{mythm}
    \begin{mypro}
        If we know a function $f$ can be represented by a power series $\displaystyle \sum_{n=0}^\infty a_n (x-a)^n$ in $(a-r,a+r)$, then its sequence of Taylor polynomial $T_n f(x;a)$ converges absolutely on $(a-r,a+r)$ and converges uniformly on $[a-R,a+R]$ for any $0<R<r$.
    \end{mypro}
    \begin{myrmk}
        Suppose $f$ is infinitely differentiable in some open neighbourhood of $x=a$ and consider the power series (Taylor series) $\displaystyle \sum_{n=0}^\infty \frac{f^{(n)}(a)}{n!} (x-a)^n$. Note
        \begin{itemize}
            \item the taylor series does not necessarily converge:
            $\displaystyle
            f(x) = \frac1{1-x}  
            $ is infinitely differentiable and it's Taylor series at $x=0$ is $\displaystyle \sum_{n=0}^\infty x^n$, which diverges for $|x|\ge 1$.

            \item if the taylor series converges, it does not necessarily converge to $f$:
            $$
            f(x) = \begin{cases}
            \exp{-\dfrac{1}{x^2}} & \text{if } x\neq 0\\[6pt]
            0 & \text{if } x=0
            \end{cases}
            $$Note $f$ is infinitely differentiable at $x=0$ and $T_n f(x)\equiv 0$ for all $n\in \N$. Thus the Taylor series of $f$ at $x=0$ converges to the zero function, which is not equal to $f$ itself except for $x=0$.
        \end{itemize}
        Combining with the Taylor formula with error terms, we know
        $$ \sum_{n=0}^\infty \frac{f^{(n)}(a)}{n!} (x-a)^n \to f(x) \Leftrightarrow E_n(x) \to 0 \text{ as } n\to \infty. $$
    \end{myrmk}
    
    \begin{mythm}
        Suppose $f$ is infinitely differentiable in an open interval $I = (a-r,a+r)$ and there is a constant $A>0$ such that $|f^{(n)}(x)| \le A^n$ for $n\ge 1$ and for any $x\in I$. Then the Taylor series generated by $f$ at $x=a$ converges to $f$ for any $x\in I$.
    \end{mythm}
    \begin{proof}
        Recall that if $f$ has continous derivative of order $n+1$ in some interval containing $a$, then the Taylor formula is the following
        $$
        f(x) = \sum_{k=0}^n \frac{f^{(k)}(a)}{k!} (x-a)^k + E_n(x) \quad \text{where } E_n(x) = \frac1{n!} \int_a^x (x-t)^n f^{(n+1)}(t) \dd t.
        $$
        We do a change of variable $t=x+(a-x)u$. Then $\displaystyle \begin{cases}
            t=a \Leftrightarrow u=1 \\
            t=x \Leftrightarrow u=0 \\
            \dd t = (a-x) \dd u
        \end{cases}
        $
        Thus the error term can be written as
        \[
        E_n(x) = \frac1{n!} (x-a)^{n+1} \int_0^1 u^n f^{(n+1)}(x+(a-x)u) \dd u.
        \]
        Note 
        $$
        \begin{aligned}
            0 \le |E_n(x)| &\le \frac{1}{n!}|x-a|^{n+1} \Big|\int_0^1 u^n f^{(n+1)}(x+(a-x)u) \dd u \Big|\\
            &\le \frac{1}{n!} |x-a|^{n+1}\int_0^1 u^n |f^{(n+1)}(x+(a-x)u)| \dd u\\
            &\le \frac{1}{n!}|x-a|^{n+1} A^{n+1} \int_0^1 u^n \dd u = \frac{(A|x-a|)^{n+1}}{(n+1)!} \to 0 \text{ as } n\to \infty.
        \end{aligned}
        $$
        By the squeeze theorem, we know $E_n(x) \to 0$ as $n\to \infty$ and thus the Taylor series generated by $f$ at $x=a$ converges to $f$ for any $x\in I$.
    \end{proof}
    \begin{myex}
        \leavevmode
        \begin{enumerate}
            \item For $f(x) = e^x$,  $|f^{(n)}(x)| = |e^x|\le A^n$ when $x\in (-r,r)$ and let $A=e^r$. Since $r$ can be any positive real number, actually
            $$
            1+x+\frac{x^2}{2!}+\cdots + \frac{x^n}{n!} + \cdots = e^x \quad \text{ for } x \in \R.
            $$
            \item For $f(x) = \sin x$, $|f^{(n)}(x)| \le 1$ for any $n\ge 1$ and any $x\in \R$. Thus
            $$
            \sin x = x - \frac{x^3}{3!} + \frac{x^5}{5!} - \cdots + (-1)^n \frac{x^{2n+1}}{(2n+1)!} + \cdots \quad \text{ for } x\in \R.
            $$
            \item For $f(x) = \cos x$, $|f^{(n)}(x)| \le 1$ for any $n\ge 1$ and any $x\in \R$. Thus
            $$
            \cos x = 1 - \frac{x^2}{2!} + \frac{x^4}{4!} - \cdots + (-1)^n \frac{x^{2n}}{(2n)!} + \cdots \quad \text{ for } x\in \R.
            $$
        \end{enumerate}
    \end{myex}
    \begin{myex}
        Power series may be used to solve differential equations. For example, consider the following second order ordinary differential equation $y=y(x)$:
        $$ (1-x) y'' = -2y.$$
        We can try to find a power series solution of the form $\displaystyle y = \sum_{n=0}^\infty a_n x^n$. Then we have
        \[
        (1-x) \sum_{n=2}^\infty n(n-1) a_n x^{n-2} = \sum_{n=0}^{\infty}\left[ (n+1)(n+2) a_{n+2} - n(n-1) a_n \right] x^n = \sum_{n=0}^\infty -2 a_n x^n.
        \]
        Thus we have the following recursive relation for the coefficients $a_n$:
        $$ a_{n+2} = \frac{n-2}{n+2} a_n \quad \text{for } n\ge 0. $$
        Comparing the coefficients of $x^n$ on both sides, we have 
        $$ \begin{cases}
            a_2 = -a_0 \\
            a_{2n} =  0 \cdot a_2 = 0 \quad &\text{for } n\ge 2 \\
            a_{2n+1} = \displaystyle \frac{-1}{(2n+1)(2n-1)} a_1 \quad &\text{for } n\ge 1
        \end{cases}
        $$
        Thus we have
        $$ y = a_0 (1-x^2) + a_1 \sum_{n=0}^\infty \frac{-1}{(2n+1)(2n-1)} x^{2n+1} \quad \text{for } x\in (-1,1). $$
        Thus the series gives a family of solutions to the differential equation. However, this is not necessarily all the solutions to the differential equation since we only consider power series solutions. In other words, can we have some solutions not infinitely differentiable but still satisfying the differential equation? We'll introduce several basic theorems about differential equations here.
    \end{myex}

    \newpage
    \section{Vector Algebra and Analytic Geometry}
    \subsection{Lines and Planes}
    \begin{myrmk}
        We shall use the words ``point'' and ``vector''(a vector with initial point at the origin $O=(0,0,\ldots,0)$ and ending point at the point $P$) interchangeably to denote an element $P = (p_1, p_2, \cdots, p_n) \in \R^n$.
    \end{myrmk}
    \begin{mydef}
        Suppose $P$ is a given point and $A$ is a given nonzero vector in $\R^n(n\ge 2)$. The set of all points of the form $P+tA$ where $t\in \R$ is called the \textbf{line} passing through $P$ and parallel to $A$. We denote this line by $L(P,A)$ and write $L(P,A) = \{ P+tA : t\in \R\}$. A point $Q$ is said to be \textbf{on the line} $L(P,A)$ if $Q\in L(P,A)$. We often call $A$ a \textbf{direction vector} of line $L(P,A)$.
    \end{mydef}
    \begin{mythm}
        Two lines $L(P,A)$ and $L(P,B)$ through the same point $P$ are equal sets if and only if the dirction vectors $A$ and $B$ are parallel, i.e. there is some $\lambda \in \R$ such that $A = \lambda B$ for some $\lambda \in \R_{\neq 0}$.
    \end{mythm}
    % \begin{proof}
    %     $(\Rightarrow)$ Suppose $L(P,A) = L(P,B)$, then $P+A \in L(P,B)$ and there is some $\lambda \in \R$ such that $P+A = P+\lambda B$. Thus $A = \lambda B$ for some $\lambda \in \R_{\neq 0}$. 

    %     $(\Leftarrow)$ Suppose $A = \lambda B$ for some $\lambda \in \R_{\neq 0}$. We need to show $L(P,A) \subseteqq L(P,B)$ and $L(P,B) \subseteqq L(P,A)$. Let $\tilde{P}\in L(P,A)$, then $\tilde{P} = P + tA$ for some $t\in \R$. Note $\tilde{P} = P + t \lambda B \in L(P,B)$. Thus $L(P,A) \subseteqq L(P,B)$. Similarly we can show $L(P,B) \subseteqq L(P,A)$ and thus $L(P,A) = L(P,B)$.
    % \end{proof}
    \begin{mythm}
        Two lines $L(P,A)$ and $L(Q,A)$ with the same direction vector $A$ are equal if and only if $Q\in L(P,A)$.
    \end{mythm}
    \begin{mythm}
        Given a line $L_1$ and a point $Q$ not on $L_1$, there is one and only one line $L_2$ containing $Q$ and parallel to $L_1$.
    \end{mythm}
    \begin{mythm}
        If $P\neq Q$ are two points, then there is one and only one line containing both $P$ and $Q$. And this line is simply the set $\{ P+t(Q-P) : t\in \R\}$.
    \end{mythm}
    \begin{mythm}
        We call a function $X: \R \to \R^n$ (with $n\ge 2$) a \textbf{vector-valued function} (of a real variable). If $P$ is a point and $A$ is a nonzero vector in $\R^n$, then $X(t) = P+tA$ is a vector-valued function whose image is the line $L(P,A)$. Here $t$ is often called a parameter and $X$ is called a \textbf{vector/parametric equation} of the line $L(P,A)$.
    \end{mythm}
    \begin{mydef}
        Note that distinct values of the parameter $t$ correspond to distinct points on the line. If $t_1<t_2<t_3$, then we say $X(t_2)$ is \textbf{between} $X(t_1)$ and $X(t_3)$ on the line $L(P,A)$. A pair of points $P,Q$ are called \textbf{congruent} to another pair of points $P',Q'$ if $|P-Q| = |P'-Q'|$ where $|\cdot|$ is the Euclidean distance in $\R^n$. We also call $|P-Q|$ the \textbf{distance between} $P$ and $Q$.
    \end{mydef}
    \begin{myrmk}
        If $P=(p_1,p_2,\cdots,p_n)$ and $A = (a_1,a_2,\cdots,a_n)\ne 0$, then the vector equation $X=(x_1,x_2,\cdots,x_n)$ of the line $L(P,A)$ is equivalent to (a set of) $n$ \textbf{scalar parametric equations} of the line $L(P,A)$:
        \[
        \begin{cases}
            x_1 = p_1 + t a_1 \\
            x_2 = p_2 + t a_2 \\
            \cdots \\
            x_n = p_n + t a_n
        \end{cases}
        \]
        Particularly in $\R^2$, solving the scalar parametric equations gives us
        \[
        a_2 (x_1(t) - p_1) = a_1 (x_2(t) - p_2) \Leftrightarrow N \cdot (X(t) - P) = 0 \quad \text{where } N := (a_2,-a_1).
        \]
        Here $N$ is called a \textbf{normal vector} of the line $L(P,A)$ since it is perpendicular to the direction vector $A$. Furthermore, if $a_1 \ne 0$, we can rewrite the scalar parametric equations as
        \[
        (x_2(t) - p_2) = \frac{a_2}{a_1} (x_1(t) - p_1) 
        \]
        The value $\dfrac{a_2}{a_1}$ is called the \textbf{slope} of the line.
    \end{myrmk}
    \begin{mythm}
        Suppose $L$ is a line in $\R^2$ consisting of all points $x$ satisfying $x \cdot N = P \cdot N$ where $P$ is a given point on $L$ and $N$ is a given nonzero vector for $L$. Let
        $$ d:= \frac{|P\cdot N|}{|N|}.$$
        Then $|x| \ge d$ and $|x| = d$ if and only if $x$ is the projection of $P$ onto $N$, i.e. $x = \dfrac{P\cdot N}{|N|^2} N$.
    \end{mythm}
    \begin{mydef}
        Suppose $P$ is a given point and $A$ and $B$ are two \textbf{linearly independent} vectors in $\R^n(n\ge 3)$. The set $M:=\{ P + sA + tB : s,t\in \R\}$ is called the \textbf{plane} passing through $P$ and parallel to $A$ and $B$. A point $Q$ is said to be \textbf{on the plane} $M$ if $Q\in M$.
    \end{mydef}
    \begin{mythm}
        Two planes $M_1 = \{ P + sA + tB\} $ and $M_2 = \{ P + sC + tD\}$ through the same point $P$ are equal sets if and only if $\operatorname{span} \{A,B\} = \operatorname{span} \{C,D\}$.
    \end{mythm}
    \begin{mythm}
        Two planes $M_1 = \{ P + sA + tB\}$ and $M_2 = \{ Q + sA + tB\}$ with the same direction vectors $A$ and $B$ are equal if and only if $Q\in M_1$.
    \end{mythm}
    \begin{mydef}
        Two planes $M_1 = \{ P + sA + tB\}$ and $M_2 = \{ Q + sC + tD\}$ are called \textbf{parallel} if $\operatorname{span} \{A,B\} = \operatorname{span} \{C,D\}$ but $Q\notin M_1$. We also say that a vector $X$ is \textbf{parallel} to a plane $M$ if $X\in \operatorname{span} \{A,B\}$.
    \end{mydef}
    \begin{myrmk}
        Similarly, a line is parallel to a plane if a direction vector of the line is parallel to the plane. Note a line parallel to a plane may be on the plane.
    \end{myrmk}
    \begin{mythm}
        Given a plane $M$ and a point $Q$ not on $M$, there is one and only one plane $M'$ containing $Q$ and parallel to $M$.
    \end{mythm}
    \begin{mythm}
        If $P,Q$ and $R$ are three non-collinear points, i.e., not on the same line, then there is one and only one plane containing all three points. And this plane is simply the set $\{ P + s(Q-P) + t(R-P) : s,t\in \R\}$.
    \end{mythm}
    \begin{mydef}
        A function $X: \R^2 \to \R^n$ (with $n\ge 3$) is called a \textbf{vector-valued function} of two real variables. If $P$ is a point and $A$ and $B$ are two linearly independent vectors in $\R^n$, then $X(s,t) = P + sA + tB$ is a vector-valued function whose image is the plane $M = \{ P + sA + tB : s,t\in \R\}$. Here $(s,t)$ is often called parameters and $X$ is called a \textbf{vector/parametric equation} of the plane $M$.
    \end{mydef}
    \begin{myrmk}
        Similarly a plane can also be described by a \textbf{scalar equation} of the form $X\cdot N = P\cdot N$ where $P$ is a given point on the plane and $N$ is a given nonzero vector perpendicular to the plane, i.e. $N$ is a normal vector of the plane. Suppose $N=(a,b,c)$, then $N\cdot (X(s,t)-P) = 0$ is equivalent to
        \[
        d = ap_1 + bp_2 + cp_3, \quad ax_1(s,t) + bx_2(s,t) + cx_3(s,t) = d.
        \]
    \end{myrmk}
    \subsection{Cross Product}
    \begin{mydef}
        Let $A=(a_1,a_2,a_3)$ and $B=(b_1,b_2,b_3) \in \R^3$. Their \textbf{cross product} $A \times B$ is defined to be
        $$ A \times B = (a_2 b_3 - a_3 b_2, a_3 b_1 - a_1 b_3, a_1 b_2 - a_2 b_1). $$
    \end{mydef}
    \begin{mythm}
        For $A,B,C\in \R^3$ and $k \in \R$, we have
        \begin{enumerate}[label=(\arabic*)]
            \item $A \times B = - B \times A$;
            \item $A \times (B+C) = A \times B + A \times C$ and $(B+C)\times A = B\times A + C\times A$;
            \item $(kA) \times B = k (A\times B)$ and $A\times (kB) = k (A\times B)$;
            \item $A\cdot (A\times B) = 0, B \cdot (A\times B) = 0$;
            \item $\| A \times B\| ^2 = \|A\|^2 \|B\|^2 - (A\cdot B)^2$;
            \item $A\times B = 0$ if and only if $A$ and $B$ are linearly dependent;
            \item $\|A\times B\| = \|A\|\|B\|\sin\theta$ where $\theta$ is the angle between $A$ and $B$.
            
        \end{enumerate}
    \end{mythm}
    \begin{mythm}
        Suppose $A$ and $B$ are two linearly independent vectors in $\R^3$. Then
        \begin{enumerate}
            \item $A$, $B$ and $A\times B$ are linearly independent;
            \item Any $N\in \R^3$ such that $N\cdot A = 0$ and $N\cdot B = 0$ must be parallel to $A\times B$, i.e. $ N \in \operatorname{span} \{A\times B\}$.
        \end{enumerate}
    \end{mythm}
    \begin{myrmk}
        The direction of $A\times B$ is determined by the right-hand rule, which corresponds to the common 3D-coordinate system. 
    \end{myrmk}
    \begin{mydef}
        Let $A,B,C\in \R^3$. $A\cdot B \times C := A\cdot (B\times C)$ is called the \textbf{scalar triple product} of $A,B$ and $C$.
    \end{mydef}
    \begin{mythm}
        Suppose $A,B,C\in \R^3$. Then
        \begin{enumerate}[label=(\arabic*)]
            \item $A,B$ and $C$ are linearly dependent if and only if $A\cdot B \times C = 0$;
            \item $A \cdot B \times C = B \cdot C \times A = C \cdot A \times B$.
            \item $A \times B \cdot C = A \cdot B \times C$.
        \end{enumerate}
    \end{mythm}
    \begin{mydef}
        Let $M = \{ P + sA + tB : s,t\in \R\} \in \R^3$ be the plane through $P$ and spanned by linearly independent vectors $A$ and $B$. A vector $N \in \R^3$ is said to be perpendicular to the plane $M$ if $N$ is perpendicular to both $A$ and $B$, i.e. $N\cdot A = 0$ and $N\cdot B = 0$. If, additionally, $N\ne 0$, then $N$ is called a \textbf{normal vector} of the plane $M$.
    \end{mydef}
    \begin{mythm}
        Given a plane $M = \{ P + sA + tB : s,t\in \R\}$ in $\R^3$ where $P$ is a point and $A$ and $B$ are linearly independent vectors. Let $N=A\times B$, then $N \ne 0 $ and $N$ is a normal vector of the plane $M$. Let $M' := \{ X\in \R^3 : N\cdot (X-P) = 0\}$, then $M' = M$.
    \end{mythm}
    \begin{myrmk}
        If $N=(a,b,c) \ne 0$, $P=(x_1,x_2,x_3)$ and $X=(x,y,z)$, then the scalar equation $N\cdot (X-P) = 0$ of the plane $M$ can be rewritten as
        \[
        ax + by + cz = ax_1 + bx_2 + cx_3 \quad \text{or} \quad a(x-x_1) + b(y-y_2) + c(z-z_3) = 0.
        \]
    \end{myrmk}
    \begin{mythm}
        Given a plane $M$ through a point $P$ and with a normal vector $N$, let $d:= \dfrac{\|P\cdot N\|}{\|N\|}$. Then for any $X\in M$, we have $|X| \ge d$ and $|X| = d$ if and only if $X$ is the projection of $P$ onto $N$, i.e. $X = \dfrac{P\cdot N}{\|N\|^2} N$.
    \end{mythm}
    \begin{myrmk}
        Suppose point $Q$ is not on $M$, then for any $X\in M$, the minimum value of $\| X-Q\|$ is $\dfrac{\| (Q-P)\cdot N \|}{\|N\|}$ and is achieved when $(X-Q)$ is the projection of $(Q-P)$ onto $N$. We call $\dfrac{\| (Q-P)\cdot N \|}{\|N\|}$ the \textbf{distance} from point $Q$ to the plane $M$.
    \end{myrmk}

    \newpage
    \section{Multivariable Calculus}
    \subsection{Vector-valued functions}
    \begin{mydef}
        Suppose $f_1,\cdots,f_n$ are functions fron $\R$ to $\R$. We call $$F(t) = (f_1(t), f_2(t), \cdots, f_n(t))$$ a \textbf{vector-valued function} and we define the \textbf{limits,derivatives} and \textbf{integrals} by the following:
        $$ \begin{aligned}
            \lim_{t\to p} F(t) &= \left( \lim_{t\to p} f_1(t), \cdots, \lim_{t\to p} f_n(t) \right),\\
            F'(t) &= (f_1'(t), \cdots, f_n'(t)),\\
            \int_a^b F(t) \dd t &= \left( \int_a^b f_1(t) \dd t, \cdots, \int_a^b f_n(t) \dd t \right).
        \end{aligned} $$
    \end{mydef}
    \begin{mythm}
        Suppose $F$ and $G$ are vector-valued functions from $\R$ to $\R^n$ and $u$ is a real-valued function from $\R$ to $\R$. If $F,G$ and $u$ are differentiable on in interval, then so are $F+G$, $uF$ and $F\cdot G$ and we have
        \begin{enumerate}[label=(\arabic*)]
            \item $(F+G)' = F' + G'$;
            \item $(uF)' = u' F + u F'$;
            \item $(F\cdot G)' = F'\cdot G + F\cdot G'$.
        \end{enumerate}
        For the case $n=3$, then $F\times G$ is also differentiable with $(F\times G)' = F'\times G + F \times G'$.
    \end{mythm}
    \begin{mypro}
        Suppose $F$ is a vector-valued function from $\R$ to $\R^n$. If $F$ is differentiable and be constant length/norm on an open interval $I$, i.e., $\| F(t)\| = c$ for any $t\in I$, then $F \cdot F' = 0$ on $I$. Namely, $F'(t)$ is always perpendicular to $F(t)$ for any $t\in I$.
    \end{mypro}
    \begin{proof}
        Let $g(t) := F(t)\cdot F(t) = \| F(t)\|^2 = c^2$ for any $t\in I$. Then $g'(t) = 2 F(t)\cdot F'(t) = 0$ for any $t\in I$. Thus $F(t)\cdot F'(t) = 0$ for any $t\in I$.
    \end{proof}
    \begin{mythm}
        Suppose $F:\R \to \R^n$ and $u:\R \to \R$. Let $G:=F \circ u$. Then:
        \begin{enumerate}[label=(\arabic*)]
            \item If $u$ is continuous at $t$ and $F$ is continuous at $u(t)$, then $G$ is continuous at $t$.
            \item If the derivative $u'(t)$ and $F'(u(t))$ exist, then $G'(t)$ exists and $G'(t) = u'(t) F'(u(t))$.
        \end{enumerate}
    \end{mythm}
    \begin{mythm}
        Suppose $F:\R \to \R^n$ is continuous on $[a,b]$, then:
        \begin{enumerate}[label=(\arabic*)]
            \item $F$ is integrable on $[a,b]$;
            \item For any $c\in [a,b]$, we can define the primitive or indefinite integral function $A(x) := \displaystyle \int_c^x F(t) \dd t$ for any $x\in [a,b]$. Then $A$ is differentiable on $(a,b)$ and $A'(x) = F(x)$ for any $x\in (a,b)$.
            \item If the derivative $F'(t)$ is continuous on some open interval $I$, then for every $c$ and $x$ in $I$, we have
            $$ F(x) = F(c) + \int_c^x F'(t) \dd t. $$
        \end{enumerate}
    \end{mythm}
    \begin{mythm}
        Suppose $F$ and $G$ are from $\R$ to $\R^n$ and are integrable on $[a,b]$. Then :
        \begin{enumerate}[label=(\arabic*)]
            \item for any $\alpha , \beta \in \R$, $\alpha F + \beta G$ is also integrable on $[a,b]$ with
            $$ \int_a^b (\alpha F + \beta G)(t) \dd t = \alpha \int_a^b F(t) \dd t + \beta \int_a^b G(t) \dd t; $$
            \item Let $C \in \R^n$, then the dot product $C \cdot F$ is also integrable on $[a,b]$ with
            $$ \int_a^b (C\cdot F)(t) \dd t = C \cdot \int_a^b F(t) \dd t. $$
            \item If $\| F \|$ is also integrable on $[a,b]$, then 
            $$ \left\| \int_a^b F(t) \dd t \right\| \le \int_a^b \| F(t)\| \dd t. $$
        \end{enumerate}
    \end{mythm}
    \begin{proof}
        We only prove (3) here. Let $A:= \displaystyle \int_a^b F(t) \dd t$. The case where $\left\|A\right\| = 0 $ is trivial. Suppose $\left\|A\right\| > 0$. Then
        $$ \| A\|^2 = A \cdot A = \int_a^b F(t) \cdot A \dd t \le \int_a^b \| F(t)\| \| A\| \dd t = \| A\| \int_a^b \| F(t)\| \dd t. $$
        Thus $\displaystyle \left\|\int_a^b F(t) \dd t\right\| = \| A\| \le \int_a^b \| F(t)\| \dd t$.
    \end{proof}
    \subsection{Curves}
    \begin{mydef}
        Suppose $X$ is a vector-valued function from $\R$ to $\R^n$ with $n=2$ or $n=3$. Suppose the domain of $X$ is an interval $D\subseteq \R$. If $X$ is continuous on $D$, then we call the range/image of $X$ a \textbf{curve} in $\R^n$. We also say that the curve is described (parametrically) by the function $X$ with parameter $t\in D$. We also call $D$ a parameter interval.
    \end{mydef}
    \begin{myrmk}
        The curve described by $X$ is invarient under the change of parameter. For example, if $u:\R \to \R$ and $I = \{u(t) :t \in J\}$ for some interval $J$, then $Y := X\circ u$ also describes the same curve as $X$. We say $X$ and $Y$ are equivalent in describing the curve $C$.
    \end{myrmk}
    \begin{mydef}
        Let $C$ be a curve described by $X$ (note $X$ is continuous). If the derivative $X'(t)$ exists and $X'(t) \ne 0$, then the line through point $X(t)$ with direction vector $X'(t)$ is called the \textbf{tangent line} of the curve $C$ at point $X(t)$.
    \end{mydef}
    \begin{myrmk}
        Though the derivative of $X$ and $Y=X\circ u$ are different, the tangent lines of the curve $C$ at point $X(t) = Y(s)$ are the same since $Y'(s) = u'(s) X'(u(s))$ is parallel to $X'(u(s))$. Thus the tangent line of a curve at a point is well-defined and does not depend on the choice of parameter.
    \end{myrmk}
    \begin{mydef}
        Consider a motion (of a particle) described by a vector-valued function $X: \R \to \R^n$ with $n=2$ or $n=3$. The derivative $X'(t)$ is called the \textbf{velocity} at time $t$. $\| X'(t)\|$ is called the \textbf{speed} at time $t$. The second derivative $X''(t)$ is called the \textbf{acceleration} at time $t$. In some textbooks, we use $\vec{r}(t)$ for $X(t)$, $\vec{v}(t)$ for $X'(t)$, $v(t)$ for $\| X'(t)\|$ and $\vec{a}(t)$ for $X''(t)$.
    \end{mydef}
    \begin{myex}
        Suppose $\vec{r}(t) = (r \cos(f(t)), r \sin(f(t)), bf(t))$ with $r>0$ and $b\ne 0$ fixed. This is called a \textbf{helix/helical curve} in $\R^3$. Note a helix is not contained in any plane in $\R^3$ but it is contained fully in a cylinder/cylindrical surface in $\R^3$.
    \end{myex}
    \begin{mydef}
        Consider a curve described by $X(t)$. We define the \textbf{unit tangent vector} $T$ as $$T(t) := \dfrac{X'(t)}{\| X'(t)\|} \quad \text{whenever } X'(t) \ne 0. $$ Note that $\| T(t) \| = 1$ and thus $T(t) \cdot T'(t) = 0$. We define the \textbf{principal normal vetor} $N$ as
        $$ N(t) := \frac{T'(t)}{\| T'(t)\|} \quad \text{whenever } T'(t) \ne 0. $$
        Note $\left\|N(t)\right\| =1$ and $N(t) \cdot T(t) = 0$. Furthermore, we call the plane through the point $X(t)$ and spanned by $T(t)$ and $N(t)$ the \textbf{osculating plane} of the curve at point $X(t)$. 
    \end{mydef}
    \begin{myrmk}
        Consider three points $X(t_1), X(t_2)$ and $X(t_3)$ on the curve. As $t_2 $ and $t_3$ approach $t_1$, the plane determined by $X(t_1), X(t_2)$ and $X(t_3)$, supposing the three points are not collinear, approaches the osculating plane at point $X(t_1)$. Note if the curve is a planar curve, then the osculating plane is the fixed plane containing the curve.
    \end{myrmk}
    \begin{mythm}
        For a motion described by $\vec r(t)$, the acceleration $\vec a(t) := \vec r\,''(t)$ is a linear combination of $T(t)$ and $T'(t)$ given by $$ \vec a(t) = v'(t) T(t) + v(t) T'(t). $$
        If, additionally $T'(t) \ne 0$, then
        $$ \vec a(t) = v'(t) T(t) + v(t) \| T'(t)\| N(t). $$
    \end{mythm}
    \begin{myrmk}
        We call $v'(t)$ and $v(t) \| T'(t)\|$ the tangential and normal components of the acceleration $\vec a(t)$, respectively. In some textbook, the \textbf{binormal vector} $B$ is also introduced as $B(T) := T(t) \times N(t)$. Note $\left\| B(t) \right\|=1$ and $T(t), N(t)$ and $B(t)$ together provide a moving reference frame for $\R^3$ along the curve.
    \end{myrmk}
    \begin{mydef}
        For a curve in $\R^2$ and $\R^3$ described by $\vec r(t)$ for $t\in [a,b]$ (Note $\vec r$ is required to be continuous on $[a,b]$). Let $P=\{t_0, t_1, \ldots, t_n\}$ be a partition of $[a,b]$. We denote the polygon whose vertices are $\vec r(t_0), \vec r(t_1), \ldots, \vec r(t_n)$ by $\Pi (P)$ and denote the \textbf{length of polygon} by $\| \Pi (P)\|$. Then
        $$ \| \Pi (P)\| = \sum_{i=1}^n \| \vec r(t_i) - \vec r(t_{i-1})\|. $$
        If there exists an $M>0$ such that $\| \Pi (P)\| \le M$ for any partition $P$ of $[a,b]$, then we say the curve is said to be \textbf{rectifiable} and its arc length, denoted by $\Lambda(a,b)$, is defined as
        $$ \Lambda(a,b) := \sup \{ \| \Pi (P)\| : P \text{ is a partition of } [a,b]\}. $$
        If there's no such $M$, then the curve is said to be \textbf{non-rectifiable}.
    \end{mydef}
    \begin{mythm}
        For a curve in $\R^2$ or $\R^3$ described by $\vec r(t)$ for $t\in [a,b]$. If $\vec v(t) := \vec r\,'(t)$ is continuous on $[a,b]$, then the curve is rectifiable and its arc length satisfies
        $$ \Lambda(a,b) \le \int_a^b \| \vec v(t)\| \dd t. $$

    \end{mythm}
    \begin{proof}
        Suppose $P = \{t_0, t_1, \ldots, t_n\}$ is a partition of $[a,b]$. Then
        $$ \| \Pi (P)\| = \sum_{i=1}^n \| \vec r(t_i) - \vec r(t_{i-1})\| = \sum_{i=1}^n \left\| \int_{t_{i-1}}^{t_i} \vec v(t) \dd t \right\| \le \sum_{i=1}^n \int_{t_{i-1}}^{t_i} \| \vec v(t)\| \dd t = \int_a^b \| \vec v(t)\| \dd t. $$
        Since $\Lambda (a,b) $ is the least upper bound, we have $\Lambda(a,b)\displaystyle \le \int_a^b \| \vec v(t)\| \dd t$.
    \end{proof}
    \begin{mythm}
        For a rectifiable curve in $\R^2$ or $\R^3$ described by $\vec r(t)$ for $t\in [a,b]$. If $a<c<b$ and $C_1$ and $C_2$ are the curves corresponding to $\vec r(t)$ on $[a,c]$ and $[c,b]$, respectively, then $C_1$ and $C_2$ are both rectifiable and $\Lambda(a,b) = \Lambda(a,c) + \Lambda(c,b)$.
    \end{mythm}
    \begin{proof}
        Suppose $P_1$ and $P_2$ are partitions of $[a,c]$ and $[c,b]$, respectively. Let $P = P_1 \cup P_2$ be the joint partition. Then $ \| \Pi (P_1)\| + \| \Pi (P_2)\| = \| \Pi (P)\| \le \Lambda(a,b)$. Thus $\| \Pi (P_1)\|$ and $\| \Pi (P_2)\|$ are both bounded above by $\Lambda(a,b)$ for any partitions $P_1$ and $P_2$. Thus $C_1$ and $C_2$ are both rectifiable.
        \begin{itemize}
            \item For a fixed $P_2$, $\| \Pi (P_1)\|  \le \Lambda(a,b) - \| \Pi (P_2)\|$ holds for any partition $P_1$ of $[a,c]$ and there by $$\Lambda(a,c) \le \Lambda(a,b) - \| \Pi (P_2)\| \implies \| \Pi (P_2) \| \le \Lambda(a,b) - \Lambda(a,c).$$
            \item since the above inequality holds for any partition $P_2$ of $[c,b]$, we have $\Lambda(a,b) \ge \Lambda(a,c) + \Lambda(c,b)$.
        \end{itemize}
        We still need to show $\Lambda(a,b) \le \Lambda(a,c) + \Lambda(c,b)$. Suppose $P$ is a partition of $[a,b]$. Adding in $c$ into $P$, we obtain a partition $P_1$ of $[a,c]$ and a partition $P_2$ of $[c,b]$ such that
        $$ \| \Pi (P)\| \le \| \Pi (P_1)\| + \| \Pi (P_2)\| \le \Lambda(a,c) + \Lambda(c,b). $$
        Thus $\Lambda(a,b) \le \Lambda(a,c) + \Lambda(c,b)$. Combining the two inequalities, we have $\Lambda(a,b) = \Lambda(a,c) + \Lambda(c,b)$.
    \end{proof}
    \begin{mydef}
        Suppose a curve in $\R^2$ or $\R^3$ is described by $\vec r(t)$ for $t\in [a,b]$. We define the \textbf{arc-length function} $s$ as
        $$ s(t) := \begin{cases}
            0, & t=a; \\
            \Lambda(a,t), & t\in (a,b].
        \end{cases}
        $$
    \end{mydef}
    \begin{mythm}
        For any rectifiable curve, the arc-length function $s$ is monotonically increasing on $[a,b]$.
    \end{mythm}
    \begin{mythm}
        Suppose a rectifiable curve described by $\vec r(t)$ for $t\in [a,b]$. Let $s$ denote the arc-length function. If the speed $v(t) = \| \vec r\,'(t)\|$ is continuous on $[a,b]$, then the derivative $s'(t)$ exists and is given by $s'(t) = v(t)$ for any $t\in [a,b]$.
    \end{mythm}
    \begin{proof}
        Suppose $h>0$, then
        $$ \frac{s(t+h) - s(t)}{h} \le \frac1h \int_t^{t+h} v(u) \dd u \to v(t) \quad \text{as } h\to 0^+. $$
        We also have
        $$ \frac{s(t+h) - s(t)}{h} \ge \left\|\frac{\vec r(t+h) - \vec r(t)}{h}\right\| \to v(t) \quad \text{as } h\to 0^+. $$
        By the squeeze theorem, we know $\displaystyle \lim_{h\to 0^+} \frac{s(t+h) - s(t)}{h} = v(t)$ and thus $s'(t) = v(t)$. The case where $h<0$ can be proved similarly.
    \end{proof}
    \begin{myrmk}
        Note
        $$ s(t_2) - s(t_1) = \int_{t_1}^{t_2} v(t) \dd t,$$
        provided that $v(t)$ is continuous on $[a,b]$. Particularly, $$\Lambda(a,b) = s(b) - s(a) = \int_a^b v(t) \dd t.$$
    \end{myrmk}
    \begin{myex}
        Consider the curve in $\R^2$ given by $\vec r(t) = (t, f(t))$ for $t\in [a,b]$ where $f$ is a scalar-valued function differentiable on $[a,b]$. Then the arc-length function is given by
        $$ s(t) = \int_a^t v(u) \dd u = \int_a^t \sqrt{1 + (f'(u))^2} \dd u. $$
    \end{myex}
    \begin{mydef}
        Let $T$ denote the unit tangent vector and $s$ denote the arc-length function of a rectifiable curve described by $\vec r(t)$ for $t\in [a,b]$. We call $\dfrac{\dd T}{\dd s}$ the \textbf{curvature vector} of the curve provided it's well-defined. Note
        $$ \frac{\dd T}{\dd s} = \frac{\dd T}{\dd t} \cdot \frac{\dd t}{\dd s} = \frac{T'(t)}{v(t)} = \frac{\left\|T'(t)\right\| N(t)}{v(t)}. $$
        This means that the curvature vector is parallel to the principal normal vector $N$. We call $$\kappa(t) := \left\| \dfrac{\dd T}{\dd s}\right\| = \dfrac{\left\|T'(t)\right\|}{v(t)}$$ the \textbf{curvature number} of the curve at point $\vec r(t)$.
    \end{mydef}
    \begin{mydef}
        When $\kappa(t) \ne 0$, we call $\rho(t) := \dfrac{1}{\kappa(t)}$ the \textbf{radius of curvature}. Then then circle in the osculating plane with radius $\rho(t)$ and centered at $\vec r(t) + \dfrac{1}{\kappa(t)} N(t)$ is called the \textbf{osculating circle}. Note the osculating circle is tangent to the curve at point $\vec r(t)$ and has the same curvature as the curve at point $\vec r(t)$.
    \end{mydef}

    \begin{mythm}
        For motion along a curve in $\R^2$ or $\R^3$ with velocity $\vec v(t)$, speed $v(t)$, acceleration $\vec a(t)$ and curvature $\kappa$. Then we have
        $$ \vec a(t) = v'(t) T(t) + \kappa (t) v^2(t) N(t). $$
        and
        $$ \kappa (t) = \frac{\left\|\vec a (t) \times \vec v (t)\right\|}{v^3 (t)}
        $$
    \end{mythm}
    \subsection{Curves in Polar Coordinates}
    
    \newpage
    \section{Differential Calculus of Scalar and Vector Fields}
    \subsection{Basic Concepts}
    \begin{mydef}
        A function $f: \R^n \to \R^m$ is called
        \begin{itemize}
            \item a \textbf{real-valued function} $f(t)$ if $n=m=1$;
            \item a \textbf{vector-valued function} (of a real variable) $\vec f (t)$ if $n=1$ and $m\ge 2$;
            \item a \textbf{scalar field} $f(\vec x)$ if $n\ge 2$ and $m=1$;
            \item a \textbf{vector field} $\vec f (\vec x)$ if $n\ge 2$ and $m\ge 2$. We often write it as $\vec f = (f_1(\vec x), \cdots, f_m(\vec x))$ where $f_i(\vec x)$ denotes the $i$-th component of $\vec f$.
        \end{itemize}
    \end{mydef}
    \begin{mydef}
        Let $\vec a \in \R^n$ and $r>0$. The set of all points $\vec x\in \R^n$ such that $\| \vec x - \vec a\| < r$ is called the \textbf{open ($n$-)ball} of radius $r$ centered at $\vec a$ and is denoted by $B(\vec a; r)$ or simply $B(\vec a)$.
    \end{mydef}
    \begin{mydef}
        Let $S\subseteq \R^n$ and $\vec a \in S$.
        \begin{itemize}
            \item We call $\vec a $ an \textbf{interior point} of $S$ if there exists an $r>0$ such that $B(\vec a; r) \subseteq S$.
            \item The set of all interior points of $S$ is called the \textbf{interior} of $S$ and is denoted by $\operatorname{int}(S)$.
            \item The set $S$ is called \textbf{open} if every point of $S$ is an interior point, i.e. $\operatorname{int}(S) = S$.
            \item A point $\vec b\in \R^n$ is said to be \textbf{exterior} to $S$ if there is an open $n$-ball $B(\vec b)$ such that $B(\vec b) \cap S = \varnothing$.
            \item The set of all points in $\R^n$ exterior to $S$ is called the \textbf{exterior} of $S$ and is denoted by $\operatorname{ext}(S)$.
            \item A point $\vec c\in \R^n$ is called a \textbf{boundary point} of $S$ if $c$ is $\vec c \notin \operatorname{int}(S)$ and $\vec c \notin \operatorname{ext}(S)$, i.e. for any open $n$-ball $B(\vec c)$, we have $B(\vec c) \nsubseteq S$ and $B(\vec c) \cap S \ne \varnothing$.
            \item The set of all boundary points of $S$ is called the \textbf{boundary} of $S$ and is denoted by $\partial S$.
            \item If $\partial S \subseteq S$, then $S$ is called a \textbf{closed} set.
        \end{itemize}
    \end{mydef}
    \begin{myex}
        \begin{enumerate}[label=(\arabic*),start=0]
            \item Is $\R^n$ open or closed? Is $\varnothing$ open or closed?
            \item An open $n$-ball is open. Note
            $$ \operatorname{ext}(B(\vec a; r)) = \{ \vec x\in \R^n : \| \vec x - \vec a\| > r\}, \quad \partial B(\vec a; r) = \{ \vec x\in \R^n : \| \vec x - \vec a\| = r\}. $$
            \item Let $A_1 = (a,b)$ and $A_2 = (c,d)$ where $a<b$, $c<d$ and $A_1, A_2 \subseteq \R$. Then $A_1\times A_2 = \{ (x,y) : x\in A_1, y\in A_2\}$ is an open set in $\R^2$.
        \end{enumerate}
    \end{myex}
    \begin{mydef}
        Consider $S \subseteq \R^n$ and function $\vec f:S\to \R^m$. 
        \begin{itemize}
            \item If $\vec a\in \R^n$ and $\vec b \in \R^m$, we write $$ \lim_{\vec x \to \vec a} \vec f(\vec x) = \vec b \quad \text{or} \quad \vec f(\vec x) \to \vec b \text{ as } \vec x \to \vec a $$ to mean that $$ \lim_{\left\|\vec x - \vec a\right\| \to 0} \left\|\vec f(\vec x) - \vec b\right\| = 0. $$
            \item $\vec f$ is said to be \textbf{continuous} at $\vec a$ if $f$ is defined at $\vec a$ and $\displaystyle \lim_{\vec x \to \vec a} \vec f(\vec x) = \vec f(\vec a)$.
            \item $\vec f$ is said to be \textbf{continuous} on $S$ if $\vec f$ is continuous at every point of $S$.
        \end{itemize}
    \end{mydef}
    \begin{mythm}
        If $\displaystyle \lim_{\vec x \to \vec a} \vec f(\vec x) = \vec b$ and $\displaystyle \lim_{\vec x \to \vec a} \vec g(\vec x) = \vec c$, then
        \begin{itemize}
            \item For $\alpha, \beta \in \R$, $\displaystyle \lim_{\vec x \to \vec a} (\alpha \vec f + \beta \vec g)(\vec x) = \alpha \vec b + \beta \vec c$;
            \item $\displaystyle \lim_{\vec x \to \vec a} (\vec f \cdot \vec g)(\vec x) = \vec b \cdot \vec c$;
            \item $\displaystyle \lim_{\vec x \to \vec a} \left\|\vec f (\vec x) \right\| = \left\|\vec b\right\|$.
        \end{itemize}
    \end{mythm}
    \begin{mythm}
        Suppose $\vec f$ and $\vec g$ are functions such that the composition $\vec f \circ \vec g$ is well-defined at $\vec a$. If $\vec g$ is continuous at $\vec a$ and $\vec f$ is continuous at $\vec g(\vec a)$, then $\vec f \circ \vec g$ is continuous at $\vec a$.
    \end{mythm}
    \begin{myex}
        \begin{enumerate}
            \item Suppose $\vec f=(f_1, \ldots, f_m):\R^n \to \R^m$ is a function. Then $\vec f$ is continuous at $\vec a$ if and only if each component function $f_i$ is continuous at $\vec a$ for $i=1, \ldots, m$.
            \item $\vec f(\vec x) = \vec x$ is called the \textbf{identity function}.
            \item A linear map $f:\R^n \to \R^m$. Then
            $$ \lim_{\left\|h\right\| \to 0} \frac{\| f(\vec a + h) - f(\vec a) - f(h)\|}{\| h\|} = \lim_{\left\|h\right\| \to 0} \frac{\| f(h)\|}{\| h\|}.$$
            If $\vec h \ne \vec 0$,
            \item Polynomials in $n$ variable
            $$ P(\vec x) = \sum_{i_1, \ldots, i_n} c_{i_1, \ldots, i_n} x_1^{i_1} \cdots x_n^{i_n} $$
            are continuous for all $\vec x$.
            \item Scalar fields of the form $f(\vec x) = \dfrac{P(\vec x)}{Q(\vec x)}$ where $P$ and $Q$ are polynomials and $Q(\vec x) \ne 0$ are continuous at any $\vec a$ with $Q(\vec a) \ne 0$.
            \item The function
            $$ f(x,y) = \begin{cases}
                0, & (x,y) = (0,0); \\[6pt]
                \dfrac{xy}{x^2 + y^2}, & (x,y) \ne (0,0)
            \end{cases}
            $$
            is not continuous at $(0,0)$ since $\displaystyle \lim_{x\to 0} f(x,x) = \dfrac{1}{2}$ but $f(0,0) = 0$. This indicates that continuity needs to be checked along all possible paths approaching the point, not just along some specific paths.
        \end{enumerate}
    \end{myex}
    \subsection{The Derivative of a Scalar Field}
    \begin{mydef}
        Given a scalar field $f:S\to \R$ where $S\subseteq \R^n$. Let $\vec a \in \operatorname{int} S$ and $\vec y \in \R^n$.
        \begin{itemize}
            \item The \textbf{derivative of $f$ at $\vec a$ with respect to $\vec y$} is denoted by $f'(\vec a; \vec y)$ and is defined as
            $$ f'(\vec a; \vec y) := \lim_{h\to 0} \frac{f(\vec a + h \vec y) - f(\vec a)}{h}, $$
            provided the limit exists.
            \item In particular, if $\vec y$ is a unit vector, we also call $f'(\vec a; \vec y)$ a \textbf{directional derivative} of $f$ at $\vec a$ in the direction $\vec y$.
            \item In particular, if $\vec y = \vec e_k$ is the $k$-th standard basis vector, we call $f'(\vec a; \vec e_k)$ the \textbf{partial derivative} of $f$ at $\vec a$ with respect to $x_k$ and denote it by $\dfrac{\partial f}{\partial x_k}(\vec a)$ or $D_k f(\vec a)$.
        \end{itemize}
    \end{mydef}
    \begin{myrmk}
        The partial derivatives themselves are also scalar fields, so we can also calculate the partial derivatives of the partial derivatives, which are called the \textbf{second-order partial derivatives}.
    \end{myrmk}
    \begin{mynot}
        In Leibniz's notation, we adopt the following convention for the second-order partial derivatives:
        $$ \frac{\partial^2 f}{\partial x ^2} := \pdv{x}\left( \pdv{f}{x} \right), \quad \frac{\partial^2 f}{\partial x \partial y} := \pdv{x}\left( \pdv{f}{y} \right), \quad \frac{\partial^2 f}{\partial y \partial x} := \pdv{y}\left( \pdv{f}{x} \right).$$
    \end{mynot}
    \begin{myex}
        The function
        $$ f(x,y) = \begin{cases}
            0, & (x,y) = (0,0); \\[6pt]
            \dfrac{x^2 y}{x^4 + y^2}, & (x,y) \ne (0,0)
        \end{cases}
        $$
        is not continuous at $(0,0)$ since approaching $(0,0)$ along the path $y=x^2$ gives
        $$ \lim_{x\to 0} f(x,x^2) = \lim_{x\to 0} \frac{x^4}{2x^4} = \frac{1}{2} \ne f(0,0) = 0. $$
        However, the derivative $f'(\vec 0; \vec y)$ exists for any $\vec y = (y_1, y_2)$ since
        $$ f'(\vec 0; \vec y) = \lim_{h\to 0} \frac{f(h y_1, h y_2) - f(0,0)}{h} = \lim_{h\to 0} \frac{h^3 y_1^2 y_2}{h^5 y_1^4 + h^4 y_2^2} = \lim_{h\to 0} \frac{y_1^2 y_2}{h^2 y_1^4 + h y_2^2} = 0. $$
        This indicates that the existence of directional derivatives does not necessarily imply the continuity of the function. However, if we still want differentiability to imply continuity like in the case of $f:\R \to \R$, we need more stronger definition of differentiability.
            
    \end{myex}
    \begin{mythm}
        For a fixed $\vec a, \vec y \in \R^n$, let $g(t):= f(\vec a + t \vec y)$ for $t\in \R$. Then $g'(t)$ exists if and only if $f'(\vec a + t \vec y; \vec y)$ exists, in which case $g'(t) = f'(\vec a + t \vec y; \vec y)$. Particularly, $f'(\vec a; \vec y) = g'(0)$.
    \end{mythm}
    \begin{mythm}
        Suppose $f'(\vec a + t\vec y; \vec y)$ exists for any $t\in [0,1]$. Then $f(\vec a + \vec y) - f(\vec a) = f'(\vec a + \theta \vec y; \vec y)$ for some $\theta \in (0,1)$.
    \end{mythm}
    \begin{mydef}
        Consider a scalar field $f:S\to \R$ where $S\subseteq \R^n$ and $\vec a \in \operatorname{int} S$. We say $f$ is \textbf{differentiable} at $\vec a$ if there exists a linear map $T_{\vec a}:\R^n \to \R$ and a function $E:\R^n \times \R^n \to \R$ such that
        $$ f(\vec a + \vec v) = f(\vec a) + T_{\vec a}(\vec v) + \|v\|E(\vec a, \vec v) $$
        for any $\vec v$ with $\vec a + \vec v \in B(\vec a;r)\subseteq S$ with $r>0$, where $E(\vec a, \vec v) \to 0$ as $\vec v \to \vec 0$.

        The linear map $T_{\vec a}$ is called the \textbf{total derivative} of $f$ at $\vec a$. The above equation is called the \textbf{first-order Taylor expansion} of $f(\vec a+\vec v)$, whose error term is $\| v\|E(\vec a, \vec v)=o(\|\vec v\|)$.
    \end{mydef}
    \begin{myrmk}
        Note we cannot define derivative through quotient like in the case of $f:\R \to \R$ since we do not have a notion of division in $\R^n$. Instead, inspired by Taylor's expansion, we define the derivative through linear approximation, which is more general and can be applied to functions between general normed vector spaces.
    \end{myrmk}
    \begin{mythm}
        Suppose $f$ is differentiable at $\vec a$ with total derivative $T_{\vec a}$. Then $f'(\vec a; \vec y)$ exists for any $\vec y\in \R^n$ and $T_{\vec a}(\vec y) = f'(\vec a; \vec y)$. Furthermore, 
        $$ f'(\vec a; \vec y) = \sum_{k=1}^n y_k \frac{\partial f}{\partial x_k}(\vec a). $$
    \end{mythm}
    \begin{mynot}
        The vector $\displaystyle\left(\frac{\partial f}{\partial x_1}(\vec a), \ldots, \frac{\partial f}{\partial x_n}(\vec a)\right)$ is called the \textbf{gradient} of $f$ at $\vec a$ and is denoted by $\nabla f(\vec a)$. Then we can write $$f'(\vec a; \vec y) = \nabla f(\vec a) \cdot \vec y.$$
        With this notation, we can also write the first-order Taylor expansion as
        $$ f(\vec a + \vec v) = f(\vec a) + \nabla f(\vec a) \cdot \vec v + o(\|\vec v\|). $$
        This is quite similar to the first-order Taylor expansion for $f:\R \to \R$ given by
        $$ f(a + h) = f(a) + f'(a) h + o(h). $$
    \end{mynot}
    \begin{mythm}
        If a scalar field is differentiable at $\vec a$, then it is continuous at $\vec a$.
    \end{mythm}
    \begin{mythm}
        If all the partial derivatives $D_1 f, \ldots, D_n f$ exist in some open $n-$ball $B(\vec a)$ and continuous at $\vec a$, then $f$ is differentiable at $\vec a$, in which case we also say $f$ is \textbf{continuously differentiable} at $\vec a$.
    \end{mythm}
    \begin{mythm}
        Suppose $f$ is a scalar field on an open set $S\subseteq \R^n$ and $\vec r$ is a vector-valued function mapping $J\subseteq \R$ to $S$. Define the composite function $g = f \circ \vec r$ on $J$ as $g(t) = f(\vec r(t))$ for every $t\in J$. If $f$ is differentiable at $\vec r(t)$ and $\vec r$ is differentiable at $t$, then $g$ is differentiable at $t$ and
        $$ g'(t) = \nabla f(\vec r(t)) \cdot \vec r\,'(t). $$
    \end{mythm}
    \begin{proof}
        Suppose $\vec r'(t)$ exists for $t\in J$ and let $\vec a = \vec r(t)$. Since $a\in S$ and $S$ is open, there is an open $n$-ball $B(\vec a) \subseteq S$. Note $\vec r$ is continuous at $t$, so we can consider $h\ne 0$ but with $\left\|h\right\|$ small enough so that $\vec r(t+h) \in B(\vec a)$. Let $\vec y = \vec r(t+h) - \vec r(t)$, then
        $$ g(t+h) - g(t) = f(\vec r(t+h)) - f(\vec r(t)) = f(\vec a + \vec y) - f(\vec a) = \nabla f(\vec a) \cdot \vec y + o(\|\vec y\|). $$
        Then
        $$ \frac{g(t+h) - g(t)}{h} = \nabla f(\vec a) \cdot \frac{\vec y}{h} + o\left( \frac{\|\vec y\|}{|h|}\right) = \nabla f(\vec a) \cdot \frac{\vec r(t+h) - \vec r(t)}{h} + o\left( \frac{\|\vec r(t+h) - \vec r(t)\|}{|h|}\right). $$
        Since $\vec r$ is differentiable at $t$, we have $\displaystyle \lim_{h\to 0} \frac{\vec r(t+h) - \vec r(t)}{h} = \vec r\,'(t)$ and thus $\displaystyle \lim_{h\to 0} o\left( \frac{\|\vec r(t+h) - \vec r(t)\|}{|h|}\right) = 0$. Thus
        $$ g'(t) = \lim_{h\to 0} \frac{g(t+h) - g(t)}{h} = \nabla f(\vec a) \cdot \vec r\,'(t). $$
    \end{proof}
    \begin{myex}
        If $\vec r$ describes a curve in $\R^2$ or $\R^3$ and $f:\R^2 \to \R$, $g=f\circ \vec r$, then
        $$ g'(t) = \nabla f(\vec r(t)) \cdot \vec r\,'(t) = \left\|\vec r\,'(t)\right\|\left(\nabla f(\vec r(t)) \cdot T(t)\right) $$
        where $T$ is the unit tangent vector of the curve. Here the term $\nabla f(\vec r(t)) \cdot T(t)$ is called the \textbf{directional derivative} of $f$ along the curve provided it's well-defined.
    \end{myex}

    \subsection{The Derivatives of a Vector Field}
    \begin{mydef}
        Given a vector field $\vec f:S\to \R^m$ where $S\subseteq \R^n$. Let $\vec a \in \operatorname{int} S$ and $\vec y \in \R^n$.
        \begin{itemize}
            \item The \textbf{derivative of $\vec f$ at $\vec a$ with respect to $\vec y$} is denoted by $\vec f'(\vec a; \vec y)$ and is defined as
            $$ \vec{f}\,'(\vec a; \vec y) := \lim_{h\to 0} \frac{\vec f(\vec a + h \vec y) - \vec f(\vec a)}{h}\in \R^m, $$
            provided the limit exists. Let $f_k$ denote the $k$-th component of $\vec f$, then $\vec f\,'(\vec a; \vec y) $ exists if and only if $f_k'(\vec a; \vec y)$ exists for $k=1, \ldots, m$, in which case $$\vec f\,'(\vec a; \vec y) = (f_1'(\vec a; \vec y), \ldots, f_m'(\vec a; \vec y))=\sum_{k=1}^m f_k'(\vec a; \vec y) \vec e_k.$$        
            \item We say that $\vec f$ is \textbf{differentiable} at $\vec a$ if there exists a linear map $\vec T_{\vec a}:\R^n \to \R^m$ and a function $\vec E:\R^n \times \R^n \to \R^m$ such that
            $$ \vec f(\vec a + \vec v) = \vec f(\vec a) + \vec T_{\vec a}(\vec v) + \|v\| \vec E(\vec a, \vec v) $$
            for any $\vec v$ with $\vec a + \vec v \in B(\vec a;r)\subseteq S$ with $r>0$, where $\vec E(\vec a, \vec v) \to \vec 0$ as $\vec v \to \vec 0$. The linear map $\vec T_{\vec a}$ uniquely exists and is called the \textbf{total derivative} of $\vec f$ at $\vec a$.
        \end{itemize}
    \end{mydef}
    \begin{mythm}
        If $\vec f$ is differentiable at $\vec a$ with total derivative $\vec T_{\vec a}$, then $\vec f'(\vec a; \vec y)$ exists for any $\vec y\in \R^n$ and $\vec T_{\vec a}(\vec y) = \vec f'(\vec a; \vec y)$. Furthermore, if we write $\vec f = (f_1, \ldots, f_m)$ and $\vec y = (y_1, \ldots, y_n)$, then
        $$ \vec T_{\vec a}(\vec y) = \vec f'(\vec a; \vec y) = \left( \nabla f_1(\vec a) \cdot \vec y, \ldots, \nabla f_m(\vec a) \cdot \vec y \right) = \sum_{k=1}^n \left( \nabla f_k (\vec a) \cdot \vec y \right) \vec e_k. $$
    \end{mythm}
    \begin{myrmk}
        Note we can also write $\vec T_{\vec a}(\vec y)$ as a matrix-vector product as follows:
        $$ \vec T_{\vec a}(\vec y) = \begin{pmatrix}
            \nabla f_1(\vec a) \\ \vdots \\ \nabla f_m(\vec a)
        \end{pmatrix} \vec y = \begin{pmatrix}
            \dfrac{\partial f_1}{\partial x_1}(\vec a) & \dfrac{\partial f_1}{\partial x_2}(\vec a) & \cdots & \dfrac{\partial f_1}{\partial x_n}(\vec a) \\[8pt]
            \dfrac{\partial f_2}{\partial x_1}(\vec a) & \dfrac{\partial f_2}{\partial x_2}(\vec a) & \cdots & \dfrac{\partial f_2}{\partial x_n}(\vec a) \\[8pt]
                \vdots & \vdots & \ddots & \vdots \\[8pt]
            \dfrac{\partial f_m}{\partial x_1}(\vec a) & \dfrac{\partial f_m}{\partial x_2}(\vec a) & \cdots & \dfrac{\partial f_m}{\partial x_n}(\vec a)
        \end{pmatrix} \vec y. $$
        This $m\times n$ matrix is called the \textbf{Jacobian matrix} of $\vec f$ at $\vec a$ and is denoted by $\mathbf{D}_{\vec f(\vec a)}$.
        The $kj$-th entry is actually $$\left[ \mathbf{D}_{\vec f(\vec a)} \right]_{kj} = \dfrac{\partial f_k}{\partial x_j}(\vec a) = D_j f_k(\vec a).$$

        We often simply denote the total derivative $\vec T_{\vec a}$ by $\vec f\,'(\vec a)$, which is a linear map from $\R^n$ to $\R^m$. Also we know $\vec f\,'(\vec a; \vec y) = \mathbf D_{\vec f(\vec a)} \vec y$ for any $\vec y\in \R^n$. Thus $\mathbf D_{\vec f(\vec a)}$ is the matrix representation of the linear map $\vec f\,'(\vec a)$ with respect to the standard bases of $\R^n$ and $\R^m$. We also have
        $$ \vec f(\vec a + \vec v) = \vec f(\vec a) + \mathbf D_{\vec f(\vec a)} \vec v + \left\|\vec v\right\| \vec E(\vec a, \vec v),$$
        where $\vec E(\vec a, \vec v) \to \vec 0$ as $\vec v \to \vec 0$.
    \end{myrmk}
    \begin{mythm}
        If a vector field $\vec f$ is differentiable at $\vec a$, then it is continuous at $\vec a$.
    \end{mythm}
    \begin{mylma}
        Suppose $\vec g$ is differentiable at $\vec a$. Then $\| \vec g(\vec a) (\vec v) \| \le M(\vec a) \| \vec v\|$ where $M(\vec a)=\displaystyle \sum_{k=1}^m \left\|\nabla g_k(\vec a)\right\|$
    \end{mylma}
    \begin{proof}
        $$ \| \vec g(\vec a) (\vec v) \| = \left\| \sum_{k=1}^m \left( \nabla g_k(\vec a) \cdot \vec v\right) \vec e_k\right\| \le \sum_{k=1}^m \left| \nabla g_k(\vec a) \cdot \vec v\right| \le \sum_{k=1}^m \|\nabla g_k(\vec a)\| \|\vec v\|. $$
    \end{proof}
    \begin{mythm}
        Suppose $\vec f$ and $\vec g$ are vector fields such that the composition $\vec h = \vec f \circ \vec g$ is defined in an opne $n$-ball $B(\vec a)$. If $\vec g$ is differentiable at $\vec a$ with total derivative $\vec g\,'(\vec a)$ and $\vec f$ is differentiable at $\vec b = \vec g(\vec a)$ with total derivative $\vec f\,'(\vec b)$, then $h = \vec f \circ \vec g$ is differentiable at $\vec a$ and
        $$\vec h\,'(\vec a) =  (\vec f \circ \vec g)\,'(\vec a) = \vec f\,'(\vec b) \circ \vec g\,'(\vec a). $$
    \end{mythm}
    \begin{proof}
        We shall consider $\vec h (\vec a + \vec v)$ for $\vec y $ with small $\| \vec y \|$. Note $\vec b = \vec g(\vec a)$, let $\vec v = \vec g(\vec a + \vec y) - \vec g(\vec a) = \vec g(\vec a + \vec y) - \vec b$, then $\vec h (\vec a + \vec y) - \vec h(\vec a) = \vec f(\vec b + \vec v) - \vec f(\vec b)$. 
        Since $f$ is differentiable at $\vec b=\vec g(\vec a)$, we have
        $$ \begin{aligned}
            \vec f(\vec b + \vec v) - \vec f(\vec b) &= \vec f\,'(\vec b)(\vec v) + \|\vec v\| \vec E_f(\vec b, \vec v)\\
            &= \vec f\,'(\vec b)(\vec g\,'(\vec a)(\vec y) + \vec E_g(\vec a, \vec y)) + \|\vec v\| \vec E_f(\vec b, \vec v)\\ 
            &= \vec f\,'(\vec b) \circ \vec g\,'(\vec a)(\vec y) + \|\vec y\| \vec f\,'(\vec b)(\vec E_g(\vec a, \vec y)) + \|\vec v\| \vec E_f(\vec b, \vec v)\\
            &= \vec f\,'(\vec b) \circ \vec g\,'(\vec a)(\vec y) + \|\vec y\| \vec E(\vec a, \vec y),
        \end{aligned}$$
        where $$\vec E(\vec a, \vec y) = \begin{cases}
            \vec f\,'(\vec b)(\vec E_g(\vec a, \vec y)) + \dfrac{\|\vec v\|}{\|\vec y\|} \vec E_f(\vec b, \vec v), & \vec y \ne \vec 0; \\[6pt]
            \vec 0, & \vec y = \vec 0.
        \end{cases}$$
        We still need to show that if $\vec y \to \vec 0$ (with $\vec y \ne \vec 0$), then $\vec E(\vec a, \vec y) \to \vec 0$. Note as $\vec y \to \vec 0$, the first term $\vec f\,'(\vec b)(\vec E_g(\vec a, \vec y)) \to \vec 0$ since $\vec E_g(\vec a, \vec y) \to \vec 0$ due to the differentiability at $\vec a$ and $\vec f\,'(\vec b)$ is a linear map and thus continuous. For the second term $\dfrac{\|\vec v\|}{\|\vec y\|} \vec E_f(\vec b, \vec v)$, we will show that the factor $\dfrac{\|\vec v\|}{\|\vec y\|}$ is bounded as $\vec y \to \vec 0$. Note
        $$ \left\|\vec v\right\| = \left\|\vec g(\vec a + \vec y) - \vec g(\vec a)\right\| = \left\|\vec g\,'(\vec a)(\vec y) + \|\vec y\| \vec E_g(\vec a, \vec y)\right\| \le \left\|\vec g\,'(\vec a)(\vec y)\right\| + \|\vec y\| \left\|\vec E_g(\vec a, \vec y)\right\|. $$
        By the above lemma, we have $\left\|\vec g\,'(\vec a)(\vec y)\right\| \le M(\vec a) \|\vec y\|$ for some constant $M(\vec a)$. Therefore,
        $$ \frac{\|\vec v\|}{\|\vec y\|} \le M(\vec a) + \left\|\vec E_g(\vec a, \vec y)\right\|. $$
        Since $\vec E_g(\vec a, \vec y) \to \vec 0$ as $\vec y \to \vec 0$, we know the factor is bounded as long as $\vec y$ is not so far away from $\vec 0$. Thus $\vec E(\vec a, \vec y) \to \vec 0$ as $\vec y \to \vec 0$ and we have $\vec E(\vec a, \vec y) = o(\|\vec y\|)$. Therefore,
        $$ \vec h(\vec a + \vec y) - \vec h(\vec a) = \vec f\,'(\vec b) \circ \vec g\,'(\vec a)(\vec y) + \| \vec y \| \vec E(\vec a, \vec y) = \vec f\,'(\vec b) \circ \vec g\,'(\vec a)(\vec y) + o(\|\vec y\|), $$
        which means $\vec h$ is differentiable at $\vec a$ with total derivative $\vec h\,'(\vec a) = \vec f\,'(\vec b) \circ \vec g\,'(\vec a)$.
    \end{proof}
    \begin{myrmk}
        Note the corresponding matrix form of a linear map composition $\vec h = \vec f \circ \vec g$ is given by the matrix product of the Jacobian matrices of $\vec f$ and $\vec g$:
        $$ \mathbf D_{\vec h(\vec a)} = \mathbf D_{\vec f(\vec b)} \mathbf D_{\vec g(\vec a)}. $$
        Additionally, suppose $\vec a \in \R^p$, $\vec b \in \R^n$ and $\vec f(\vec b) \in \R^m$, then $\vec h(\vec a) \in \R^m$ and the Jacobian matrix of $\vec h$ at $\vec a$ is given by
        $$ \left[ \mathbf D_{\vec h(\vec a)} \right]_{kj} = \left[ \mathbf D_{\vec f(\vec b)} \mathbf D_{\vec g(\vec a)} \right]_{kj} = \sum_{i=1}^n \left[ \mathbf D_{\vec f(\vec b)} \right]_{ki} \left[ \mathbf D_{\vec g(\vec a)} \right]_{ij} = \sum_{i=1}^n \frac{\partial f_k}{\partial x_i}(\vec b) \frac{\partial g_i}{\partial x_j}(\vec a), $$
        for $k=1, \ldots, m$ and $j=1, \ldots, p$. Equivalently, we can write the chain rule as
        $$ D_j h_k(\vec a) = \sum_{i=1}^n D_i f_k(\vec b) D_j g_i(\vec a). $$
    \end{myrmk}
    \begin{myex}

        For the special case $\vec a\in \R^p$, $\vec b = \vec g(\vec a) \in \R^n$ and $\vec f(\vec b) \in \R$, the chain rule simplifies to
        $$ D_j h(\vec a) = \sum_{i=1}^n D_i f(\vec b) D_j g_i(\vec a),\quad j=1, \ldots, p. $$
        \begin{enumerate}[label=(\arabic*), start=1]
            \item In Section 5.2, we cover the case where $p=1$:
            $$ h'(t) = \sum_{i=1}^n D_i f(\vec b) g_i'(t) = \nabla f(\vec b) \cdot \vec g\,'(t). $$
            \item Consider $p=2$ and $n=2$, $\vec a = (s,t)$, $\vec b = \vec g(\vec a) = (x,y)$ with $x = g_1(s,t)$ and $y = g_2(s,t)$, then
            $$ D_1 h(s,t) = D_1 f(x,y) D_1 g_1(s,t) + D_2 f(x,y) D_1 g_2(s,t), $$
            or equivalently,
            $$ \pdv{h}{s} = \pdv{f}{x} \pdv{x}{s} + \pdv{f}{y} \pdv{y}{s}. $$
        \end{enumerate}
    \end{myex}
    \begin{mythm}
        Suppose $F$ is a scalar field differentiable on an open set $T\subseteq \R^n$ and that the equation $F(x_1, \ldots, x_n) = 0$ defines $x_n$ implicitly as a differentiable function of $x_1, \ldots, x_{n-1}$, say $x_n = f(x_1, \ldots, x_{n-1})$ for all $(x_1, \ldots, x_{n-1})$ in some open set $S\subseteq \R^{n-1}$. Then for each $k=1,2, \ldots, n-1$, the partial derivative $D_k f$ or $\displaystyle \pdv{f}{x_k}$ exists and is given by
        $$ \frac{\partial f}{\partial x_k} = -\frac{D_k F(x_1, \ldots, x_{n-1}, f(x_1, \ldots, x_{n-1}))}{D_n F(x_1, \ldots, x_{n-1}, f(x_1, \ldots, x_{n-1}))} = -\frac{\partial F(x_1, \ldots,x_{n-1}, f(x_1, \ldots, x_{n-1})) / \partial x_k}{\partial F(x_1, \ldots, x_{n-1}, f(x_1, \ldots, x_{n-1})) / \partial x_n}. $$
    \end{mythm}

    \newpage
    \section{Applications of Differential Calculus}
    \subsection{Applications to Geometry}
    \begin{mydef}
        Suppose $f$ is a scalar field on $S\subseteq \R^n$ and $c\in \R$ is a constant. We call the set $L(c) := \{\vec x \in S: f(\vec x) = c\}$ a \textbf{level set} of $f$. In particular, if $n=2$, we call $L(c)$ a \textbf{level curve} of $f$; if $n=3$, we call $L(c)$ a \textbf{level surface} of $f$.
    \end{mydef}
    \begin{myex}
        Suppose $f:\R^2 \to \R$. For any curve $\vec r(t) $ in $\R^2$, if we let $g$ be defined as $g(t) = f[\vec r(t)]$, then $$g'(t) = \nabla f(\vec r(t)) \cdot \vec r\,'(t) = \| \vec r\,'(t)\| \left( \nabla f(\vec r(t)) \cdot T(t) \right).$$
        Particularly, if this curve happens to be the level curve given by $f(x,y)  = c$. Suppose level curve $L(c) := \{ (x,y) : f(x,y) = c\}$ is also described by $\vec r_c(t)$. Then $g(t) = f[\vec r_c(t)] = c$ for any $t$, which means $g'(t) = 0$ for any $t$. Thus
        $$
        \| \vec r_c\,'(t)\| \left( \nabla f(\vec r_c(t)) \cdot T(t) \right) = 0.
        $$
        On this level curve, $\nabla f \cdot T = 0$, and we say: the gradient vector is always normal to the curve $C$; the directional derivative of $f$ along any level curve of $f$ is zero. By Cauchy-Schwarz inequality, we have
        $$ |\nabla f(x,y) \cdot \vec y | \le \|\nabla f(x,y)\| \|\vec y\|, $$
        maximized when $\vec y$ is parallel to $\nabla f(x,y)$, i.e., $\vec y$ is perpendicular to the tangent vector of a level curve.
    \end{myex}
    \begin{myex}
        Suppose $f:\R^3 \to \R$ differentiable. Suppose $L(c)$ is a level surface and $\vec a$ is a point on the level surface. Let $\vec r (t)$ describes a curve that passes through $\vec a$ and lies completely with in $L(c)$. If we consider all curves passing through $\vec a $ and living within $L(c)$, then the tangent vectors of all therse curves at $\vec a$ determin/live in a plane that passes throught $\vec a$ and has a normal vector equal to $\nabla f(\vec a)$ provided that $\nabla f(\vec a) \ne \vec 0$. This plane is can be described by the equation
        $$ \nabla f(\vec a) \cdot (\vec x - \vec a) = 0. $$
    \end{myex}

    \subsection{Applications to Optimization Problems}
    \begin{mydef}
        A scalar field is said to
        \begin{itemize}
            \item have an \textbf{absolute (or global maximum)} at $\vec a \in S \subseteq \R^n$ if $f(\vec a) \ge f(\vec x)$ for any $\vec x \in S$, in which case we also call $f(\vec a)$ the \textbf{absolute maximum value} of $f$ on $S$;
            \item have an \textbf{absolute (or global minimum)} at $\vec a \in S$ if $f(\vec a) \le f(\vec x)$ for any $\vec x \in S$, in which case we also call $f(\vec a)$ the \textbf{absolute minimum value} of $f$ on $S$;
            \item have a \textbf{relative (or local maximum)} at $\vec a \in S$ if there exists $r>0$ such that $f(\vec a) \ge f(\vec x)$ for any $\vec x \in B(\vec a; r)\subseteq S$, in which case we also call $f(\vec a)$ the \textbf{local maximum value} of $f$ at $\vec a$;
            \item have a \textbf{relative (or local minimum)} at $\vec a \in S$ if there exists $r>0$ such that $f(\vec a) \le f(\vec x)$ for any $\vec x \in B(\vec a; r)\subseteq S$, in which case we also call $f(\vec a)$ the \textbf{local minimum value} of $f$ at $\vec a$.
        \end{itemize}
        If $f$ has either a relative maximum or a relative minimum at $\vec a$, then we say $f$ has an \textbf{extremum} at $\vec a$.
    \end{mydef}
    \begin{mythm}
        If $f$ has an extremum at $\vec a\in \int S$ and $f$ is differentiable at $\vec a$, then $\nabla f(\vec a) = \vec 0$.
    \end{mythm}
    \begin{myrmk}
        If $f$ is differentiable at $\vec a$ and $\nabla f(\vec a) = \vec 0$, then $f$ does not necessarily have an extremum at $\vec a$. Assume $z=f(x,y) = x^2 - y^2$, then $\nabla f(0,0) = \vec 0$ but $f$ does not have an extremum at $(0,0)$ since $f$ can be positive or negative in any neighborhood of $(0,0)$.
    \end{myrmk}
    \begin{mydef}
        Suppose scalar field $f$ is differentiable at $\vec a$. If $\nabla f(\vec a) = \vec 0$, then we say $\vec a$ is a \textbf{stationary point} of $f$. A stationary point is called a \textbf{saddle point} if every open $n$-ball $B(\vec a)$ contains both points $\vec x$ with $f(\vec x) > f(\vec a)$ and points $\vec y$ with $f(\vec y) < f(\vec a)$.
    \end{mydef}
    \begin{myrmk}
        Suppose $f$ has a stationary point at $\vec a$, then the nature of this stationary point is determined by the sign of $f(\vec x) - f(\vec a)$ for $\vec x$ close to $\vec a$. Suppose $\vec x = \vec a + \vec y$, then
        $$ f(\vec a + \vec y) - f(\vec a) = \nabla f(\vec a) \cdot \vec y + \| \vec y\| E(\vec a, \vec y),$$
        where $E(\vec a, \vec y) \to 0$ as $\vec y \to \vec 0$.
    \end{myrmk}
    \begin{mythm}
        Suppose $f$ is a scalar field with continuous second-order partial derivatives $D_{ij} f$ in an open $n$-ball $B(\vec a)$. Then for all $\vec y \in \R^n$ such that $(\vec a + \vec y) \in B(\vec a)$, we have
        $$ 
        f(\vec a + \vec y) - f(\vec a ) = \nabla f(\vec a) \cdot \vec y + \frac{1}{2!} \vec y H(\vec a + c \vec y) \vec y^T,
        $$
        where $c\in (0,1)$ and $H(\vec x)$ is called the \textbf{Hessian matrix} whose $ij$-th entry is given by $D_{ij} f(\vec x) = \displaystyle \frac{\partial^2 f}{\partial x_i \partial x_j}(\vec x)$. Note that here $\vec y$ is treated as a row vector. We can also write the second-order Taylor expansion as
        $$
        f(\vec a + \vec y) = f(\vec a) + \nabla f(\vec a) \cdot \vec y + \frac{1}{2!} \vec y H(\vec a) \vec y^T + \| \vec y \|^2 E_2(\vec a, \vec y),
        $$
        where $E_2(\vec a, \vec y) \to \vec 0$ as $\vec y \to \vec 0$.
    \end{mythm}
    \begin{myrmk}
        Since the Hessian matrix is symmetric, we know it has an orthonormal eigenbasis and $\vec y H(\vec a + c \vec y) \vec y^T$ is a quadratic form associated with a unique symmetric matrix. If we know all the eigenvalues of $H(\vec a)$ are positive, then
        $$ f(\vec a + \vec y) - f(\vec a) = \frac{1}{2} \vec y H(\vec a) \vec y^T + \| \vec y\|^2 E_2(\vec a, \vec y) > 0 $$
        for $\vec y$ close enough to $\vec 0$.
    \end{myrmk}
    \begin{proof}
        Fix $\vec y$ and define $g(u) = f(\vec a + u\vec y)$ for $u \in [-1,1]$. Then
        $$ f(\vec a + \vec y ) = g(1) - g(0) = g'(0) + \frac{1}{2} g''(c)$$
        where $c\in (0,1)$ by the second-order Taylor formula for $g$ on $[0,1]$ with the Lagrange form of the remainder.

        Let $\vec r = \vec a + u \vec y$, then $\vec r\,'(u) = \vec y$ and
        $$ g'(u) = \nabla f(\vec r) \cdot \vec r\,'(u) = \sum_{j=1} D_j f(\vec r) y_j. $$
        Thus
        $$ g'(0) = \sum_{j=1}^{n=1} D_j f(\vec a) y_j = \nabla f(\vec a) \cdot \vec y = 0. $$

    \end{proof}
    \begin{mylma}
        Suppose $f$ is a scalar field on $\R^2$ with $D_1 f, D_2 f, D_{12} f, D_{21} f$ exist on an open set $S$. If $(a,b) \in S$ and $D_{12} f$ and $D_{21} f$ are both continuous at $(a,b)$, then $D_{12} f(a,b) = D_{21} f(a,b)$.
    \end{mylma}
    \begin{mylma}
        Suppose $A\in \mathbf M_{n\times n}(\R)$ is a symmetric matrix and $\vec y $ is a row vector in $\R^n$. Let $Q(\vec y) = \vec y A \vec y^T$. Then
        \begin{enumerate}
            \item $Q(\vec y) > 0$ for any $\vec y \ne \vec 0$ if and only if all the eigenvalues of $A$ are positive;
            \item $Q(\vec y) < 0$ for any $\vec y \ne \vec 0$ if and only if all the eigenvalues of $A$ are negative;
            \item $Q(\vec y)$ takes both positive and negative values for $\vec y \ne \vec 0$ if and only if $A$ has both positive and negative eigenvalues.
        \end{enumerate}
    \end{mylma}
    \begin{mythm}
        Suppose $f$ is a scalar field with continuous second-order partial derivatives $D_{ij} f$ in an open $n$-ball $B(\vec a)$ and Let $H(\vec a)$ denote the Hessian matrix of $f$ at $\vec a$, which is a stationary point of $f$. Then
        \begin{enumerate}
            \item If all the eigenvalues of $H(\vec a)$ are positive, i.e., $H(\vec a)$ is positive definite, then $f$ has a local minimum at $\vec a$;
            \item If all the eigenvalues of $H(\vec a)$ are negative, i.e., $H(\vec a)$ is negative definite, then $f$ has a local maximum at $\vec a$;
            \item If $H(\vec a)$ has both positive and negative eigenvalues, then $f$ has a saddle point at $\vec a$.
        \end{enumerate}
    \end{mythm}
    \begin{mythm}
        Given $\vec a = (a_1, \ldots, a_n), \vec b = (b_1, \ldots, b_n) \in \R^n$ with $a_i < b_i$ for $i=1, \ldots, n$. We define $\left[ \vec a, \vec b \right] := [a_1, b_1] \times \cdots \times [a_n, b_n]$. If a scalar field $f$ is continuous at every $\vec x \in [\vec a, \vec b]$, then
        \begin{itemize}
            \item $f$ is bounded on $[\vec a, \vec b]$, i.e., there exists $M>0$ such that $|f(\vec x)| \le M$ for any $\vec x \in [\vec a, \vec b]$;
            \item There exists $\vec c$ and $\vec d$ such that
            $$ f(\vec c) = \sup_{\vec x \in [\vec a, \vec b]} f(\vec x),\quad f(\vec d) = \inf_{\vec x \in [\vec a, \vec b]} f(\vec x). $$
        \end{itemize}
    \end{mythm}
    \begin{mydef}
        Suppose $\left[ \vec a, \vec b \right]$ and $P_k$ is a partition of $[a_k, b_k]$ for $k=1, \ldots, n$. We call $P := P_1 \times \cdots \times P_n$ a \textbf{partition} of $[\vec a, \vec b]$. 
    \end{mydef}
    \begin{mydef}
        Let $f$ be a continuous scalar field on $[\vec a, \vec b] \subseteq \R^n$ and let $M(f)$ and $m(f)$ denote the absolute maximum and minimum values of $f$ on $[\vec a, \vec b]$, respectively. We call $M(f) - m(f)$ the \textbf{span} of $f$ on $[\vec a, \vec b]$.
    \end{mydef}
    \begin{mythm}
        Suppose $f$ is a scalar field continuous on $[\vec a, \vec b]$. Then, for every $\varepsilon > 0$, there is a partition of $[\vec a, \vec b]$ such that the span of $f$ on every subinterval of this partition is less than $\varepsilon$.
    \end{mythm}

    \newpage
    \section{Line Integrals}
    \subsection{Basic Definitions and Properties}
    \begin{mydef}
        A function $\vec r : [a,b] \to \R^n$ that is continuous on $[a,b]$ is also called a \textbf{path} in $\R^n$. If $\vec r \, '$ exists and is continuous on $(a,b)$, then we say $\vec r$ is a \textbf{smooth path}. If there is a partition of $[a,b]$ such that $\vec r$ is smooth on each subinterval of this partition, then we say $\vec r$ is a \textbf{piecewise smooth path}.
    \end{mydef}
    \begin{mydef}
        Let $\vec r:[a,b] \to \R^n$ be a piecewise smooth path and $\vec f$ be a vector field defined and bounded on the image of $\vec r$. The line integral of $\vec f$ along $\vec r$ is defined as
        $$ \int_{\vec r} \vec f \cdot \dd\vec r := \int_a^b \vec f(\vec r(t)) \cdot \vec r\,'(t) \dd t $$
        provided the limit on the right-hand side exists.
    \end{mydef}
    \begin{myrmk}
        Suppose curve $C$ is described by $\vec r(t)$ for $t\in [a,b]$ and $\vec r$ is a piecewise smooth path, then we can also write $\displaystyle \int \vec f \cdot \dd \vec r$ as
        \begin{enumerate}
            \item $\displaystyle \int_C \vec f \cdot \dd \vec r$;
            \item $\displaystyle \int_{\vec a}^{\vec b} \vec f \cdot \dd \vec r$ where $\vec a = \vec r(a)$ and $\vec b = \vec r(b)$;
            \item $\displaystyle \oint \vec f \cdot \dd \vec r$ or $\displaystyle \oint_C \vec f \cdot \dd \vec r$ if $C$ is a closed curve, i.e., $\vec r(a) = \vec r(b)$;
            \item $\displaystyle \int_{C} f_1(x,y) \dd x + f_2(x,y) \dd y$ if $\vec f = (f_1(x,y), f_2(x,y))$ with $x=r_1(t)$ and $y=r_2(t)$.
        \end{enumerate}
    \end{myrmk}
    \begin{myrmk}
        Since we define line integrals in terms of regular integrals of a scalar function on $[a,b]$, line integrals have the following basic properties of regular integral.
    \end{myrmk}
    \begin{mythm}
        Suppose $\vec f$ and $\vec g$ are vector fields defined and bounded on the image of a piecewise smooth path $\vec r$. Then
        \begin{enumerate}[label=(\arabic*)]
            \item Linearity: $$ \int (\alpha \vec f + \beta \vec g) \cdot \dd \vec r = \alpha \int \vec f \cdot \dd \vec r + \beta \int \vec g \cdot \dd \vec r$$ for any scalar $\alpha, \beta$;
            \item Additivity: if curve $C$ can be regarded as two pieces $C_1$ and $C_2$ joined together, then $$ \int_C \vec f \cdot \dd \vec r = \int_{C_1} \vec f \cdot \dd \vec r + \int_{C_2} \vec f \cdot \dd \vec r; $$
        \end{enumerate}
    \end{mythm}
    \begin{mydef}
        Suppose $\vec \beta:[a,b] \to \R^n$ is a continuous path and $u:\R \to \R$ is differentiable on $[c,d]$ with $\left\{ u(s) : s\in [c,d] \right\} = [a,b]$ and $u'(s) \ne 0$ for any $s\in [c,d]$. Then $\vec \gamma: [c,d] \to \R^n$ defined as $\vec \gamma(s) := \vec \beta [u(s)]$ for $s\in [c,d]$ is also a continuous path with the same image as $\vec \beta$. We call $\vec \beta$ and $\vec \gamma$ \textbf{equivalent paths} and $u$ a \textbf{change of parameter function}. Namely, $\vec \gamma$ and $\vec \beta$ describe the same curve but with different parametrizations. If $u'(s) >0$ for $s\in [c,d]$, then we say $\vec \gamma$ and $\vec \beta$ trace out $C$ in the same direction.
    \end{mydef}
\end{document}